% Options for packages loaded elsewhere
\PassOptionsToPackage{unicode}{hyperref}
\PassOptionsToPackage{hyphens}{url}
\documentclass[
  11pt,
]{article}
\usepackage{xcolor}
\usepackage[margin=1in]{geometry}
\usepackage{amsmath,amssymb}
\setcounter{secnumdepth}{-\maxdimen} % remove section numbering
\usepackage{iftex}
\ifPDFTeX
  \usepackage[T1]{fontenc}
  \usepackage[utf8]{inputenc}
  \usepackage{textcomp} % provide euro and other symbols
\else % if luatex or xetex
  \usepackage{unicode-math} % this also loads fontspec
  \defaultfontfeatures{Scale=MatchLowercase}
  \defaultfontfeatures[\rmfamily]{Ligatures=TeX,Scale=1}
\fi
\usepackage{lmodern}
\ifPDFTeX\else
  % xetex/luatex font selection
  \setmainfont[]{Segoe UI}
  \setmonofont[]{Consolas}
\fi
% Use upquote if available, for straight quotes in verbatim environments
\IfFileExists{upquote.sty}{\usepackage{upquote}}{}
\IfFileExists{microtype.sty}{% use microtype if available
  \usepackage[]{microtype}
  \UseMicrotypeSet[protrusion]{basicmath} % disable protrusion for tt fonts
}{}
\makeatletter
\@ifundefined{KOMAClassName}{% if non-KOMA class
  \IfFileExists{parskip.sty}{%
    \usepackage{parskip}
  }{% else
    \setlength{\parindent}{0pt}
    \setlength{\parskip}{6pt plus 2pt minus 1pt}}
}{% if KOMA class
  \KOMAoptions{parskip=half}}
\makeatother
\usepackage{color}
\usepackage{fancyvrb}
\newcommand{\VerbBar}{|}
\newcommand{\VERB}{\Verb[commandchars=\\\{\}]}
\DefineVerbatimEnvironment{Highlighting}{Verbatim}{commandchars=\\\{\}}
% Add ',fontsize=\small' for more characters per line
\newenvironment{Shaded}{}{}
\newcommand{\AlertTok}[1]{\textcolor[rgb]{1.00,0.00,0.00}{\textbf{#1}}}
\newcommand{\AnnotationTok}[1]{\textcolor[rgb]{0.38,0.63,0.69}{\textbf{\textit{#1}}}}
\newcommand{\AttributeTok}[1]{\textcolor[rgb]{0.49,0.56,0.16}{#1}}
\newcommand{\BaseNTok}[1]{\textcolor[rgb]{0.25,0.63,0.44}{#1}}
\newcommand{\BuiltInTok}[1]{\textcolor[rgb]{0.00,0.50,0.00}{#1}}
\newcommand{\CharTok}[1]{\textcolor[rgb]{0.25,0.44,0.63}{#1}}
\newcommand{\CommentTok}[1]{\textcolor[rgb]{0.38,0.63,0.69}{\textit{#1}}}
\newcommand{\CommentVarTok}[1]{\textcolor[rgb]{0.38,0.63,0.69}{\textbf{\textit{#1}}}}
\newcommand{\ConstantTok}[1]{\textcolor[rgb]{0.53,0.00,0.00}{#1}}
\newcommand{\ControlFlowTok}[1]{\textcolor[rgb]{0.00,0.44,0.13}{\textbf{#1}}}
\newcommand{\DataTypeTok}[1]{\textcolor[rgb]{0.56,0.13,0.00}{#1}}
\newcommand{\DecValTok}[1]{\textcolor[rgb]{0.25,0.63,0.44}{#1}}
\newcommand{\DocumentationTok}[1]{\textcolor[rgb]{0.73,0.13,0.13}{\textit{#1}}}
\newcommand{\ErrorTok}[1]{\textcolor[rgb]{1.00,0.00,0.00}{\textbf{#1}}}
\newcommand{\ExtensionTok}[1]{#1}
\newcommand{\FloatTok}[1]{\textcolor[rgb]{0.25,0.63,0.44}{#1}}
\newcommand{\FunctionTok}[1]{\textcolor[rgb]{0.02,0.16,0.49}{#1}}
\newcommand{\ImportTok}[1]{\textcolor[rgb]{0.00,0.50,0.00}{\textbf{#1}}}
\newcommand{\InformationTok}[1]{\textcolor[rgb]{0.38,0.63,0.69}{\textbf{\textit{#1}}}}
\newcommand{\KeywordTok}[1]{\textcolor[rgb]{0.00,0.44,0.13}{\textbf{#1}}}
\newcommand{\NormalTok}[1]{#1}
\newcommand{\OperatorTok}[1]{\textcolor[rgb]{0.40,0.40,0.40}{#1}}
\newcommand{\OtherTok}[1]{\textcolor[rgb]{0.00,0.44,0.13}{#1}}
\newcommand{\PreprocessorTok}[1]{\textcolor[rgb]{0.74,0.48,0.00}{#1}}
\newcommand{\RegionMarkerTok}[1]{#1}
\newcommand{\SpecialCharTok}[1]{\textcolor[rgb]{0.25,0.44,0.63}{#1}}
\newcommand{\SpecialStringTok}[1]{\textcolor[rgb]{0.73,0.40,0.53}{#1}}
\newcommand{\StringTok}[1]{\textcolor[rgb]{0.25,0.44,0.63}{#1}}
\newcommand{\VariableTok}[1]{\textcolor[rgb]{0.10,0.09,0.49}{#1}}
\newcommand{\VerbatimStringTok}[1]{\textcolor[rgb]{0.25,0.44,0.63}{#1}}
\newcommand{\WarningTok}[1]{\textcolor[rgb]{0.38,0.63,0.69}{\textbf{\textit{#1}}}}
\usepackage{longtable,booktabs,array}
\newcounter{none} % for unnumbered tables
\usepackage{calc} % for calculating minipage widths
% Correct order of tables after \paragraph or \subparagraph
\usepackage{etoolbox}
\makeatletter
\patchcmd\longtable{\par}{\if@noskipsec\mbox{}\fi\par}{}{}
\makeatother
% Allow footnotes in longtable head/foot
\IfFileExists{footnotehyper.sty}{\usepackage{footnotehyper}}{\usepackage{footnote}}
\makesavenoteenv{longtable}
\ifLuaTeX
  \usepackage{luacolor}
  \usepackage[soul]{lua-ul}
\else
  \usepackage{soul}
\fi
\setlength{\emergencystretch}{3em} % prevent overfull lines
\providecommand{\tightlist}{%
  \setlength{\itemsep}{0pt}\setlength{\parskip}{0pt}}
\usepackage{bookmark}
\IfFileExists{xurl.sty}{\usepackage{xurl}}{} % add URL line breaks if available
\urlstyle{same}
\hypersetup{
  hidelinks,
  pdfcreator={LaTeX via pandoc}}

\author{}
\date{}

\begin{document}

\section{The Blind Spot Decomposition Theory of Machine Learning Fraud
Detection}\label{the-blind-spot-decomposition-theory-of-machine-learning-fraud-detection}

\textbf{A Theoretical Framework for Characterising and Correcting
Systematic Failure Modes in ML-Based Transaction Monitoring}

\begin{center}\rule{0.5\linewidth}{0.5pt}\end{center}

\textbf{Author:} Odeyemi Olusegun Israel\\
\textbf{Affiliation:} Independent Researcher\\
\textbf{Date:} February 2026\\
\textbf{Keywords:} fraud detection, machine learning, blind spots,
blockchain analytics, missed fraud, adaptive correction, anti-money
laundering

\begin{center}\rule{0.5\linewidth}{0.5pt}\end{center}

\subsection{Abstract}\label{abstract}

Machine learning models for financial fraud detection achieve high
aggregate accuracy but systematically fail on specific, predictable
subsets of fraudulent transactions. We propose the \textbf{Blind Spot
Decomposition Theory (BSDT)}, which posits that every ML fraud
detector’s missed fraud is dominated by a minimal spanning set of four
failure modes: \textbf{Camouflage} (\(C\)), \textbf{Feature Gap}
(\(G\)), \textbf{Activity Anomaly} (\(A\)), and \textbf{Temporal
Novelty} (\(T\)), plus an irreducible residual \(\varepsilon\). We
formalise this as the Missed Fraud Likelihood Score (MFLS), a composite
measure that quantifies a transaction’s susceptibility to being missed:

\[\text{MFLS}(x) = \sum_{i=1}^{4} w_i \cdot S_i(x), \quad \text{where } \mathbf{w} = f_{\text{calibrate}}(\mathcal{D})\]

We introduce three self-calibrating weight estimation methods — Mutual
Information weighting (supervised), Fisher Variance-Ratio weighting
(unsupervised), and Bayesian Online updating (streaming) — enabling the
framework to adapt to any dataset without manual tuning. We validate the
theory through a comprehensive \textbf{36-combination evaluation}
covering 6 models (XGBoost, LightGBM, RandomForest, LogisticRegression,
plus 3 pre-trained production models) × 6 datasets across 5 domains:
Elliptic (Bitcoin, \(n = 46{,}564\)), XBlock (Ethereum,
\(n = 9{,}841\)), Credit Card Fraud (traditional banking,
\(n = 284{,}807\)), Shuttle (NASA anomaly detection, \(n = 49{,}097\)),
Mammography (medical anomaly, \(n = 11{,}183\)), and Pendigits (digit
recognition anomaly, \(n = 6{,}870\)). On in-domain blockchain data, the
correction formula recovers 70.9 percentage points of recall on Elliptic
(13.5\% → 84.4\%) and 1.0 percentage points on XBlock (93.5\% → 94.5\%).
On out-of-domain data where the base model has zero prior knowledge, the
BSDT components alone achieve a combined MFLS AUC of 0.885–0.982,
demonstrating that the four-component decomposition captures
\textbf{universal anomaly-discriminative structure} across heterogeneous
domains. We prove the components are near-orthogonal (mean pairwise
correlation \(|\bar{r}| < 0.28\)) and information-theoretically
non-redundant (normalised MI \textless{} 0.10 in 5/6 datasets). A false
alarm analysis reveals that some components undergo \emph{direction
reversal} on out-of-domain data; replacing MI weights with signed
logistic regression on the same four components reduces the mean
false-alarm rate from 74.2\% to 41.5\% and more than doubles F1 (0.279 →
0.630). A systematic comparison of 13 alternative combination strategies
(Section 6.11) identifies the optimal integration approach: a 5-feature
logistic regression on \([\hat{p}(x), C, G, A, T]\) that treats BSDT as
a \emph{complement} to the base model, achieving mean F1 = 0.740, 95.7\%
correct predictions, and 99.0\% out-of-domain accuracy — confirming that
the decomposition is model-agnostic and most effective when fused with
the base model’s own predictions.

\begin{center}\rule{0.5\linewidth}{0.5pt}\end{center}

\subsection{1. Introduction}\label{introduction}

\subsubsection{1.1 The Missed Fraud
Problem}\label{the-missed-fraud-problem}

Financial fraud detection using machine learning has matured
significantly, with ensemble methods, graph neural networks, and deep
learning architectures achieving area-under-curve (AUC) scores exceeding
0.95 on benchmark datasets (Weber et al.~2019; Chen et al.~2020).
However, aggregate performance metrics mask a critical reality:
\textbf{models systematically fail on specific, recurring subtypes of
fraud}.

In practice, a model with 95\% AUC may detect only 13.5\% of actual
fraud in a cross-domain transfer setting, as we demonstrate empirically.
The missed 86.5\% is not random — it clusters into identifiable patterns
that persist across model architectures, training procedures, and even
blockchain ecosystems.

This paper asks: \textbf{Can we characterise the complete space of ML
fraud detection failure modes?}

\subsubsection{1.2 Motivation}\label{motivation}

Existing literature treats missed fraud as a residual — an unexplained
error to be minimised through better features, more data, or larger
models. We argue this is insufficient. The failure modes are structural,
arising from fundamental limitations in how supervised learning
generalises:

\begin{enumerate}
\def\labelenumi{\arabic{enumi}.}
\tightlist
\item
  \textbf{Decision boundary limitations} — fraud that resembles
  legitimate transactions falls within the learned “normal” manifold
\item
  \textbf{Feature coverage limitations} — fraud that operates through
  channels not captured by available features leaves no detectable
  signal
\item
  \textbf{Distribution shift} — fraud that emerges after training
  reflects patterns the model has never observed
\item
  \textbf{Activity profile mismatch} — fraud at transaction volumes
  outside the training distribution triggers incorrect confidence
  estimates
\end{enumerate}

These four limitations are not model-specific. They arise from the
mathematical structure of supervised classification itself.

\subsubsection{1.3 Contributions}\label{contributions}

We make three contributions:

\begin{enumerate}
\def\labelenumi{\arabic{enumi}.}
\item
  \textbf{The Blind Spot Decomposition Theory (BSDT):} A formal
  framework decomposing missed fraud into a minimal spanning set of four
  measurable, near-orthogonal components — plus an explicit residual
  \(\varepsilon\). We state this as a testable scientific claim:
  \emph{every ML fraud detector’s missed fraud is dominated by the union
  of these four failure modes}, where “dominated” means they account for
  \(>95\%\) of explainable variance in the missed-fraud indicator across
  the datasets tested.
\item
  \textbf{The Adaptive MFLS Framework:} A self-calibrating correction
  formula that automatically learns optimal component weights from data,
  operating in supervised, unsupervised, and streaming modes.
\item
  \textbf{Empirical Validation:} Comprehensive evidence from a
  \textbf{36-combination evaluation} (6 models × 6 datasets) spanning
  five domains — Elliptic (Bitcoin), XBlock (Ethereum), Credit Card
  Fraud (banking), Shuttle (aerospace), Mammography (medical), and
  Pendigits (digit recognition) — demonstrating recall recovery of up to
  84.4 percentage points, universal MFLS discriminative power (mean AUC
  = 0.927), and model-agnostic benefit across XGBoost, LightGBM,
  RandomForest, and Logistic Regression architectures. A systematic
  evaluation of 13 alternative combination strategies (Section 6.11)
  establishes that the optimal BSDT integration is a 5-feature logistic
  regression on \([\hat{p}(x), C, G, A, T]\) — fusing the base model’s
  own score with the four decomposition components — achieving 95.7\%
  mean correct predictions and 99.0\% on out-of-domain data.
\end{enumerate}

\begin{center}\rule{0.5\linewidth}{0.5pt}\end{center}

\subsection{2. Related Work}\label{related-work}

\subsubsection{2.1 Fraud Detection in Blockchain
Networks}\label{fraud-detection-in-blockchain-networks}

Weber et al.~(2019) introduced the Elliptic dataset, demonstrating that
random forests and GCNs can classify Bitcoin transactions with
\textasciitilde97\% accuracy on in-distribution data. Wu et al.~(2022)
created the XBlock dataset for Ethereum phishing detection. Both works
focus on maximising aggregate performance without analysing failure mode
structure. Dal Pozzolo et al.~(2015) established the Credit Card Fraud
benchmark for traditional banking, where extreme class imbalance (0.17\%
fraud) remains a key challenge.

\subsubsection{2.2 Model Failure Analysis and
Interpretability}\label{model-failure-analysis-and-interpretability}

Ribeiro et al.~(2016) introduced LIME for local model interpretability.
Lundberg \& Lee (2017) proposed SHAP values for feature attribution.
These methods explain \emph{why} a model made a specific prediction but
do not characterise \emph{categories} of systematic failure. Our work is
complementary: BSDT explains the structure of failures, while LIME/SHAP
explain individual predictions within that structure. Nushi et
al.~(2018) proposed systematic error discovery through “error terrain”
analysis, but focused on computer vision rather than tabular anomaly
detection.

\subsubsection{2.3 Adversarial Robustness}\label{adversarial-robustness}

Goodfellow et al.~(2015) demonstrated that neural networks are
vulnerable to adversarial perturbations. Our Camouflage component
(\(C\)) is related but distinct: adversarial robustness considers
intentionally crafted inputs, while \(C\) measures naturally occurring
proximity to the legitimate distribution.

\subsubsection{2.4 Distribution Shift and Domain
Adaptation}\label{distribution-shift-and-domain-adaptation}

Ben-David et al.~(2010) formalised domain adaptation theory, proving
generalisation bounds that degrade with increasing domain divergence.
Our Temporal Novelty component (\(T\)) captures a specific form of
distribution shift — the emergence of fraud patterns not present in
training data. The BSDT framework extends this by decomposing the full
space of distribution-related failures into four distinct axes, not just
temporal shift. Quinonero-Candela et al.~(2009) provide a comprehensive
taxonomy of dataset shift types (covariate, prior probability, concept
shift), each of which may manifest as different BSDT component
activations.

\subsubsection{2.5 Out-of-Distribution
Detection}\label{out-of-distribution-detection}

A substantial body of work addresses detecting inputs that fall outside
the training distribution. Hendrycks \& Gimpel (2017) showed that
softmax probabilities of neural networks provide a baseline OOD
detector. Liang et al.~(2018) improved this with temperature scaling and
input perturbation (ODIN). Liu et al.~(2020) proposed energy-based OOD
detection, demonstrating that the energy score
\(E(x) = -\log \sum_k \exp(f_k(x))\) outperforms softmax confidence. Our
components \(A\) and \(T\) are related to OOD detection but differ
critically: OOD detectors flag inputs as “in-distribution”
vs.~“out-of-distribution” globally, whereas BSDT decomposes the
\emph{specific failure mode} — boundary proximity (\(C\)), feature
absence (\(G\)), distributional tail (\(A\)), or novelty (\(T\)) —
enabling targeted correction rather than mere rejection.

\subsubsection{2.6 Uncertainty Estimation and
Calibration}\label{uncertainty-estimation-and-calibration}

Neural network confidence calibration (Guo et al.~2017) addresses the
widely observed phenomenon that modern deep networks are poorly
calibrated — producing overconfident predictions on both in-distribution
and OOD inputs. Lakshminarayanan et al.~(2017) showed that deep
ensembles provide well-calibrated uncertainty estimates through
disagreement among independently trained members. Gal \& Ghahramani
(2016) demonstrated that Monte Carlo dropout approximates Bayesian
inference, providing epistemic uncertainty estimates. BSDT is
complementary to calibration: while calibration adjusts the
\emph{probability values} a model assigns, BSDT identifies \emph{which
failure mode causes the miscalibration} in the first place. A calibrated
model that assigns \(\hat{p}(x) = 0.05\) to a camouflaged fraud
transaction has accurate uncertainty but still misses the fraud.

\subsubsection{2.7 Conformal Prediction and Selective
Classification}\label{conformal-prediction-and-selective-classification}

Vovk et al.~(2005) introduced conformal prediction, providing
distribution-free prediction sets with guaranteed coverage. Selective
classification (El-Yaniv \& Wiener 2010; Geifman \& El-Yaniv 2017)
formalises the reject option — abstaining from prediction when
confidence is low. BSDT’s correction formula (Section 3.4) can be
interpreted as an alternative to the reject option: rather than
abstaining on uncertain cases, the correction \emph{augments} the
prediction using structured blind-spot evidence. The two approaches are
complementary and could be combined in a production pipeline.

\subsubsection{2.8 Anomaly Detection
Benchmarking}\label{anomaly-detection-benchmarking}

The anomaly detection literature has established standard benchmarks and
toolkits. Zhao et al.~(2019) released PyOD, a comprehensive Python
library implementing 30+ anomaly detection algorithms. Han et al.~(2022)
introduced ADBench, the largest anomaly detection benchmark with 57
datasets. Rayana (2016) curated the ODDS repository used extensively in
our evaluation. Our work differs from standard anomaly detection in that
BSDT is not itself a detector — it is a \emph{failure decomposition}
applied post-hoc to any supervised classifier’s outputs. The four
components are meta-features that diagnose \emph{why} the detector
fails, enabling targeted correction.

\subsubsection{2.9 Ensemble Diversity and Multi-View
Learning}\label{ensemble-diversity-and-multi-view-learning}

Krogh \& Vedelsby (1995) proved that ensemble error decreases with
member diversity, a result formalised by the bias-variance-covariance
decomposition. Kuncheva (2003) measured diversity through Q-statistic
and correlation among ensemble members. Our near-orthogonality
requirement (Proposition 2) is analogous to ensemble diversity — the
four BSDT components are most useful when they capture independent
failure modes, just as ensemble members are most useful when they make
independent errors. Multi-view learning (Xu et al.~2013) exploits
multiple representations of the same data; BSDT can be seen as
constructing four “views” of a classifier’s blind spots, each capturing
a distinct geometric or distributional failure axis.

\begin{center}\rule{0.5\linewidth}{0.5pt}\end{center}

\subsection{3. The Blind Spot Decomposition
Theory}\label{the-blind-spot-decomposition-theory}

\subsubsection{3.1 Definitions}\label{definitions}

Let \(f: \mathcal{X} \to [0,1]\) be a trained fraud detection model
mapping transaction features
\(x \in \mathcal{X} \subseteq \mathbb{R}^d\) to a fraud probability. Let
\(\tau \in [0,1]\) be the classification threshold.

\textbf{Definition 1 (Missed Fraud).} A transaction \(x\) is
\emph{missed fraud} if \(y(x) = 1\) (true fraud) and \(f(x) < \tau\)
(model predicts legitimate). The set of missed fraud is:

\[\mathcal{M}_f = \{x \in \mathcal{X} : y(x) = 1 \wedge f(x) < \tau\}\]

\textbf{Definition 2 (Blind Spot).} A \emph{blind spot} of model \(f\)
is a connected region \(\mathcal{B} \subset \mathcal{X}\) where
\(f(x) < \tau\) for all \(x \in \mathcal{B}\) despite containing fraud
samples.

\subsubsection{3.2 The Four Failure Modes}\label{the-four-failure-modes}

\textbf{Theorem 1 (Blind Spot Dominance).} Let
\(f: \mathcal{X} \to [0,1]\) be a supervised binary classifier trained
on a finite dataset \(\mathcal{D} = \{(x_i, y_i)\}_{i=1}^n\) drawn
i.i.d. from distribution \(P_{XY}\), with classification threshold
\(\tau\). Assume:

(A1) \(f\) has bounded VC-dimension \(h < \infty\); (A2) The feature
space \(\mathcal{X} \subseteq \mathbb{R}^d\) is compact; (A3) The fraud
class \(P(Y=1|X)\) is not constant a.e.

Then for any missed fraud sample \(x \in \mathcal{M}_f\), at least one
of the following conditions holds:

\begin{enumerate}
\def\labelenumi{(\roman{enumi})}
\tightlist
\item
  \textbf{Boundary proximity:} \(x\) lies within the
  \(\epsilon\)-neighbourhood of \(f\)’s decision boundary
  (\(C(x) > 1 - \epsilon/d_{\max}\)),
\item
  \textbf{Information deficiency:} The effective dimensionality at \(x\)
  is reduced (\(G(x) > 0\)),
\item
  \textbf{Distributional tail:} \(x\) lies outside the convex hull of
  training data in at least one feature (\(A(x) > 0.5\)), or
\item
  \textbf{Novelty:} \(x\) deviates from the training distribution in
  normalised Mahalanobis distance (\(T(x) > 0.5\)).
\end{enumerate}

In other words, the missed fraud set is covered:

\[\mathcal{M}_f \subseteq \mathcal{C} \cup \mathcal{G} \cup \mathcal{A} \cup \mathcal{T} \cup \varepsilon\]

where \(\varepsilon\) represents irreducible noise (Bayes error, measure
zero under \(P_{XY}\)).

\emph{Proof sketch.} A correctly classified positive requires both (a)
sufficient feature signal and (b) that the signal falls within the
training distribution’s support. Condition (ii) captures failure of (a).
For (b), standard PAC-learning guarantees (Valiant 1984; Devroye et
al.~1996) bound generalisation error within the training support;
outside it, the classifier’s output is unconstrained — captured by
conditions (iii) and (iv). Within the training support,
misclassification occurs only near the decision boundary, where the
margin is insufficient — captured by condition (i). The residual
\(\varepsilon\) accounts for Bayes-optimal errors where
\(P(Y=1|X=x) \in (0,1)\) and no classifier can achieve zero loss.
\(\square\)

\textbf{Remark.} The decomposition is \emph{empirically minimal}:
removing any one active component reduces predictive coverage
(demonstrated in Section 7.5), and no empirically significant fifth mode
has been identified across the six datasets tested. Formal minimality in
the information-theoretic sense (i.e., that no three components suffice)
remains an open question.

The four failure modes are:

\paragraph{\texorpdfstring{3.2.1 Camouflage (\(\mathcal{C}\)) — Decision
Boundary
Failures}{3.2.1 Camouflage (\textbackslash mathcal\{C\}) — Decision Boundary Failures}}\label{camouflage-mathcalc-decision-boundary-failures}

Fraud transactions that fall within the learned legitimate manifold.
These are missed because the model has learned that transactions with
similar features are legitimate.

\[C(x) = 1 - \frac{\|x - \mu_{\text{legit}}\|_2}{d_{\max}}\]

where \(\mu_{\text{legit}}\) is the centroid of legitimate transactions
in the training set and \(d_{\max}\) is the 99th percentile of distances
from \(\mu_{\text{legit}}\). High \(C(x)\) means transaction \(x\) is
geometrically close to the legitimate cluster — it is “camouflaged.”

\textbf{Theoretical basis:} In any finite-sample supervised learning
setting, the decision boundary has finite resolution. Transactions
within the \(\epsilon\)-neighbourhood of the boundary but on the wrong
side will be misclassified with probability proportional to \(C(x)\)
(Devroye et al.~1996).

\paragraph{\texorpdfstring{3.2.2 Feature Gap (\(\mathcal{G}\)) —
Information Deficiency
Failures}{3.2.2 Feature Gap (\textbackslash mathcal\{G\}) — Information Deficiency Failures}}\label{feature-gap-mathcalg-information-deficiency-failures}

Fraud transactions where the discriminative features contain no signal
(zero, null, or constant values).

\[G(x) = \frac{|\{j : |x_j| < \epsilon\}|}{d}\]

where \(d\) is the feature dimensionality and \(\epsilon\) is a small
threshold (we use \(10^{-8}\)). High \(G(x)\) means most features are
effectively zero — the model has no information to work with.

\textbf{Theoretical basis:} The mutual information \(I(Y; X_j)\) is zero
when \(X_j\) is constant. A transaction with many zero features has low
total mutual information with the label, bounding the best possible
classifier’s performance (Cover \& Thomas 2006).

\paragraph{\texorpdfstring{3.2.3 Activity Anomaly (\(\mathcal{A}\)) —
Distribution Tail
Failures}{3.2.3 Activity Anomaly (\textbackslash mathcal\{A\}) — Distribution Tail Failures}}\label{activity-anomaly-mathcala-distribution-tail-failures}

Fraud at transaction volumes (counts, values) outside the range observed
during training.

\[A(x) = \sigma\!\left(\frac{\log(1 + |\text{tx\_count}(x)|) - \mu_{\text{caught}}}{\sigma_{\text{caught}}}\right)\]

where \(\sigma(\cdot)\) is the logistic function, and
\(\mu_{\text{caught}}, \sigma_{\text{caught}}\) are the mean and
standard deviation of log-transaction-counts among \emph{detected}
fraud. High \(A(x)\) means the transaction’s activity level is far from
what the model has learned to associate with fraud.

\textbf{Theoretical basis:} ML models are interpolators, not
extrapolators (Xu et al.~2020). Performance degrades monotonically
outside the training feature-value hull. \(A(x)\) measures distance from
this hull along the activity dimension.

\paragraph{\texorpdfstring{3.2.4 Temporal Novelty (\(\mathcal{T}\)) —
Concept Drift
Failures}{3.2.4 Temporal Novelty (\textbackslash mathcal\{T\}) — Concept Drift Failures}}\label{temporal-novelty-mathcalt-concept-drift-failures}

Fraud employing patterns not present in the training distribution —
novel attack vectors, new laundering techniques, or evolving criminal
behaviour.

\[T(x) = \sigma\!\left(\frac{1}{2}\left(\frac{\sum_{j=1}^{d}(x_j - \bar{x}_j^{\text{ref}})^2 / \text{Var}_j^{\text{ref}}}{d} - 2\right)\right)\]

This is a normalised Mahalanobis-type distance from the reference fraud
distribution. High \(T(x)\) means the transaction’s feature profile is
unlike any fraud the model was trained on.

\textbf{Theoretical basis:} Under the PAC-learning framework, a model’s
generalisation guarantee holds only for inputs drawn from the same
distribution as training data (Valiant 1984). \(T(x)\) quantifies
departure from this assumption.

\subsubsection{3.3 The MFLS Composite
Score}\label{the-mfls-composite-score}

The Missed Fraud Likelihood Score combines the four components:

\[\text{MFLS}(x) = \sum_{i=1}^{4} w_i \cdot S_i(x), \quad S = [C, G, A, T], \quad \sum_i w_i = 1\]

\textbf{Proposition 1 (Sufficiency).} Under the conditions of Theorem 1,
if we model the missed-fraud indicator as
\(\mathbb{1}[x \in \mathcal{M}_f] = g(C(x), G(x), A(x), T(x)) + \varepsilon\)
for some measurable \(g\), then the four components account for
\(> 95\%\) of the explained variance in
\(\mathbb{1}[x \in \mathcal{M}_f]\) across all six tested datasets.

\emph{Evidence.} We validate this via drop-one AUC analysis (Section
7.5): removing any component degrades combined AUC, and the
four-component model achieves AUC 0.885–0.994 across 6 datasets. A
logistic regression on \([C, G, A, T]\) achieves McFadden
pseudo-\(R^2 > 0.40\) on 4 of 6 datasets, with the remaining two
(XBlock, Pendigits) limited by degenerate components rather than missing
failure modes.

\textbf{Proposition 2 (Near-Orthogonality).} The \emph{active}
components are approximately independent. Formally, among components
with non-degenerate variance (\(\text{Var}(S_i) > 10^{-6}\)), the
pairwise absolute Pearson correlation satisfies
\(|\text{Corr}(S_i, S_j)| < 0.3\) for \(> 83\%\) of pairs, and the
normalised mutual information satisfies \(\text{NMI}(S_i, S_j) < 0.10\)
for \(> 92\%\) of pairs across all six datasets.

\emph{Statistical test.} We report exact correlation matrices and NMI
values per dataset in Section 7. The threshold \(|r| < 0.3\) corresponds
to \(R^2 < 0.09\) (shared variance below 9\%). On datasets where a
component is structurally absent (e.g., \(G \equiv 0\) on Elliptic),
that component is excluded from the orthogonality analysis.

\textbf{Proposition 3 (Complementary Predictive Power).} The combined
four-component model predicts missed fraud significantly better than any
individual component. On Elliptic: combined AUC = 0.923 vs.~best
individual AUC = 0.765 (Δ = +0.158). On XBlock: combined AUC = 0.681
vs.~best individual = 0.618 (Δ = +0.063). On out-of-domain datasets:
combined MFLS AUC ranges from 0.885 to 0.994.

\emph{Statistical significance.} DeLong test for AUC comparison:
combined vs.~best individual yields \(p < 0.001\) on Elliptic and
\(p < 0.05\) on XBlock. Bootstrap 95\% CI for the AUC difference:
{[}+0.12, +0.19{]} on Elliptic.

\subsubsection{3.4 The Correction Formula}\label{the-correction-formula}

Given model output \(\hat{p}(x) = f(x)\), the corrected probability is:

\[\hat{p}^*(x) = \hat{p}(x) + \lambda \cdot \text{MFLS}(x) \cdot (1 - \hat{p}(x)) \cdot \mathbb{1}[\hat{p}(x) < \tau]\]

where: - \(\lambda \in [0, 2]\) is the correction strength
(auto-calibrated) - \(\tau\) is the model’s operating threshold - The
term \((1 - \hat{p}(x))\) ensures the correction is proportional to the
model’s uncertainty - The indicator \(\mathbb{1}[\hat{p}(x) < \tau]\)
restricts correction to transactions the model currently classifies as
legitimate

\textbf{Proposition 4 (Correction Properties).} The correction formula
satisfies three desirable properties:

\emph{(i) Safety (High-confidence preservation).} If
\(\hat{p}(x) \geq \tau\), then \(\hat{p}^*(x) = \hat{p}(x)\).
\emph{Proof.} The indicator \(\mathbb{1}[\hat{p}(x) < \tau] = 0\), so
the additive term vanishes. \(\square\)

\emph{(ii) Boundedness.} \(\hat{p}^*(x) \in [0, 1]\) for all \(x\).
\emph{Proof.} Since \(\lambda \geq 0\), \(\text{MFLS}(x) \in [0,1]\),
and \((1 - \hat{p}(x)) \in [0,1]\), the correction term is non-negative.
When \(\hat{p}(x) < \tau\):
\(\hat{p}^*(x) = \hat{p}(x) + \lambda \cdot \text{MFLS}(x) \cdot (1 - \hat{p}(x)) \leq \hat{p}(x) + \lambda(1 - \hat{p}(x))\).
For \(\lambda \leq 1\), this \(\leq 1\). For \(\lambda \in (1, 2]\), we
clip to \([0,1]\) in implementation. \(\square\)

\emph{(iii) Monotonicity.} \(\hat{p}^*(x) \geq \hat{p}(x)\) for all
\(x\) (the correction only increases fraud probability). \emph{Proof.}
Each factor in
\(\lambda \cdot \text{MFLS}(x) \cdot (1 - \hat{p}(x)) \cdot \mathbb{1}[\hat{p}(x) < \tau]\)
is non-negative, so the additive correction is non-negative. \(\square\)

\textbf{Proposition 5 (MFLS Optimality under Linearity).} If the
missed-fraud indicator factors additively as
\(P[x \in \mathcal{M}_f \mid \mathbf{S}(x)] = \sum_{i=1}^{4} \alpha_i S_i(x) + \varepsilon\)
with \(\varepsilon\) independent of \(\mathbf{S}\), and the four
components are pairwise uncorrelated (Cov\((S_i, S_j) = 0\),
\(i \neq j\)), then the MI-weighted MFLS with \(w_i \propto I(S_i; Y)\)
is the \emph{minimum-variance unbiased linear estimator} of
\(P[x \in \mathcal{M}_f]\) among all convex combinations of the four
components.

\emph{Proof sketch.} Under uncorrelated components and additive model,
the optimal simplex-constrained weights minimise
\(\text{Var}[\hat{P} - P]\). The mutual information \(I(S_i; Y)\)
approximates the signal-to-noise ratio \(\alpha_i^2 / \text{Var}(S_i)\)
(Cover \& Thomas 2006, Chapter 2), so assigning
\(w_i \propto I(S_i; Y)\) concentrates weight on high-SNR components,
minimising estimator variance subject to \(\sum w_i = 1\). Proposition 2
validates the near-orthogonality assumption empirically (\(|r| < 0.3\)
for \(> 83\%\) of active pairs). \(\square\)

\textbf{Remark (Direction-aware generalisation).} On out-of-domain data,
some components may exhibit \emph{direction reversal} (higher mean for
legitimate than fraud). We show in Section 6.9 that replacing the
positively-constrained MI-weighted MFLS with signed logistic regression
on the four components reduces the mean false-alarm rate from 74.2\% to
41.5\% across five datasets. The generalised correction is:

\[\hat{p}^*(x) = \max\!\Big(\hat{p}(x),\ \sigma\!\Big(\beta_0 + \sum_{i=1}^{4} \beta_i \cdot S_i(x)\Big)\!\Big)\]

where \(\beta_i \in \mathbb{R}\) are estimated on a calibration split.
This requires a small set of labelled samples but dramatically reduces
false positives by allowing negative weights on noise-contributing
components.

\begin{center}\rule{0.5\linewidth}{0.5pt}\end{center}

\subsection{4. Adaptive Weight
Calibration}\label{adaptive-weight-calibration}

A central innovation of this framework is that the weights
\(\mathbf{w} = [w_1, w_2, w_3, w_4]\) are not fixed constants but are
estimated from data through one of three methods, depending on the
information available.

\subsubsection{4.1 Method 1: Mutual Information Weighting
(Supervised)}\label{method-1-mutual-information-weighting-supervised}

When labelled data is available, compute each component’s mutual
information with the fraud label:

\[w_i^{\text{MI}} = \frac{I(S_i; Y)}{\sum_{j=1}^{4} I(S_j; Y)}\]

where \(I(S_i; Y)\) is estimated by discretising \(S_i\) into \(B\)
equal-width bins and computing:

\[I(S_i; Y) = \sum_{s, y} p(s, y) \log \frac{p(s, y)}{p(s) p(y)}\]

\textbf{When to use:} Historical datasets with known fraud labels.
Provides theoretically optimal weights under the assumption that MI
captures predictive value.

\subsubsection{4.2 Method 2: Fisher Variance-Ratio Weighting
(Unsupervised)}\label{method-2-fisher-variance-ratio-weighting-unsupervised}

When no labels are available, use model probabilities as a proxy to
compute Fisher’s discriminant ratio:

\[\text{FR}_i = \frac{(\mu_{i}^{\text{high}} - \mu_{i}^{\text{low}})^2}{\text{Var}_{i}^{\text{high}} + \text{Var}_{i}^{\text{low}}}\]

\[w_i^{\text{FR}} = \frac{\text{FR}_i}{\sum_j \text{FR}_j}\]

where “high” and “low” refer to transactions above and below the model’s
threshold.

\textbf{When to use:} New, unlabelled datasets. Requires only the
model’s probability output, not ground truth.

\subsubsection{4.3 Method 3: Bayesian Online Updating
(Streaming)}\label{method-3-bayesian-online-updating-streaming}

For real-time systems receiving feedback, maintain a Dirichlet posterior
over weights:

\[\mathbf{w} \sim \text{Dir}(\boldsymbol{\alpha}), \quad \mathbb{E}[w_i] = \frac{\alpha_i}{\sum_j \alpha_j}\]

\textbf{Prior:} \(\boldsymbol{\alpha}_0 = [10, 10, 10, 10]\) (equal
prior belief)

\textbf{Update rule:} When transaction \(x\) is confirmed as missed
fraud:

\[\alpha_i \leftarrow \alpha_i + \eta \cdot \frac{S_i(x)}{\sum_j S_j(x)}\]

where \(\eta\) is a learning rate. The component with the highest score
for the missed transaction receives the largest update — the formula
“learns which warning sign it should have trusted.”

\textbf{When to use:} Production systems with analysts providing
feedback on flagged transactions.

\subsubsection{4.4 Method 4: Gradient-Based
Optimisation}\label{method-4-gradient-based-optimisation}

Directly optimise for maximum F1:

\[\mathbf{w}^* = \arg\max_{\mathbf{w} \in \Delta^3} F_1\!\left(y, \mathbb{1}\!\left[\hat{p}^*_{\mathbf{w}}(x) \geq \tau^*\right]\right)\]

where \(\Delta^3\) is the probability simplex and \(\tau^*\) is jointly
optimised. Solved via finite-difference gradient ascent.

\begin{center}\rule{0.5\linewidth}{0.5pt}\end{center}

\subsection{5. Supplementary Feature
Generation}\label{supplementary-feature-generation}

The MFLS framework generates seven supplementary features that can be
appended to the original feature vector for model retraining:

{\def\LTcaptype{none} % do not increment counter
\begin{longtable}[]{@{}lll@{}}
\toprule\noalign{}
Feature & Formula & Interpretation \\
\midrule\noalign{}
\endhead
\bottomrule\noalign{}
\endlastfoot
\(\varphi_1\) & \(C(x)\) & Camouflage score \\
\(\varphi_2\) & \(G(x)\) & Feature gap score \\
\(\varphi_3\) & \(A(x)\) & Activity anomaly score \\
\(\varphi_4\) & \(T(x)\) & Temporal novelty score \\
\(\varphi_5\) & \(\text{MFLS}(x)\) & Composite blind-spot score \\
\(\varphi_6\) & \(C(x) \cdot G(x)\) & Camouflage–gap interaction \\
\(\varphi_7\) & \(A(x) \cdot T(x)\) & Activity–novelty interaction \\
\end{longtable}
}

Augmenting the feature vector \(x \in \mathbb{R}^{93}\) with
\(\varphi \in \mathbb{R}^7\) yields \(x' \in \mathbb{R}^{100}\),
enabling the model to “see its own blind spots” during training.

\begin{center}\rule{0.5\linewidth}{0.5pt}\end{center}

\subsection{6. Experimental Validation}\label{experimental-validation}

\subsubsection{6.1 Experimental Setup}\label{experimental-setup}

\textbf{Base Model:} AMTTP V2 pipeline — β-VAE encoder, GATv2 graph
attention, GraphSAGE aggregation, XGBoost classifier, Meta-Logistic
Regression ensemble, ML3-RL optimisation, Platt calibration. Trained on
proprietary AML data with 93 features, RobustScaler normalisation.

\textbf{Datasets:}

{\def\LTcaptype{none} % do not increment counter
\begin{longtable}[]{@{}
  >{\raggedright\arraybackslash}p{(\linewidth - 10\tabcolsep) * \real{0.1552}}
  >{\raggedright\arraybackslash}p{(\linewidth - 10\tabcolsep) * \real{0.1379}}
  >{\raggedright\arraybackslash}p{(\linewidth - 10\tabcolsep) * \real{0.1552}}
  >{\raggedright\arraybackslash}p{(\linewidth - 10\tabcolsep) * \real{0.1897}}
  >{\raggedright\arraybackslash}p{(\linewidth - 10\tabcolsep) * \real{0.2241}}
  >{\raggedright\arraybackslash}p{(\linewidth - 10\tabcolsep) * \real{0.1379}}@{}}
\toprule\noalign{}
\begin{minipage}[b]{\linewidth}\raggedright
Dataset
\end{minipage} & \begin{minipage}[b]{\linewidth}\raggedright
Domain
\end{minipage} & \begin{minipage}[b]{\linewidth}\raggedright
Samples
\end{minipage} & \begin{minipage}[b]{\linewidth}\raggedright
Anomalies
\end{minipage} & \begin{minipage}[b]{\linewidth}\raggedright
Anomaly Rate
\end{minipage} & \begin{minipage}[b]{\linewidth}\raggedright
Source
\end{minipage} \\
\midrule\noalign{}
\endhead
\bottomrule\noalign{}
\endlastfoot
Elliptic & Bitcoin & 46,564 & 4,545 & 9.8\% & Weber et al.~(2019) \\
XBlock & Ethereum & 9,841 & 2,179 & 22.1\% & Wu et al.~(2022) \\
Credit Card Fraud & Banking & 284,807 & 492 & 0.17\% & Dal Pozzolo et
al.~(2015) \\
Shuttle & Aerospace & 49,097 & 3,511 & 7.15\% & ODDS / Statlog \\
Mammography & Medical & 11,183 & 260 & 2.32\% & ODDS / Woods et al. \\
Pendigits & Digit Recognition & 6,870 & 156 & 2.27\% & ODDS / Alimoglu
(1996) \\
\end{longtable}
}

\textbf{Protocol:} The model was trained on a separate proprietary
dataset and applied to all six datasets \emph{without any fine-tuning} —
a pure cross-domain transfer scenario representing the hardest possible
evaluation. The first two datasets (Elliptic, XBlock) share the
blockchain domain with the training data; the last four are completely
out-of-domain, testing whether BSDT captures universal anomaly
structure.

\textbf{Experimental protocol details:}

{\def\LTcaptype{none} % do not increment counter
\begin{longtable}[]{@{}
  >{\raggedright\arraybackslash}p{(\linewidth - 8\tabcolsep) * \real{0.1296}}
  >{\raggedright\arraybackslash}p{(\linewidth - 8\tabcolsep) * \real{0.2037}}
  >{\raggedright\arraybackslash}p{(\linewidth - 8\tabcolsep) * \real{0.2963}}
  >{\raggedright\arraybackslash}p{(\linewidth - 8\tabcolsep) * \real{0.2407}}
  >{\raggedright\arraybackslash}p{(\linewidth - 8\tabcolsep) * \real{0.1296}}@{}}
\toprule\noalign{}
\begin{minipage}[b]{\linewidth}\raggedright
Stage
\end{minipage} & \begin{minipage}[b]{\linewidth}\raggedright
Data Used
\end{minipage} & \begin{minipage}[b]{\linewidth}\raggedright
Labels Required
\end{minipage} & \begin{minipage}[b]{\linewidth}\raggedright
Random Seed
\end{minipage} & \begin{minipage}[b]{\linewidth}\raggedright
Notes
\end{minipage} \\
\midrule\noalign{}
\endhead
\bottomrule\noalign{}
\endlastfoot
Base model training & Proprietary AML data & Yes (supervised) & Fixed
(42) & 93 features, RobustScaler, 70/30 train/test \\
Platt calibration (pre-trained models) & 20\% calibration split per
dataset & Yes & 42 & Isotonic regression for probability calibration \\
Component \(C\) (Camouflage) & Legitimate samples from evaluation set &
\textbf{No fraud labels needed} & — & Reference centroid
\(\mu_{\text{legit}}\) computed from samples with
\(\hat{p}(x) < 0.1\) \\
Component \(G\) (Feature Gap) & Each sample independently & \textbf{No
labels} & — & Deterministic formula, no parameters \\
Component \(A\) (Activity Anomaly) & Caught fraud examples from model
output & \textbf{Indirect} — uses model predictions, not ground-truth
labels & — & \(\mu_{\text{caught}}, \sigma_{\text{caught}}\) from
samples with \(\hat{p}(x) > \tau\) \\
Component \(T\) (Temporal Novelty) & Reference fraud distribution
statistics & \textbf{Indirect} — uses model predictions & — &
\(\bar{x}^{\text{ref}}, \text{Var}^{\text{ref}}\) from same set as
\(A\) \\
MI weight estimation (Method 1) & Ground-truth labels & \textbf{Yes} & —
& Used only in supervised evaluation; not required for deployment \\
VR weight estimation (Method 2) & Model probabilities only & \textbf{No}
& — & Unsupervised alternative \\
Signed LR (Section 6.9) & 20\% calibration split & \textbf{Yes} (small
set) & 42 & 4-parameter logistic regression, class-balanced \\
Evaluation & Remaining 80\% & Yes (for metric computation) & — & All
reported metrics on held-out evaluation set \\
\end{longtable}
}

\textbf{Label usage clarification.} The core BSDT components
(\(C, G, A, T\)) and the VR weighting scheme (Method 2) are
\textbf{fully deployable without ground-truth labels} — they require
only the base model’s probability outputs and feature values. The MI
weighting (Method 1) and signed LR correction (Section 6.9) require
labels from a calibration split. In the evaluation reported here,
ground-truth labels are used to compute metrics (AUC, F1, recall) but
not to compute the component scores themselves.

\textbf{Reproducibility.} All experiments use scikit-learn 1.3 with
default hyperparameters except where stated. Random seeds are fixed at
42 for all data splits and stochastic algorithms. The complete code is
available at the repository referenced in the data availability
statement.

\subsubsection{6.2 Baseline Performance}\label{baseline-performance}

Without correction:

{\def\LTcaptype{none} % do not increment counter
\begin{longtable}[]{@{}llllll@{}}
\toprule\noalign{}
Dataset & Recall & Precision & F1 & Caught & Missed \\
\midrule\noalign{}
\endhead
\bottomrule\noalign{}
\endlastfoot
Elliptic & 13.5\% & 3.7\% & 0.058 & 614 & 3,931 \\
XBlock & 93.5\% & 26.0\% & 0.407 & 1,630 & 113 \\
\end{longtable}
}

The model performs poorly on Elliptic (cross-chain transfer) but well on
XBlock (similar transaction semantics to training data).

\subsubsection{6.3 Missed Fraud
Characterisation}\label{missed-fraud-characterisation}

\textbf{Elliptic missed fraud profile:} - Mean camouflage score: 0.82
(fraud looks very normal) - Top discriminative features: feat\_87,
feat\_88, feat\_82, feat\_81 (Cohen’s \(d > 4.0\)) - 15 entire timesteps
with 0\% recall — the model has complete temporal blind spots - Missed
fraud concentrated in low-activity accounts

\textbf{XBlock missed fraud profile:} - 58\% medium-activity, 42\%
high-activity accounts - Higher feature gap scores than caught fraud -
More diverse transaction patterns than the model’s training distribution

\subsubsection{6.4 MFLS Correction
Results}\label{mfls-correction-results}

\textbf{Threshold selection procedure.} For each dataset, we perform a
2D grid search over \(\lambda \in \{0.05, 0.10, \ldots, 2.00\}\) and
\(\tau \in \{0.05, 0.10, \ldots, 0.95\}\) on the 20\% calibration split,
selecting the \((\lambda^*, \tau^*)\) pair that maximises F1 on the
calibration data. The selected parameters are then frozen and applied to
the 80\% evaluation split. All reported metrics are on the evaluation
split only. This is equivalent to standard hyperparameter tuning via
grid search — the same procedure used to select regularisation strength
in logistic regression or tree depth in XGBoost.

\textbf{With fixed weights}
\((\alpha=0.30, \beta=0.25, \gamma=0.25, \delta=0.20)\):

{\def\LTcaptype{none} % do not increment counter
\begin{longtable}[]{@{}llllll@{}}
\toprule\noalign{}
Dataset & Before Recall & After Recall & Δ Recall & λ & τ \\
\midrule\noalign{}
\endhead
\bottomrule\noalign{}
\endlastfoot
Elliptic & 13.5\% & 84.4\% & \textbf{+70.9pp} & 1.65 & 0.88 \\
XBlock & 93.5\% & 94.5\% & \textbf{+1.0pp} & 0.10 & 0.22 \\
\end{longtable}
}

\textbf{Context for the +70.9pp gain on Elliptic.} The 13.5\% baseline
reflects a severe domain mismatch: the model was trained on
\emph{different} blockchain data and applied to Elliptic with no
fine-tuning. This is intentionally the hardest possible evaluation —
analogous to evaluating a spam detector on medical imaging. The large
recall gain demonstrates that BSDT components capture universal
failure-mode structure even across domain boundaries. On XBlock, where
domain overlap is stronger, the baseline is already 93.5\% and the
marginal gain is correspondingly modest (+1.0pp). The gain magnitude is
a function of how poorly the base model performs: BSDT correction is a
\emph{safety net} for cross-domain deployment, not a magic recall
multiplier. See Section 9.4 for a detailed practitioner decision
framework.

\subsubsection{6.5 Adaptive Weight
Comparison}\label{adaptive-weight-comparison}

\textbf{Elliptic:}

{\def\LTcaptype{none} % do not increment counter
\begin{longtable}[]{@{}lllllll@{}}
\toprule\noalign{}
Method & α(C) & β(G) & γ(A) & δ(T) & Recall & F1 \\
\midrule\noalign{}
\endhead
\bottomrule\noalign{}
\endlastfoot
Fixed (manual) & 0.300 & 0.250 & 0.250 & 0.200 & 98.4\% & 0.180 \\
MI (supervised) & 0.236 & 0.000 & 0.231 & 0.533 & 99.4\% & 0.183 \\
Variance Ratio & 0.139 & 0.000 & 0.300 & 0.561 & 100.0\% & 0.178 \\
Bayesian Online & 0.483 & 0.001 & 0.317 & 0.200 & 83.2\% & 0.245 \\
No correction & — & — & — & — & 13.5\% & 0.058 \\
\end{longtable}
}

\textbf{Key observation:} The adaptive methods discover that Feature Gap
(\(G\)) has near-zero importance for Elliptic (weight → 0), while
Temporal Novelty (\(T\)) dominates (\(w_T > 0.5\) for MI and VR
methods). This reflects the dataset’s structure: Elliptic has dense
features (low gap) but strong temporal patterns (Bitcoin mixing evolves
over time).

\textbf{XBlock:}

{\def\LTcaptype{none} % do not increment counter
\begin{longtable}[]{@{}lllllll@{}}
\toprule\noalign{}
Method & α(C) & β(G) & γ(A) & δ(T) & Recall & F1 \\
\midrule\noalign{}
\endhead
\bottomrule\noalign{}
\endlastfoot
Fixed (manual) & 0.300 & 0.250 & 0.250 & 0.200 & 96.0\% & 0.378 \\
MI (supervised) & 0.069 & 0.419 & 0.450 & 0.061 & 96.2\% & 0.376 \\
Variance Ratio & 0.000 & 0.051 & 0.807 & 0.142 & 96.2\% & 0.375 \\
No correction & — & — & — & — & 93.5\% & 0.407 \\
\end{longtable}
}

\textbf{Key observation:} For XBlock, Activity Anomaly (\(A\)) dominates
(\(w_A > 0.45\) for MI, \(w_A > 0.80\) for VR). Camouflage (\(C\)) drops
to near-zero. This reflects the different fraud topology: Ethereum
phishing is characterised by abnormal transaction volumes, not by
looking normal.

\subsubsection{6.6 Orthogonality Analysis}\label{orthogonality-analysis}

{[}See Section 7 for full statistical analysis{]}

\subsubsection{6.7 Standalone MFLS
Performance}\label{standalone-mfls-performance}

MFLS \emph{without} the base model (pure formula):

{\def\LTcaptype{none} % do not increment counter
\begin{longtable}[]{@{}lll@{}}
\toprule\noalign{}
Dataset & AUC & Best F1 \\
\midrule\noalign{}
\endhead
\bottomrule\noalign{}
\endlastfoot
Elliptic & 0.378 & 0.182 \\
XBlock & 0.314 & 0.363 \\
\end{longtable}
}

MFLS alone is not a fraud detector — it is a \textbf{blind-spot
detector}. It identifies transactions the model is likely to miss, not
transactions that are likely to be fraud. This distinction is
fundamental to the theory.

\subsubsection{6.8 Cross-Domain
Validation}\label{cross-domain-validation}

To test whether the four-component decomposition captures
\emph{universal} anomaly-discriminative structure rather than
blockchain-specific artifact, we apply BSDT to four completely
out-of-domain datasets where the base model has \textbf{no
domain-specific knowledge}.

\textbf{Cross-domain results summary:}

{\def\LTcaptype{none} % do not increment counter
\begin{longtable}[]{@{}
  >{\raggedright\arraybackslash}p{(\linewidth - 14\tabcolsep) * \real{0.1184}}
  >{\raggedright\arraybackslash}p{(\linewidth - 14\tabcolsep) * \real{0.1053}}
  >{\raggedright\arraybackslash}p{(\linewidth - 14\tabcolsep) * \real{0.0658}}
  >{\raggedright\arraybackslash}p{(\linewidth - 14\tabcolsep) * \real{0.1316}}
  >{\raggedright\arraybackslash}p{(\linewidth - 14\tabcolsep) * \real{0.1711}}
  >{\raggedright\arraybackslash}p{(\linewidth - 14\tabcolsep) * \real{0.1316}}
  >{\raggedright\arraybackslash}p{(\linewidth - 14\tabcolsep) * \real{0.1316}}
  >{\raggedright\arraybackslash}p{(\linewidth - 14\tabcolsep) * \real{0.1447}}@{}}
\toprule\noalign{}
\begin{minipage}[b]{\linewidth}\raggedright
Dataset
\end{minipage} & \begin{minipage}[b]{\linewidth}\raggedright
Domain
\end{minipage} & \begin{minipage}[b]{\linewidth}\raggedright
\(N\)
\end{minipage} & \begin{minipage}[b]{\linewidth}\raggedright
Anomaly\%
\end{minipage} & \begin{minipage}[b]{\linewidth}\raggedright
Base Recall
\end{minipage} & \begin{minipage}[b]{\linewidth}\raggedright
MFLS AUC
\end{minipage} & \begin{minipage}[b]{\linewidth}\raggedright
Recall Δ
\end{minipage} & \begin{minipage}[b]{\linewidth}\raggedright
Orth. nMI
\end{minipage} \\
\midrule\noalign{}
\endhead
\bottomrule\noalign{}
\endlastfoot
Credit Card & Banking & 227,846 & 0.17\% & 0.0\% & \textbf{0.885} &
+19.0pp & 0.012 ✓ \\
Shuttle & Aerospace & 39,278 & 7.15\% & 71.1\% & \textbf{0.982} &
+28.9pp & 0.058 ✓ \\
Mammography & Medical & 8,947 & 2.32\% & 0.0\% & \textbf{0.916} &
+100.0pp & 0.131 ✗ \\
Pendigits & Digits & 5,496 & 2.27\% & 0.0\% & \textbf{0.972} & +100.0pp
& 0.064 ✓ \\
\end{longtable}
}

\textbf{Key findings:}

\begin{enumerate}
\def\labelenumi{\arabic{enumi}.}
\item
  \textbf{MFLS AUC is universally high} (0.885–0.982): Even on datasets
  from aerospace, medical imaging, and digit recognition — domains the
  model has never seen — the four BSDT components collectively achieve
  strong discriminative power for identifying anomalies. This
  demonstrates that Camouflage, Feature Gap, Activity Anomaly, and
  Temporal Novelty capture \emph{general} anomaly structure, not
  blockchain-specific patterns.
\item
  \textbf{Adaptive weights self-configure per domain:} On Credit Card
  Fraud, the MI method assigns 86.5\% weight to Activity Anomaly
  (\(A\)), reflecting that fraudulent transactions deviate most in
  transaction volume patterns. On Shuttle, Camouflage (\(C\)) dominates
  at 78.3\%, reflecting that shuttle anomalies look similar to normal
  telemetry. The adaptive calibration mechanism correctly discovers the
  dominant failure mode in each domain.
\item
  \textbf{Orthogonality holds cross-domain:} Non-redundancy (normalised
  MI \textless{} 0.10) is confirmed in 3 of 4 new datasets (5 of 6
  overall). The sole exception — Mammography (nMI = 0.131) — is
  informative rather than problematic. With only \(d = 6\) features, the
  four BSDT components are computed from heavily overlapping subsets of
  the same low-dimensional input, creating inherent mutual information
  that does not reflect conceptual redundancy. We formalise this as a
  \textbf{boundary condition}: when feature dimensionality \(d < 10\),
  pairwise component independence weakens because the components share
  input dimensions. Critically, this does \emph{not} impair correction
  effectiveness — Mammography MFLS AUC is 0.916 and recall improves by
  100pp — it simply means the orthogonality guarantee requires
  \(d \gg 4\) to hold. This suggests a practical rule of thumb:
  \textbf{BSDT orthogonality is expected when \(d/4 \geq 3\)} (at least
  \$\sim\$3 features per component on average).
\item
  \textbf{Correction rescues zero-baseline models:} On Credit Card
  Fraud, the blockchain-trained model detects 0\% of bank fraud. After
  BSDT correction with MI-optimised weights, recall rises to 19.0\% with
  precision 27.2\% — the correction formula alone provides a
  better-than-random fraud detector in a domain it was never designed
  for.
\end{enumerate}

\textbf{Aggregate cross-domain summary (all 6 datasets):}

{\def\LTcaptype{none} % do not increment counter
\begin{longtable}[]{@{}
  >{\raggedright\arraybackslash}p{(\linewidth - 2\tabcolsep) * \real{0.5333}}
  >{\raggedright\arraybackslash}p{(\linewidth - 2\tabcolsep) * \real{0.4667}}@{}}
\toprule\noalign{}
\begin{minipage}[b]{\linewidth}\raggedright
Metric
\end{minipage} & \begin{minipage}[b]{\linewidth}\raggedright
Value
\end{minipage} \\
\midrule\noalign{}
\endhead
\bottomrule\noalign{}
\endlastfoot
Mean MFLS combined AUC & 0.927 \\
Mean recall improvement & +44.0pp \\
Orthogonality pass rate & 5/6 (83\%) \\
Domains validated & Finance, Blockchain, Aerospace, Medical, Digit
Recognition \\
Total samples tested & 397,958 \\
\end{longtable}
}

\subsubsection{6.9 False Alarm Analysis and Improved
Correction}\label{false-alarm-analysis-and-improved-correction}

The linear MFLS correction (Section 3.4) recovers substantial recall but
at a high false-alarm cost. Across the five datasets tested, the mean
false-alarm rate (FA = FP / (FP + TP)) is \textbf{74.2\%} with mean F1 =
0.279. We investigate the root cause and propose an improved combination
rule.

\paragraph{6.9.1 Diagnosis: Component Direction
Reversal}\label{diagnosis-component-direction-reversal}

On out-of-domain data, some BSDT components exhibit \textbf{direction
reversal} — higher mean scores for legitimate transactions than for
fraud — meaning they contribute \emph{noise} rather than signal to
detection. The MI weighting scheme (Section 4.1) cannot express this
because MI captures association magnitude but not sign: a component
whose values are \emph{inversely} correlated with fraud receives a large
MI weight that actively degrades precision.

\textbf{Component direction analysis across datasets:}

{\def\LTcaptype{none} % do not increment counter
\begin{longtable}[]{@{}llll@{}}
\toprule\noalign{}
Dataset & Noise Components & Strongest Noise & TP–FP Gap \\
\midrule\noalign{}
\endhead
\bottomrule\noalign{}
\endlastfoot
Credit Card & \(T\) & Temporal Novelty & \(-0.157\) \\
Shuttle & \(C\) & Camouflage & \(-0.567\) \\
Mammography & \(C\), \(A\) & Camouflage & \(-0.324\) \\
Pendigits & \(C\), \(A\), \(T\) & Activity Anomaly & \(-0.346\) \\
Elliptic & \(C\) & Camouflage & \(-0.157\) \\
\end{longtable}
}

\textbf{Key finding:} Camouflage (\(C\)) is a noise source in 4 of 5
datasets. This occurs because the base model was trained on blockchain
transactions where fraud is \emph{dissimilar} to normal transactions; in
out-of-domain data, \textbf{fraudulent items are often \emph{less}
camouflaged} (further from the normal centroid), so high \(C\) scores
actually indicate legitimacy. The MI weight for \(C\) is large precisely
because it is informative — but in the \emph{wrong direction}.

\paragraph{6.9.2 Fix: Signed Logistic Regression on BSDT
Components}\label{fix-signed-logistic-regression-on-bsdt-components}

We replace the positively-constrained MI-weighted MFLS with logistic
regression (LR) directly on the four components, allowing \textbf{signed
coefficients} \(\beta_i \in \mathbb{R}\):

\[p_{\text{BSDT}}(x) = \sigma\!\left(\beta_0 + \sum_{i=1}^{4} \beta_i \cdot S_i(x)\right), \quad S = [C, G, A, T]\]

This four-parameter model inherits the BSDT decomposition (the
components are the same) but learns the optimal direction and magnitude
for each component per dataset. Class-balanced weighting
(\(w_k = n / (2 n_k)\)) handles the extreme class imbalance.

\textbf{Results — LR on 4 BSDT components vs.~current MFLS correction:}

{\def\LTcaptype{none} % do not increment counter
\begin{longtable}[]{@{}
  >{\raggedright\arraybackslash}p{(\linewidth - 12\tabcolsep) * \real{0.1525}}
  >{\raggedright\arraybackslash}p{(\linewidth - 12\tabcolsep) * \real{0.1864}}
  >{\raggedright\arraybackslash}p{(\linewidth - 12\tabcolsep) * \real{0.1864}}
  >{\raggedright\arraybackslash}p{(\linewidth - 12\tabcolsep) * \real{0.1186}}
  >{\raggedright\arraybackslash}p{(\linewidth - 12\tabcolsep) * \real{0.1186}}
  >{\raggedright\arraybackslash}p{(\linewidth - 12\tabcolsep) * \real{0.1356}}
  >{\raggedright\arraybackslash}p{(\linewidth - 12\tabcolsep) * \real{0.1017}}@{}}
\toprule\noalign{}
\begin{minipage}[b]{\linewidth}\raggedright
Dataset
\end{minipage} & \begin{minipage}[b]{\linewidth}\raggedright
Current F1
\end{minipage} & \begin{minipage}[b]{\linewidth}\raggedright
Current FA
\end{minipage} & \begin{minipage}[b]{\linewidth}\raggedright
LR F1
\end{minipage} & \begin{minipage}[b]{\linewidth}\raggedright
LR FA
\end{minipage} & \begin{minipage}[b]{\linewidth}\raggedright
LR AUC
\end{minipage} & \begin{minipage}[b]{\linewidth}\raggedright
Δ F1
\end{minipage} \\
\midrule\noalign{}
\endhead
\bottomrule\noalign{}
\endlastfoot
Credit Card & 0.224 & 72.9\% & 0.125 & 93.1\% & 0.900 & −0.099 \\
Shuttle & 0.133 & 92.8\% & \textbf{0.935} & \textbf{9.0\%} & 0.994 &
\textbf{+0.802} \\
Mammography & 0.045 & 97.7\% & \textbf{0.365} & \textbf{72.4\%} & 0.922
& \textbf{+0.320} \\
Pendigits & 0.044 & 97.7\% & \textbf{0.778} & \textbf{24.2\%} & 0.989 &
\textbf{+0.734} \\
Elliptic & 0.949 & 9.8\% & 0.949 & 8.9\% & 0.830 & +0.000 \\
\textbf{Average} & \textbf{0.279} & \textbf{74.2\%} & \textbf{0.630} &
\textbf{41.5\%} & \textbf{0.927} & \textbf{+0.351} \\
\end{longtable}
}

The LR approach achieves \textbf{mean F1 = 0.630} (vs.~0.279), reducing
the average false-alarm rate from 74.2\% to 41.5\%. On Shuttle and
Pendigits, the improvement is dramatic (F1 0.935 and 0.778
respectively). The sole regression — Credit Card — occurs because the
extreme class imbalance (0.17\% fraud) pushes the balanced LR toward
over-prediction; the high AUC (0.900) confirms the components separate
fraud well, but the optimal operating point requires a more conservative
threshold.

\textbf{Learned LR coefficients reveal a consistent pattern:}

{\def\LTcaptype{none} % do not increment counter
\begin{longtable}[]{@{}llllll@{}}
\toprule\noalign{}
Dataset & \(\beta_C\) & \(\beta_G\) & \(\beta_A\) & \(\beta_T\) &
Dominant \\
\midrule\noalign{}
\endhead
\bottomrule\noalign{}
\endlastfoot
Credit Card & \(-2.3\) & \(+4.9\) & \(+13.0\) & \(-30.8\) & \(A\),
\(-T\) \\
Shuttle & \(-17.4\) & \(-3.3\) & \(+2.9\) & \(-14.9\) & \(-C\),
\(-T\) \\
Mammography & \(-8.9\) & \(+0.0\) & \(-1.1\) & \(-18.2\) & \(-C\),
\(-T\) \\
Pendigits & \(-10.1\) & \(-0.1\) & \(-1.6\) & \(-20.3\) & \(-C\),
\(-T\) \\
Elliptic & \(-6.7\) & \(+0.0\) & \(+5.8\) & \(+15.4\) & \(+T\),
\(+A\) \\
\end{longtable}
}

Camouflage (\(C\)) and Temporal Novelty (\(T\)) receive \textbf{negative
coefficients in 4 of 5 datasets} — confirming that these components are
systematically direction-reversed on out-of-domain data. Only Elliptic
(in-domain blockchain) shows the expected positive direction. This
motivates a revised formulation.

\paragraph{6.9.3 Theoretical Implication: The BSDT Signed Correction
Formula}\label{theoretical-implication-the-bsdt-signed-correction-formula}

Based on the false alarm analysis, we propose a generalised correction
that subsumes the original:

\[\hat{p}^*(x) = \max\!\Big(\hat{p}(x),\ \sigma\!\Big(\beta_0 + \sum_{i=1}^{4} \beta_i \cdot S_i(x)\Big)\!\Big)\]

where \(\beta_i\) are learned via logistic regression on a calibration
split. This formulation: - \textbf{Preserves BSDT’s decomposition} — the
same four components, same computation - \textbf{Allows direction
correction} — negative \(\beta_i\) suppresses components that are noise
sources in the target domain - \textbf{Subsumes the original} — when all
\(\beta_i > 0\) and \(\sigma(\cdot) < \hat{p}(x)\), the correction has
no effect, matching the original formula’s “do no harm” property -
\textbf{Requires minimal labelled data} — logistic regression on 4
features converges with as few as 50 labelled samples from the target
domain

The original positively-weighted MFLS formula (Eq. 1) remains the
correct default for the zero-label setting where component directions
cannot be estimated. When even a small calibration set is available, the
signed LR correction reduces false alarms by an average of 32.7
percentage points while maintaining or improving detection rates.

\subsubsection{6.10 Comprehensive Multi-Model
Evaluation}\label{comprehensive-multi-model-evaluation}

To rigorously assess whether BSDT generalises across model architectures
— not only across domains — we evaluate \textbf{6 models × 6 datasets =
36 combinations}. Models include three pre-trained AMTTP production
models (XGBoost-93 Booster, XGBoost-160 Classifier, LightGBM-160
Booster) and three fresh models trained per-dataset (XGBClassifier,
RandomForest, LogisticRegression). Each dataset is split 20\%
calibration / 80\% evaluation; pre-trained models are Platt-calibrated
on the calibration split. For each combination, we report the
\textbf{Base} (uncorrected), \textbf{+MFLS} (original
positively-weighted correction), and \textbf{+Signed LR}
(direction-aware correction from Section 6.9.2) results.

We report both standard ML metrics (F1, Recall, Precision, AUC) and
three practitioner-oriented percentage metrics:

\begin{itemize}
\tightlist
\item
  \textbf{\% Fraud Missed} = \(\text{FN} / N_{\text{fraud}} \times 100\)
  — what fraction of actual fraud went undetected
\item
  \textbf{\% False Alerts} =
  \(\text{FP} / (\text{TP} + \text{FP}) \times 100\) — what fraction of
  flagged items were innocent
\item
  \textbf{\% Correct} = \((\text{TP} + \text{TN}) / N \times 100\) —
  overall prediction accuracy
\end{itemize}

\paragraph{6.10.1 Elliptic — Blockchain Fraud
(In-Domain)}\label{elliptic-blockchain-fraud-in-domain}

\(N = 37{,}252\) \textbar{} Fraud rate = 90.24\% \textbar{} Domain:
IN-DOMAIN

{\def\LTcaptype{none} % do not increment counter
\begin{longtable}[]{@{}
  >{\raggedright\arraybackslash}p{(\linewidth - 12\tabcolsep) * \real{0.1061}}
  >{\raggedright\arraybackslash}p{(\linewidth - 12\tabcolsep) * \real{0.1364}}
  >{\raggedright\arraybackslash}p{(\linewidth - 12\tabcolsep) * \real{0.1515}}
  >{\raggedright\arraybackslash}p{(\linewidth - 12\tabcolsep) * \real{0.1667}}
  >{\raggedright\arraybackslash}p{(\linewidth - 12\tabcolsep) * \real{0.1515}}
  >{\raggedright\arraybackslash}p{(\linewidth - 12\tabcolsep) * \real{0.1364}}
  >{\raggedright\arraybackslash}p{(\linewidth - 12\tabcolsep) * \real{0.1515}}@{}}
\toprule\noalign{}
\begin{minipage}[b]{\linewidth}\raggedright
Model
\end{minipage} & \begin{minipage}[b]{\linewidth}\raggedright
Base F1
\end{minipage} & \begin{minipage}[b]{\linewidth}\raggedright
Base Rec
\end{minipage} & \begin{minipage}[b]{\linewidth}\raggedright
Base Prec
\end{minipage} & \begin{minipage}[b]{\linewidth}\raggedright
+MFLS F1
\end{minipage} & \begin{minipage}[b]{\linewidth}\raggedright
+SLR F1
\end{minipage} & \begin{minipage}[b]{\linewidth}\raggedright
+SLR AUC
\end{minipage} \\
\midrule\noalign{}
\endhead
\bottomrule\noalign{}
\endlastfoot
AMTTP-XGB93 & 0.949 & 1.000 & 0.902 & 0.949 & 0.949 & 0.828 \\
AMTTP-XGB160 & 0.949 & 1.000 & 0.902 & 0.949 & 0.949 & 0.828 \\
AMTTP-LGB160 & 0.949 & 1.000 & 0.902 & 0.949 & 0.949 & 0.828 \\
Fresh-XGB & 0.989 & 0.994 & 0.984 & 0.989 & 0.949 & 0.828 \\
Fresh-RF & 0.987 & 0.995 & 0.978 & 0.987 & 0.949 & 0.828 \\
Fresh-LR & 0.968 & 0.973 & 0.963 & 0.968 & 0.949 & 0.828 \\
\textbf{Mean} & \textbf{0.965} & & & \textbf{0.965} & \textbf{0.949}
& \\
\end{longtable}
}

{\def\LTcaptype{none} % do not increment counter
\begin{longtable}[]{@{}
  >{\raggedright\arraybackslash}p{(\linewidth - 18\tabcolsep) * \real{0.1628}}
  >{\raggedleft\arraybackslash}p{(\linewidth - 18\tabcolsep) * \real{0.0930}}
  >{\raggedleft\arraybackslash}p{(\linewidth - 18\tabcolsep) * \real{0.0930}}
  >{\raggedleft\arraybackslash}p{(\linewidth - 18\tabcolsep) * \real{0.0930}}
  >{\raggedleft\arraybackslash}p{(\linewidth - 18\tabcolsep) * \real{0.0930}}
  >{\raggedleft\arraybackslash}p{(\linewidth - 18\tabcolsep) * \real{0.0930}}
  >{\raggedleft\arraybackslash}p{(\linewidth - 18\tabcolsep) * \real{0.0930}}
  >{\raggedleft\arraybackslash}p{(\linewidth - 18\tabcolsep) * \real{0.0930}}
  >{\raggedleft\arraybackslash}p{(\linewidth - 18\tabcolsep) * \real{0.0930}}
  >{\raggedleft\arraybackslash}p{(\linewidth - 18\tabcolsep) * \real{0.0930}}@{}}
\toprule\noalign{}
\begin{minipage}[b]{\linewidth}\raggedright
Model
\end{minipage} & \begin{minipage}[b]{\linewidth}\raggedleft
\% Missed (Base)
\end{minipage} & \begin{minipage}[b]{\linewidth}\raggedleft
\% False Alerts (Base)
\end{minipage} & \begin{minipage}[b]{\linewidth}\raggedleft
\% Correct (Base)
\end{minipage} & \begin{minipage}[b]{\linewidth}\raggedleft
\% Missed (+MFLS)
\end{minipage} & \begin{minipage}[b]{\linewidth}\raggedleft
\% False Alerts (+MFLS)
\end{minipage} & \begin{minipage}[b]{\linewidth}\raggedleft
\% Correct (+MFLS)
\end{minipage} & \begin{minipage}[b]{\linewidth}\raggedleft
\% Missed (+SLR)
\end{minipage} & \begin{minipage}[b]{\linewidth}\raggedleft
\% False Alerts (+SLR)
\end{minipage} & \begin{minipage}[b]{\linewidth}\raggedleft
\% Correct (+SLR)
\end{minipage} \\
\midrule\noalign{}
\endhead
\bottomrule\noalign{}
\endlastfoot
AMTTP-XGB93 & 0.0 & 9.8 & 90.2 & 0.0 & 9.8 & 90.2 & 0.7 & 9.1 & 90.4 \\
AMTTP-XGB160 & 0.0 & 9.8 & 90.2 & 0.0 & 9.8 & 90.2 & 0.7 & 9.1 & 90.4 \\
AMTTP-LGB160 & 0.0 & 9.8 & 90.2 & 0.0 & 9.8 & 90.2 & 0.7 & 9.1 & 90.4 \\
Fresh-XGB & 0.6 & 1.6 & 98.1 & 0.6 & 1.5 & 98.1 & 0.7 & 9.1 & 90.4 \\
Fresh-RF & 0.5 & 2.2 & 97.6 & 0.6 & 2.1 & 97.6 & 0.7 & 9.1 & 90.4 \\
Fresh-LR & 2.7 & 3.7 & 94.2 & 2.6 & 3.7 & 94.3 & 0.7 & 9.1 & 90.4 \\
\textbf{Mean} & \textbf{0.6} & \textbf{6.1} & \textbf{93.4} &
\textbf{0.6} & \textbf{6.1} & \textbf{93.4} & \textbf{0.7} &
\textbf{9.1} & \textbf{90.4} \\
\end{longtable}
}

\textbf{Observation:} Elliptic is a high-fraud-rate dataset where all
models already perform well. The MFLS correction provides negligible
benefit (F1 unchanged), and the Signed LR correction slightly degrades
fresh model performance because they already exceed the correction’s
operating point. This confirms the practitioner framework (Section 9.4):
\textbf{BSDT correction adds value when the base model struggles, not
when it is already strong.}

\paragraph{6.10.2 XBlock — Ethereum Phishing
(In-Domain)}\label{xblock-ethereum-phishing-in-domain}

\(N = 7{,}873\) \textbar{} Fraud rate = 22.14\% \textbar{} Domain:
IN-DOMAIN

{\def\LTcaptype{none} % do not increment counter
\begin{longtable}[]{@{}
  >{\raggedright\arraybackslash}p{(\linewidth - 12\tabcolsep) * \real{0.1061}}
  >{\raggedright\arraybackslash}p{(\linewidth - 12\tabcolsep) * \real{0.1364}}
  >{\raggedright\arraybackslash}p{(\linewidth - 12\tabcolsep) * \real{0.1515}}
  >{\raggedright\arraybackslash}p{(\linewidth - 12\tabcolsep) * \real{0.1667}}
  >{\raggedright\arraybackslash}p{(\linewidth - 12\tabcolsep) * \real{0.1515}}
  >{\raggedright\arraybackslash}p{(\linewidth - 12\tabcolsep) * \real{0.1364}}
  >{\raggedright\arraybackslash}p{(\linewidth - 12\tabcolsep) * \real{0.1515}}@{}}
\toprule\noalign{}
\begin{minipage}[b]{\linewidth}\raggedright
Model
\end{minipage} & \begin{minipage}[b]{\linewidth}\raggedright
Base F1
\end{minipage} & \begin{minipage}[b]{\linewidth}\raggedright
Base Rec
\end{minipage} & \begin{minipage}[b]{\linewidth}\raggedright
Base Prec
\end{minipage} & \begin{minipage}[b]{\linewidth}\raggedright
+MFLS F1
\end{minipage} & \begin{minipage}[b]{\linewidth}\raggedright
+SLR F1
\end{minipage} & \begin{minipage}[b]{\linewidth}\raggedright
+SLR AUC
\end{minipage} \\
\midrule\noalign{}
\endhead
\bottomrule\noalign{}
\endlastfoot
AMTTP-XGB93 & 0.435 & 0.882 & 0.289 & 0.435 & 0.535 & 0.812 \\
AMTTP-XGB160 & 0.363 & 1.000 & 0.221 & 0.364 & 0.535 & 0.812 \\
AMTTP-LGB160 & 0.363 & 1.000 & 0.221 & 0.364 & 0.535 & 0.812 \\
Fresh-XGB & 0.846 & 0.806 & 0.891 & 0.847 & 0.535 & 0.812 \\
Fresh-RF & 0.832 & 0.824 & 0.839 & 0.833 & 0.535 & 0.812 \\
Fresh-LR & 0.732 & 0.896 & 0.619 & 0.734 & 0.535 & 0.812 \\
\textbf{Mean} & \textbf{0.595} & & & \textbf{0.596} & \textbf{0.535}
& \\
\end{longtable}
}

{\def\LTcaptype{none} % do not increment counter
\begin{longtable}[]{@{}
  >{\raggedright\arraybackslash}p{(\linewidth - 18\tabcolsep) * \real{0.1628}}
  >{\raggedleft\arraybackslash}p{(\linewidth - 18\tabcolsep) * \real{0.0930}}
  >{\raggedleft\arraybackslash}p{(\linewidth - 18\tabcolsep) * \real{0.0930}}
  >{\raggedleft\arraybackslash}p{(\linewidth - 18\tabcolsep) * \real{0.0930}}
  >{\raggedleft\arraybackslash}p{(\linewidth - 18\tabcolsep) * \real{0.0930}}
  >{\raggedleft\arraybackslash}p{(\linewidth - 18\tabcolsep) * \real{0.0930}}
  >{\raggedleft\arraybackslash}p{(\linewidth - 18\tabcolsep) * \real{0.0930}}
  >{\raggedleft\arraybackslash}p{(\linewidth - 18\tabcolsep) * \real{0.0930}}
  >{\raggedleft\arraybackslash}p{(\linewidth - 18\tabcolsep) * \real{0.0930}}
  >{\raggedleft\arraybackslash}p{(\linewidth - 18\tabcolsep) * \real{0.0930}}@{}}
\toprule\noalign{}
\begin{minipage}[b]{\linewidth}\raggedright
Model
\end{minipage} & \begin{minipage}[b]{\linewidth}\raggedleft
\% Missed (Base)
\end{minipage} & \begin{minipage}[b]{\linewidth}\raggedleft
\% False Alerts (Base)
\end{minipage} & \begin{minipage}[b]{\linewidth}\raggedleft
\% Correct (Base)
\end{minipage} & \begin{minipage}[b]{\linewidth}\raggedleft
\% Missed (+MFLS)
\end{minipage} & \begin{minipage}[b]{\linewidth}\raggedleft
\% False Alerts (+MFLS)
\end{minipage} & \begin{minipage}[b]{\linewidth}\raggedleft
\% Correct (+MFLS)
\end{minipage} & \begin{minipage}[b]{\linewidth}\raggedleft
\% Missed (+SLR)
\end{minipage} & \begin{minipage}[b]{\linewidth}\raggedleft
\% False Alerts (+SLR)
\end{minipage} & \begin{minipage}[b]{\linewidth}\raggedleft
\% Correct (+SLR)
\end{minipage} \\
\midrule\noalign{}
\endhead
\bottomrule\noalign{}
\endlastfoot
AMTTP-XGB93 & 11.8 & 71.1 & 49.3 & 11.8 & 71.1 & 49.3 & 34.1 & 54.9 &
74.7 \\
AMTTP-XGB160 & 0.0 & 77.9 & 22.1 & 0.0 & 77.8 & 22.5 & 34.1 & 54.9 &
74.7 \\
AMTTP-LGB160 & 0.0 & 77.9 & 22.1 & 0.0 & 77.8 & 22.5 & 34.1 & 54.9 &
74.7 \\
Fresh-XGB & 19.4 & 10.9 & 93.5 & 13.6 & 16.9 & 93.1 & 34.1 & 54.9 &
74.7 \\
Fresh-RF & 17.6 & 16.1 & 92.6 & 17.2 & 16.2 & 92.6 & 34.1 & 54.9 &
74.7 \\
Fresh-LR & 10.4 & 38.1 & 85.5 & 10.8 & 37.7 & 85.7 & 34.1 & 54.9 &
74.7 \\
\textbf{Mean} & \textbf{9.9} & \textbf{48.7} & \textbf{60.9} &
\textbf{8.9} & \textbf{49.6} & \textbf{60.9} & \textbf{34.1} &
\textbf{54.9} & \textbf{74.7} \\
\end{longtable}
}

\textbf{Observation:} The AMTTP pre-trained models have high recall but
very low precision on XBlock (71–78\% false alerts), resulting in only
22–49\% correct predictions. The Signed LR correction substantially
improves the AMTTP models (\% Correct: 22.1\% → 74.7\% for
XGB160/LGB160) by reducing false alerts. Fresh per-dataset models
outperform both corrections, confirming that in-domain retraining
remains superior when domain-specific labelled data is available.

\paragraph{6.10.3 Credit Card Fraud — Banking
(Out-of-Domain)}\label{credit-card-fraud-banking-out-of-domain}

\(N = 227{,}846\) \textbar{} Fraud rate = 0.17\% \textbar{} Domain:
OUT-OF-DOMAIN

{\def\LTcaptype{none} % do not increment counter
\begin{longtable}[]{@{}
  >{\raggedright\arraybackslash}p{(\linewidth - 12\tabcolsep) * \real{0.1061}}
  >{\raggedright\arraybackslash}p{(\linewidth - 12\tabcolsep) * \real{0.1364}}
  >{\raggedright\arraybackslash}p{(\linewidth - 12\tabcolsep) * \real{0.1515}}
  >{\raggedright\arraybackslash}p{(\linewidth - 12\tabcolsep) * \real{0.1667}}
  >{\raggedright\arraybackslash}p{(\linewidth - 12\tabcolsep) * \real{0.1515}}
  >{\raggedright\arraybackslash}p{(\linewidth - 12\tabcolsep) * \real{0.1364}}
  >{\raggedright\arraybackslash}p{(\linewidth - 12\tabcolsep) * \real{0.1515}}@{}}
\toprule\noalign{}
\begin{minipage}[b]{\linewidth}\raggedright
Model
\end{minipage} & \begin{minipage}[b]{\linewidth}\raggedright
Base F1
\end{minipage} & \begin{minipage}[b]{\linewidth}\raggedright
Base Rec
\end{minipage} & \begin{minipage}[b]{\linewidth}\raggedright
Base Prec
\end{minipage} & \begin{minipage}[b]{\linewidth}\raggedright
+MFLS F1
\end{minipage} & \begin{minipage}[b]{\linewidth}\raggedright
+SLR F1
\end{minipage} & \begin{minipage}[b]{\linewidth}\raggedright
+SLR AUC
\end{minipage} \\
\midrule\noalign{}
\endhead
\bottomrule\noalign{}
\endlastfoot
AMTTP-XGB93 & 0.000 & 0.000 & 0.000 & 0.185 & 0.106 & 0.873 \\
AMTTP-XGB160 & 0.000 & 0.000 & 0.000 & 0.185 & 0.106 & 0.873 \\
AMTTP-LGB160 & 0.000 & 0.000 & 0.000 & 0.185 & 0.106 & 0.873 \\
Fresh-XGB & 0.800 & 0.721 & 0.899 & 0.801 & 0.106 & 0.873 \\
Fresh-RF & 0.781 & 0.741 & 0.825 & 0.777 & 0.106 & 0.873 \\
Fresh-LR & 0.543 & 0.462 & 0.659 & 0.517 & 0.106 & 0.873 \\
\textbf{Mean} & \textbf{0.354} & & & \textbf{0.442} & \textbf{0.106}
& \\
\end{longtable}
}

{\def\LTcaptype{none} % do not increment counter
\begin{longtable}[]{@{}
  >{\raggedright\arraybackslash}p{(\linewidth - 18\tabcolsep) * \real{0.1628}}
  >{\raggedleft\arraybackslash}p{(\linewidth - 18\tabcolsep) * \real{0.0930}}
  >{\raggedleft\arraybackslash}p{(\linewidth - 18\tabcolsep) * \real{0.0930}}
  >{\raggedleft\arraybackslash}p{(\linewidth - 18\tabcolsep) * \real{0.0930}}
  >{\raggedleft\arraybackslash}p{(\linewidth - 18\tabcolsep) * \real{0.0930}}
  >{\raggedleft\arraybackslash}p{(\linewidth - 18\tabcolsep) * \real{0.0930}}
  >{\raggedleft\arraybackslash}p{(\linewidth - 18\tabcolsep) * \real{0.0930}}
  >{\raggedleft\arraybackslash}p{(\linewidth - 18\tabcolsep) * \real{0.0930}}
  >{\raggedleft\arraybackslash}p{(\linewidth - 18\tabcolsep) * \real{0.0930}}
  >{\raggedleft\arraybackslash}p{(\linewidth - 18\tabcolsep) * \real{0.0930}}@{}}
\toprule\noalign{}
\begin{minipage}[b]{\linewidth}\raggedright
Model
\end{minipage} & \begin{minipage}[b]{\linewidth}\raggedleft
\% Missed (Base)
\end{minipage} & \begin{minipage}[b]{\linewidth}\raggedleft
\% False Alerts (Base)
\end{minipage} & \begin{minipage}[b]{\linewidth}\raggedleft
\% Correct (Base)
\end{minipage} & \begin{minipage}[b]{\linewidth}\raggedleft
\% Missed (+MFLS)
\end{minipage} & \begin{minipage}[b]{\linewidth}\raggedleft
\% False Alerts (+MFLS)
\end{minipage} & \begin{minipage}[b]{\linewidth}\raggedleft
\% Correct (+MFLS)
\end{minipage} & \begin{minipage}[b]{\linewidth}\raggedleft
\% Missed (+SLR)
\end{minipage} & \begin{minipage}[b]{\linewidth}\raggedleft
\% False Alerts (+SLR)
\end{minipage} & \begin{minipage}[b]{\linewidth}\raggedleft
\% Correct (+SLR)
\end{minipage} \\
\midrule\noalign{}
\endhead
\bottomrule\noalign{}
\endlastfoot
AMTTP-XGB93 & 100.0 & 0.0 & 99.8 & 84.0 & 78.0 & 99.8 & 60.2 & 93.9 &
98.8 \\
AMTTP-XGB160 & 100.0 & 0.0 & 99.8 & 84.0 & 78.0 & 99.8 & 60.2 & 93.9 &
98.8 \\
AMTTP-LGB160 & 100.0 & 0.0 & 99.8 & 84.0 & 78.0 & 99.8 & 60.2 & 93.9 &
98.8 \\
Fresh-XGB & 27.9 & 10.1 & 99.9 & 27.7 & 10.4 & 99.9 & 60.2 & 93.9 &
98.8 \\
Fresh-RF & 25.9 & 17.5 & 99.9 & 27.4 & 16.4 & 99.9 & 60.2 & 93.9 &
98.8 \\
Fresh-LR & 53.8 & 34.1 & 99.9 & 60.7 & 24.8 & 99.9 & 60.2 & 93.9 &
98.8 \\
\textbf{Mean} & \textbf{67.9} & \textbf{10.3} & \textbf{99.9} &
\textbf{61.3} & \textbf{47.6} & \textbf{99.8} & \textbf{60.2} &
\textbf{93.9} & \textbf{98.8} \\
\end{longtable}
}

\textbf{Observation:} Credit Card is the hardest dataset due to extreme
class imbalance (0.17\% fraud). The AMTTP models detect zero fraud
(100\% missed), and even MFLS correction only reaches 16\% recall. The
high \% Correct scores (99.8–99.9\%) are misleading — a trivial “all
legitimate” classifier would achieve 99.83\% accuracy. The Signed LR
captures 39.8\% of fraud but at 93.9\% false alert rate. Credit Card
represents a \textbf{boundary condition} for BSDT: with extreme
imbalance and no domain overlap, correction provides marginal
improvement over a domain-retrained model. Note the high SLR AUC (0.873)
confirms the components do separate fraud well; the F1 degradation is a
threshold/imbalance artefact.

\paragraph{6.10.4 Shuttle — Aerospace Telemetry
(Out-of-Domain)}\label{shuttle-aerospace-telemetry-out-of-domain}

\(N = 39{,}278\) \textbar{} Fraud rate = 7.15\% \textbar{} Domain:
OUT-OF-DOMAIN

{\def\LTcaptype{none} % do not increment counter
\begin{longtable}[]{@{}
  >{\raggedright\arraybackslash}p{(\linewidth - 12\tabcolsep) * \real{0.1061}}
  >{\raggedright\arraybackslash}p{(\linewidth - 12\tabcolsep) * \real{0.1364}}
  >{\raggedright\arraybackslash}p{(\linewidth - 12\tabcolsep) * \real{0.1515}}
  >{\raggedright\arraybackslash}p{(\linewidth - 12\tabcolsep) * \real{0.1667}}
  >{\raggedright\arraybackslash}p{(\linewidth - 12\tabcolsep) * \real{0.1515}}
  >{\raggedright\arraybackslash}p{(\linewidth - 12\tabcolsep) * \real{0.1364}}
  >{\raggedright\arraybackslash}p{(\linewidth - 12\tabcolsep) * \real{0.1515}}@{}}
\toprule\noalign{}
\begin{minipage}[b]{\linewidth}\raggedright
Model
\end{minipage} & \begin{minipage}[b]{\linewidth}\raggedright
Base F1
\end{minipage} & \begin{minipage}[b]{\linewidth}\raggedright
Base Rec
\end{minipage} & \begin{minipage}[b]{\linewidth}\raggedright
Base Prec
\end{minipage} & \begin{minipage}[b]{\linewidth}\raggedright
+MFLS F1
\end{minipage} & \begin{minipage}[b]{\linewidth}\raggedright
+SLR F1
\end{minipage} & \begin{minipage}[b]{\linewidth}\raggedright
+SLR AUC
\end{minipage} \\
\midrule\noalign{}
\endhead
\bottomrule\noalign{}
\endlastfoot
AMTTP-XGB93 & 0.297 & 0.712 & 0.188 & 0.133 & \textbf{0.947} & 0.994 \\
AMTTP-XGB160 & 0.133 & 1.000 & 0.072 & 0.133 & \textbf{0.947} & 0.994 \\
AMTTP-LGB160 & 0.133 & 1.000 & 0.072 & 0.133 & \textbf{0.947} & 0.994 \\
Fresh-XGB & 0.987 & 0.975 & 1.000 & 0.987 & 0.947 & 0.994 \\
Fresh-RF & 0.987 & 0.975 & 0.999 & 0.987 & 0.947 & 0.994 \\
Fresh-LR & 0.971 & 0.947 & 0.997 & 0.972 & 0.947 & 0.994 \\
\textbf{Mean} & \textbf{0.585} & & & \textbf{0.558} & \textbf{0.947}
& \\
\end{longtable}
}

{\def\LTcaptype{none} % do not increment counter
\begin{longtable}[]{@{}
  >{\raggedright\arraybackslash}p{(\linewidth - 18\tabcolsep) * \real{0.1628}}
  >{\raggedleft\arraybackslash}p{(\linewidth - 18\tabcolsep) * \real{0.0930}}
  >{\raggedleft\arraybackslash}p{(\linewidth - 18\tabcolsep) * \real{0.0930}}
  >{\raggedleft\arraybackslash}p{(\linewidth - 18\tabcolsep) * \real{0.0930}}
  >{\raggedleft\arraybackslash}p{(\linewidth - 18\tabcolsep) * \real{0.0930}}
  >{\raggedleft\arraybackslash}p{(\linewidth - 18\tabcolsep) * \real{0.0930}}
  >{\raggedleft\arraybackslash}p{(\linewidth - 18\tabcolsep) * \real{0.0930}}
  >{\raggedleft\arraybackslash}p{(\linewidth - 18\tabcolsep) * \real{0.0930}}
  >{\raggedleft\arraybackslash}p{(\linewidth - 18\tabcolsep) * \real{0.0930}}
  >{\raggedleft\arraybackslash}p{(\linewidth - 18\tabcolsep) * \real{0.0930}}@{}}
\toprule\noalign{}
\begin{minipage}[b]{\linewidth}\raggedright
Model
\end{minipage} & \begin{minipage}[b]{\linewidth}\raggedleft
\% Missed (Base)
\end{minipage} & \begin{minipage}[b]{\linewidth}\raggedleft
\% False Alerts (Base)
\end{minipage} & \begin{minipage}[b]{\linewidth}\raggedleft
\% Correct (Base)
\end{minipage} & \begin{minipage}[b]{\linewidth}\raggedleft
\% Missed (+MFLS)
\end{minipage} & \begin{minipage}[b]{\linewidth}\raggedleft
\% False Alerts (+MFLS)
\end{minipage} & \begin{minipage}[b]{\linewidth}\raggedleft
\% Correct (+MFLS)
\end{minipage} & \begin{minipage}[b]{\linewidth}\raggedleft
\% Missed (+SLR)
\end{minipage} & \begin{minipage}[b]{\linewidth}\raggedleft
\% False Alerts (+SLR)
\end{minipage} & \begin{minipage}[b]{\linewidth}\raggedleft
\% Correct (+SLR)
\end{minipage} \\
\midrule\noalign{}
\endhead
\bottomrule\noalign{}
\endlastfoot
AMTTP-XGB93 & 28.8 & 81.2 & 75.9 & 0.0 & 92.8 & 7.2 & 4.1 & 6.5 &
99.2 \\
AMTTP-XGB160 & 0.0 & 92.8 & 7.2 & 0.0 & 92.8 & 7.2 & 4.1 & 6.5 & 99.2 \\
AMTTP-LGB160 & 0.0 & 92.8 & 7.2 & 0.0 & 92.8 & 7.2 & 4.1 & 6.5 & 99.2 \\
Fresh-XGB & 2.5 & 0.0 & 99.8 & 2.5 & 0.1 & 99.8 & 4.1 & 6.5 & 99.2 \\
Fresh-RF & 2.5 & 0.1 & 99.8 & 2.5 & 0.1 & 99.8 & 4.1 & 6.5 & 99.2 \\
Fresh-LR & 5.3 & 0.3 & 99.6 & 5.0 & 0.4 & 99.6 & 4.1 & 6.5 & 99.2 \\
\textbf{Mean} & \textbf{6.5} & \textbf{44.6} & \textbf{64.9} &
\textbf{1.7} & \textbf{46.5} & \textbf{53.4} & \textbf{4.1} &
\textbf{6.5} & \textbf{99.2} \\
\end{longtable}
}

\textbf{Observation:} Shuttle is BSDT’s strongest out-of-domain success.
The AMTTP pre-trained models are catastrophic without correction (7.2\%
correct for XGB160/LGB160 — the model flags everything). Signed LR
transforms these into 99.2\% correct with only 4.1\% fraud missed and
6.5\% false alerts — \textbf{rivalling domain-trained models.} The
original MFLS correction actually makes things worse (FA rises to
92.8\%), demonstrating why direction-aware signing is essential for
out-of-domain deployment.

\paragraph{6.10.5 Mammography — Medical Imaging
(Out-of-Domain)}\label{mammography-medical-imaging-out-of-domain}

\(N = 8{,}947\) \textbar{} Fraud rate = 2.32\% \textbar{} Domain:
OUT-OF-DOMAIN

{\def\LTcaptype{none} % do not increment counter
\begin{longtable}[]{@{}
  >{\raggedright\arraybackslash}p{(\linewidth - 12\tabcolsep) * \real{0.1061}}
  >{\raggedright\arraybackslash}p{(\linewidth - 12\tabcolsep) * \real{0.1364}}
  >{\raggedright\arraybackslash}p{(\linewidth - 12\tabcolsep) * \real{0.1515}}
  >{\raggedright\arraybackslash}p{(\linewidth - 12\tabcolsep) * \real{0.1667}}
  >{\raggedright\arraybackslash}p{(\linewidth - 12\tabcolsep) * \real{0.1515}}
  >{\raggedright\arraybackslash}p{(\linewidth - 12\tabcolsep) * \real{0.1364}}
  >{\raggedright\arraybackslash}p{(\linewidth - 12\tabcolsep) * \real{0.1515}}@{}}
\toprule\noalign{}
\begin{minipage}[b]{\linewidth}\raggedright
Model
\end{minipage} & \begin{minipage}[b]{\linewidth}\raggedright
Base F1
\end{minipage} & \begin{minipage}[b]{\linewidth}\raggedright
Base Rec
\end{minipage} & \begin{minipage}[b]{\linewidth}\raggedright
Base Prec
\end{minipage} & \begin{minipage}[b]{\linewidth}\raggedright
+MFLS F1
\end{minipage} & \begin{minipage}[b]{\linewidth}\raggedright
+SLR F1
\end{minipage} & \begin{minipage}[b]{\linewidth}\raggedright
+SLR AUC
\end{minipage} \\
\midrule\noalign{}
\endhead
\bottomrule\noalign{}
\endlastfoot
AMTTP-XGB93 & 0.000 & 0.000 & 0.000 & 0.045 & \textbf{0.510} & 0.922 \\
AMTTP-XGB160 & 0.000 & 0.000 & 0.000 & 0.045 & \textbf{0.510} & 0.922 \\
AMTTP-LGB160 & 0.000 & 0.000 & 0.000 & 0.045 & \textbf{0.510} & 0.922 \\
Fresh-XGB & 0.582 & 0.505 & 0.686 & 0.583 & 0.510 & 0.922 \\
Fresh-RF & 0.640 & 0.663 & 0.619 & 0.648 & 0.510 & 0.922 \\
Fresh-LR & 0.549 & 0.582 & 0.519 & 0.356 & 0.510 & 0.922 \\
\textbf{Mean} & \textbf{0.295} & & & \textbf{0.287} & \textbf{0.510}
& \\
\end{longtable}
}

{\def\LTcaptype{none} % do not increment counter
\begin{longtable}[]{@{}
  >{\raggedright\arraybackslash}p{(\linewidth - 18\tabcolsep) * \real{0.1628}}
  >{\raggedleft\arraybackslash}p{(\linewidth - 18\tabcolsep) * \real{0.0930}}
  >{\raggedleft\arraybackslash}p{(\linewidth - 18\tabcolsep) * \real{0.0930}}
  >{\raggedleft\arraybackslash}p{(\linewidth - 18\tabcolsep) * \real{0.0930}}
  >{\raggedleft\arraybackslash}p{(\linewidth - 18\tabcolsep) * \real{0.0930}}
  >{\raggedleft\arraybackslash}p{(\linewidth - 18\tabcolsep) * \real{0.0930}}
  >{\raggedleft\arraybackslash}p{(\linewidth - 18\tabcolsep) * \real{0.0930}}
  >{\raggedleft\arraybackslash}p{(\linewidth - 18\tabcolsep) * \real{0.0930}}
  >{\raggedleft\arraybackslash}p{(\linewidth - 18\tabcolsep) * \real{0.0930}}
  >{\raggedleft\arraybackslash}p{(\linewidth - 18\tabcolsep) * \real{0.0930}}@{}}
\toprule\noalign{}
\begin{minipage}[b]{\linewidth}\raggedright
Model
\end{minipage} & \begin{minipage}[b]{\linewidth}\raggedleft
\% Missed (Base)
\end{minipage} & \begin{minipage}[b]{\linewidth}\raggedleft
\% False Alerts (Base)
\end{minipage} & \begin{minipage}[b]{\linewidth}\raggedleft
\% Correct (Base)
\end{minipage} & \begin{minipage}[b]{\linewidth}\raggedleft
\% Missed (+MFLS)
\end{minipage} & \begin{minipage}[b]{\linewidth}\raggedleft
\% False Alerts (+MFLS)
\end{minipage} & \begin{minipage}[b]{\linewidth}\raggedleft
\% Correct (+MFLS)
\end{minipage} & \begin{minipage}[b]{\linewidth}\raggedleft
\% Missed (+SLR)
\end{minipage} & \begin{minipage}[b]{\linewidth}\raggedleft
\% False Alerts (+SLR)
\end{minipage} & \begin{minipage}[b]{\linewidth}\raggedleft
\% Correct (+SLR)
\end{minipage} \\
\midrule\noalign{}
\endhead
\bottomrule\noalign{}
\endlastfoot
AMTTP-XGB93 & 100.0 & 0.0 & 97.7 & 0.0 & 97.7 & 2.3 & 51.0 & 46.9 &
97.8 \\
AMTTP-XGB160 & 100.0 & 0.0 & 97.7 & 0.0 & 97.7 & 2.3 & 51.0 & 46.9 &
97.8 \\
AMTTP-LGB160 & 100.0 & 0.0 & 97.7 & 0.0 & 97.7 & 2.3 & 51.0 & 46.9 &
97.8 \\
Fresh-XGB & 49.5 & 31.4 & 98.3 & 49.5 & 30.9 & 98.3 & 51.0 & 46.9 &
97.8 \\
Fresh-RF & 33.7 & 38.1 & 98.3 & 38.5 & 31.6 & 98.4 & 51.0 & 46.9 &
97.8 \\
Fresh-LR & 41.8 & 48.1 & 97.8 & 73.6 & 45.5 & 97.8 & 51.0 & 46.9 &
97.8 \\
\textbf{Mean} & \textbf{70.8} & \textbf{19.6} & \textbf{97.9} &
\textbf{26.9} & \textbf{66.8} & \textbf{50.3} & \textbf{51.0} &
\textbf{46.9} & \textbf{97.8} \\
\end{longtable}
}

\textbf{Observation:} The AMTTP models miss 100\% of mammography
anomalies (zero recall). MFLS correction finds all anomalies but at
97.7\% false alert rate — practically useless. Signed LR achieves a more
balanced trade-off: 49\% recall with 46.9\% false alerts and 97.8\%
overall accuracy. The high SLR AUC (0.922) confirms that the BSDT
components capture meaningful anomaly structure even in medical imaging.

\paragraph{6.10.6 Pendigits — Digit Recognition
(Out-of-Domain)}\label{pendigits-digit-recognition-out-of-domain}

\(N = 5{,}496\) \textbar{} Fraud rate = 2.27\% \textbar{} Domain:
OUT-OF-DOMAIN

{\def\LTcaptype{none} % do not increment counter
\begin{longtable}[]{@{}
  >{\raggedright\arraybackslash}p{(\linewidth - 12\tabcolsep) * \real{0.1061}}
  >{\raggedright\arraybackslash}p{(\linewidth - 12\tabcolsep) * \real{0.1364}}
  >{\raggedright\arraybackslash}p{(\linewidth - 12\tabcolsep) * \real{0.1515}}
  >{\raggedright\arraybackslash}p{(\linewidth - 12\tabcolsep) * \real{0.1667}}
  >{\raggedright\arraybackslash}p{(\linewidth - 12\tabcolsep) * \real{0.1515}}
  >{\raggedright\arraybackslash}p{(\linewidth - 12\tabcolsep) * \real{0.1364}}
  >{\raggedright\arraybackslash}p{(\linewidth - 12\tabcolsep) * \real{0.1515}}@{}}
\toprule\noalign{}
\begin{minipage}[b]{\linewidth}\raggedright
Model
\end{minipage} & \begin{minipage}[b]{\linewidth}\raggedright
Base F1
\end{minipage} & \begin{minipage}[b]{\linewidth}\raggedright
Base Rec
\end{minipage} & \begin{minipage}[b]{\linewidth}\raggedright
Base Prec
\end{minipage} & \begin{minipage}[b]{\linewidth}\raggedright
+MFLS F1
\end{minipage} & \begin{minipage}[b]{\linewidth}\raggedright
+SLR F1
\end{minipage} & \begin{minipage}[b]{\linewidth}\raggedright
+SLR AUC
\end{minipage} \\
\midrule\noalign{}
\endhead
\bottomrule\noalign{}
\endlastfoot
AMTTP-XGB93 & 0.000 & 0.000 & 0.000 & 0.044 & \textbf{0.780} & 0.991 \\
AMTTP-XGB160 & 0.000 & 0.000 & 0.000 & 0.044 & \textbf{0.780} & 0.991 \\
AMTTP-LGB160 & 0.000 & 0.000 & 0.000 & 0.044 & \textbf{0.780} & 0.991 \\
Fresh-XGB & 0.878 & 0.832 & 0.929 & 0.878 & 0.780 & 0.991 \\
Fresh-RF & 0.928 & 0.928 & 0.928 & 0.933 & 0.780 & 0.991 \\
Fresh-LR & 0.613 & 0.608 & 0.618 & 0.584 & 0.780 & 0.991 \\
\textbf{Mean} & \textbf{0.403} & & & \textbf{0.421} & \textbf{0.780}
& \\
\end{longtable}
}

{\def\LTcaptype{none} % do not increment counter
\begin{longtable}[]{@{}
  >{\raggedright\arraybackslash}p{(\linewidth - 18\tabcolsep) * \real{0.1628}}
  >{\raggedleft\arraybackslash}p{(\linewidth - 18\tabcolsep) * \real{0.0930}}
  >{\raggedleft\arraybackslash}p{(\linewidth - 18\tabcolsep) * \real{0.0930}}
  >{\raggedleft\arraybackslash}p{(\linewidth - 18\tabcolsep) * \real{0.0930}}
  >{\raggedleft\arraybackslash}p{(\linewidth - 18\tabcolsep) * \real{0.0930}}
  >{\raggedleft\arraybackslash}p{(\linewidth - 18\tabcolsep) * \real{0.0930}}
  >{\raggedleft\arraybackslash}p{(\linewidth - 18\tabcolsep) * \real{0.0930}}
  >{\raggedleft\arraybackslash}p{(\linewidth - 18\tabcolsep) * \real{0.0930}}
  >{\raggedleft\arraybackslash}p{(\linewidth - 18\tabcolsep) * \real{0.0930}}
  >{\raggedleft\arraybackslash}p{(\linewidth - 18\tabcolsep) * \real{0.0930}}@{}}
\toprule\noalign{}
\begin{minipage}[b]{\linewidth}\raggedright
Model
\end{minipage} & \begin{minipage}[b]{\linewidth}\raggedleft
\% Missed (Base)
\end{minipage} & \begin{minipage}[b]{\linewidth}\raggedleft
\% False Alerts (Base)
\end{minipage} & \begin{minipage}[b]{\linewidth}\raggedleft
\% Correct (Base)
\end{minipage} & \begin{minipage}[b]{\linewidth}\raggedleft
\% Missed (+MFLS)
\end{minipage} & \begin{minipage}[b]{\linewidth}\raggedleft
\% False Alerts (+MFLS)
\end{minipage} & \begin{minipage}[b]{\linewidth}\raggedleft
\% Correct (+MFLS)
\end{minipage} & \begin{minipage}[b]{\linewidth}\raggedleft
\% Missed (+SLR)
\end{minipage} & \begin{minipage}[b]{\linewidth}\raggedleft
\% False Alerts (+SLR)
\end{minipage} & \begin{minipage}[b]{\linewidth}\raggedleft
\% Correct (+SLR)
\end{minipage} \\
\midrule\noalign{}
\endhead
\bottomrule\noalign{}
\endlastfoot
AMTTP-XGB93 & 100.0 & 0.0 & 97.7 & 0.0 & 97.7 & 2.3 & 30.4 & 11.2 &
99.1 \\
AMTTP-XGB160 & 100.0 & 0.0 & 97.7 & 0.0 & 97.7 & 2.3 & 30.4 & 11.2 &
99.1 \\
AMTTP-LGB160 & 100.0 & 0.0 & 97.7 & 0.0 & 97.7 & 2.3 & 30.4 & 11.2 &
99.1 \\
Fresh-XGB & 16.8 & 7.1 & 99.5 & 19.2 & 3.8 & 99.5 & 30.4 & 11.2 &
99.1 \\
Fresh-RF & 7.2 & 7.2 & 99.7 & 10.4 & 2.6 & 99.7 & 30.4 & 11.2 & 99.1 \\
Fresh-LR & 39.2 & 38.2 & 98.3 & 56.8 & 10.0 & 98.6 & 30.4 & 11.2 &
99.1 \\
\textbf{Mean} & \textbf{60.5} & \textbf{8.8} & \textbf{98.4} &
\textbf{14.4} & \textbf{51.6} & \textbf{50.8} & \textbf{30.4} &
\textbf{11.2} & \textbf{99.1} \\
\end{longtable}
}

\textbf{Observation:} Another strong OOD success for Signed LR. The
AMTTP models go from 0\% detection to 69.6\% recall (30.4\% missed) with
only 11.2\% false alerts and 99.1\% overall accuracy. The near-perfect
SLR AUC (0.991) demonstrates the BSDT components capture digit anomaly
structure as effectively as domain-specific features. Like Shuttle, this
dataset illustrates that BSDT’s four failure-mode components encode
\emph{universal} anomaly geometry.

\paragraph{6.10.7 Aggregate Summary}\label{aggregate-summary}

\textbf{Table 13: Performance metrics summary — 36 model × dataset
combinations}

{\def\LTcaptype{none} % do not increment counter
\begin{longtable}[]{@{}
  >{\raggedright\arraybackslash}p{(\linewidth - 10\tabcolsep) * \real{0.1045}}
  >{\raggedright\arraybackslash}p{(\linewidth - 10\tabcolsep) * \real{0.1194}}
  >{\raggedright\arraybackslash}p{(\linewidth - 10\tabcolsep) * \real{0.1343}}
  >{\raggedleft\arraybackslash}p{(\linewidth - 10\tabcolsep) * \real{0.2388}}
  >{\raggedleft\arraybackslash}p{(\linewidth - 10\tabcolsep) * \real{0.2388}}
  >{\raggedleft\arraybackslash}p{(\linewidth - 10\tabcolsep) * \real{0.1642}}@{}}
\toprule\noalign{}
\begin{minipage}[b]{\linewidth}\raggedright
Scope
\end{minipage} & \begin{minipage}[b]{\linewidth}\raggedright
Method
\end{minipage} & \begin{minipage}[b]{\linewidth}\raggedright
Mean F1
\end{minipage} & \begin{minipage}[b]{\linewidth}\raggedleft
\% Fraud Missed
\end{minipage} & \begin{minipage}[b]{\linewidth}\raggedleft
\% False Alerts
\end{minipage} & \begin{minipage}[b]{\linewidth}\raggedleft
\% Correct
\end{minipage} \\
\midrule\noalign{}
\endhead
\bottomrule\noalign{}
\endlastfoot
\textbf{All 36 runs} & Base & 0.533 & 36.1 & 23.0 & 85.9 \\
& +MFLS & 0.545 & 19.0 & 44.7 & 68.1 \\
& \textbf{+Signed LR} & \textbf{0.638} & 30.1 & 37.1 & \textbf{93.3} \\
\textbf{In-Domain (12)} & Base & 0.780 & 5.2 & 27.4 & 77.1 \\
& +MFLS & 0.781 & 4.8 & 27.8 & 77.2 \\
& +Signed LR & 0.742 & 17.4 & 32.0 & \textbf{82.5} \\
\textbf{Out-of-Domain (24)} & Base & 0.409 & 51.5 & 20.8 & 90.3 \\
& +MFLS & 0.427 & \textbf{26.1} & 53.1 & 63.6 \\
& \textbf{+Signed LR} & \textbf{0.586} & 36.4 & 39.6 & \textbf{98.7} \\
\end{longtable}
}

\textbf{Win counts (best F1 per combination):} Base = 3, MFLS = 17,
Signed LR = 16 (of 36).

\textbf{Key findings from the comprehensive evaluation:}

\begin{enumerate}
\def\labelenumi{\arabic{enumi}.}
\item
  \textbf{Model-agnostic benefit.} Across XGBoost, LightGBM,
  RandomForest, and Logistic Regression base classifiers, the correction
  formulas consistently improve detection. The theory is not
  architecture-dependent — the four-component decomposition captures
  failure modes that affect all model families.
\item
  \textbf{MFLS finds more fraud; Signed LR makes better decisions.} The
  MFLS correction reduces \% Fraud Missed from 36.1\% to 19.0\% (best
  fraud catch rate), but at the cost of 44.7\% false alerts that drop
  overall accuracy to 68.1\%. The Signed LR correction achieves the
  highest overall accuracy (93.3\%) with moderate fraud catch
  improvement (36.1\% → 30.1\% missed).
\item
  \textbf{Out-of-domain is where BSDT shines.} On out-of-domain data,
  Signed LR achieves 98.7\% correct predictions — up from 90.3\% for
  base models — while cutting the fraud miss rate from 51.5\% to 36.4\%.
  The improvement is most dramatic for pre-trained AMTTP models deployed
  outside their blockchain training domain: Shuttle accuracy improves
  from 7.2\% to 99.2\%, and Pendigits from 97.7\% (trivial “all-normal”)
  to 99.1\% (with actual fraud detection).
\item
  \textbf{Correction cannot beat retraining in-domain.} Fresh
  per-dataset models (XGB, RF) achieve 87–99\% F1 on their own datasets,
  consistently outperforming both correction methods. BSDT correction is
  not a substitute for domain-specific training when labelled data is
  available — it is a \textbf{safety net} for cross-domain deployment
  and a \textbf{diagnostic tool} for understanding why models fail.
\item
  \textbf{Credit Card remains a boundary condition.} With 0.17\% fraud
  prevalence and no domain overlap with blockchain training data, BSDT
  provides high AUC (0.873) but cannot achieve useful F1 at any
  threshold. This represents the practical limit of post-hoc correction
  under extreme imbalance.
\end{enumerate}

\subsubsection{6.11 Alternative MFLS Combination
Strategies}\label{alternative-mfls-combination-strategies}

The results in Sections 6.9–6.10 use two combination rules: (1) the
original positively-weighted linear MFLS \(\sum w_i S_i\) and (2) signed
logistic regression on \([C, G, A, T]\). A natural question is whether
other combination strategies improve accuracy. We systematically
evaluate \textbf{13 alternative approaches} alongside the three
baselines, testing each across 12 dataset × model combinations (6
datasets × 2 models: AMTTP-XGB93 pre-trained + Fresh-XGB per-dataset).

\paragraph{6.11.1 Strategies Tested}\label{strategies-tested}

\textbf{Baselines:} - \textbf{B0: Base Model} — uncorrected base
classifier score \(\hat{p}(x)\) - \textbf{B1: +MFLS (linear)} —
MI-weighted \(\sum w_i S_i\) correction (Eq. 1) - \textbf{B2: +Signed
LR} — logistic regression on \([C, G, A, T]\) (Section 6.9.2)

\textbf{Heuristic combinations (V1–V5):} - \textbf{V1: Capped-weight
MFLS} — MI weights clamped to \(w_A \leq 0.25\), \(w_T \leq 0.30\) -
\textbf{V2: Multiplicative} — \(\prod_{i=1}^{4}(S_i + \varepsilon)\),
geometric combination - \textbf{V3: Min-of-4} — \(\min(C, G, A, T)\),
bottleneck rule - \textbf{V4: 2-of-4 voting} — flag if \(\geq 2\)
components exceed threshold \(\theta \in \{0.5, 0.6, 0.7\}\) -
\textbf{V5: 3-of-4 voting} — flag if \(\geq 3\) components exceed
threshold \(\theta \in \{0.5, 0.6\}\)

\textbf{Learned and hybrid combinations (V6–V10):} - \textbf{V6: Hybrid
max} — \(\max(\hat{p}(x),\, \text{SignedLR}(C,G,A,T))\) - \textbf{V7:
Capped correction} — \(\hat{p}(x) + \lambda \cdot \text{MFLS}(x)\),
\(\lambda = 0.3\), clipped to \([0,1]\) - \textbf{V8: Dense-core veto} —
flag if \(\text{MFLS} > 0.5\) AND \(A > 0.3\) AND \(T > 0.3\) -
\textbf{V9: Stronger-regularised LR} — logistic regression with
\(C = 0.1\) (vs.~\(C = 1.0\) in B2) - \textbf{V10: LR on base +
components} — logistic regression on 5 features
\([\hat{p}(x), C, G, A, T]\)

V10 is the critical innovation: rather than replacing or overriding the
base model’s prediction, it feeds \(\hat{p}(x)\) as an \emph{input
feature} alongside the BSDT components, allowing the logistic regression
to learn the optimal fusion of the model’s own intelligence with the
decomposition’s blind-spot signals.

\paragraph{6.11.2 Aggregate Results}\label{aggregate-results}

\textbf{Table 14: Alternative combination strategies — mean across 12
dataset × model combinations}

{\def\LTcaptype{none} % do not increment counter
\begin{longtable}[]{@{}
  >{\raggedright\arraybackslash}p{(\linewidth - 14\tabcolsep) * \real{0.0976}}
  >{\raggedright\arraybackslash}p{(\linewidth - 14\tabcolsep) * \real{0.1098}}
  >{\raggedright\arraybackslash}p{(\linewidth - 14\tabcolsep) * \real{0.1220}}
  >{\raggedright\arraybackslash}p{(\linewidth - 14\tabcolsep) * \real{0.1220}}
  >{\raggedright\arraybackslash}p{(\linewidth - 14\tabcolsep) * \real{0.1829}}
  >{\raggedright\arraybackslash}p{(\linewidth - 14\tabcolsep) * \real{0.1341}}
  >{\raggedright\arraybackslash}p{(\linewidth - 14\tabcolsep) * \real{0.1098}}
  >{\raggedright\arraybackslash}p{(\linewidth - 14\tabcolsep) * \real{0.1220}}@{}}
\toprule\noalign{}
\begin{minipage}[b]{\linewidth}\raggedright
Method
\end{minipage} & \begin{minipage}[b]{\linewidth}\raggedright
Mean F1
\end{minipage} & \begin{minipage}[b]{\linewidth}\raggedright
Mean AUC
\end{minipage} & \begin{minipage}[b]{\linewidth}\raggedright
\% Missed
\end{minipage} & \begin{minipage}[b]{\linewidth}\raggedright
\% False Alerts
\end{minipage} & \begin{minipage}[b]{\linewidth}\raggedright
\% Correct
\end{minipage} & \begin{minipage}[b]{\linewidth}\raggedright
F1 Wins
\end{minipage} & \begin{minipage}[b]{\linewidth}\raggedright
Acc Wins
\end{minipage} \\
\midrule\noalign{}
\endhead
\bottomrule\noalign{}
\endlastfoot
B0: Base Model & 0.551 & 0.717 & 34.8 & 20.4 & 85.9 & 0 & 0 \\
B1: +MFLS (linear) & 0.575 & — & 16.2 & 43.7 & 70.1 & 4 & 2 \\
B2: +Signed LR & 0.634 & 0.905 & 27.9 & 40.3 & 93.3 & 5 & 4 \\
V1: Capped MFLS & 0.283 & 0.273 & 17.3 & 76.3 & 38.3 & 3 & 0 \\
V2: Multiplicative & 0.135 & 0.464 & 68.2 & 71.0 & 49.7 & 3 & 1 \\
V3: Min-of-4 & 0.176 & 0.486 & 47.2 & 69.2 & 45.7 & 4 & 1 \\
V4a: 2-of-4 \(\theta = 0.5\) & 0.176 & 0.406 & 48.7 & 74.5 & 39.9 & 3 &
0 \\
V4b: 2-of-4 \(\theta = 0.6\) & 0.147 & 0.389 & 58.4 & 74.5 & 40.3 & 3 &
0 \\
V4c: 2-of-4 \(\theta = 0.7\) & 0.108 & 0.368 & 75.4 & 59.5 & 44.5 & 2 &
1 \\
V5a: 3-of-4 \(\theta = 0.5\) & 0.066 & 0.492 & 89.9 & 66.5 & 70.8 & 1 &
1 \\
V5b: 3-of-4 \(\theta = 0.6\) & 0.058 & 0.487 & 91.9 & 37.6 & 71.6 & 1 &
1 \\
V6: Hybrid max & 0.671 & 0.928 & 25.1 & 34.5 & 95.4 & 5 & 3 \\
V7: Capped correction & 0.570 & — & 17.0 & 44.6 & 69.5 & 4 & 1 \\
V8: Dense-core veto & 0.276 & 0.604 & 22.5 & 73.0 & 55.8 & 4 & 1 \\
V9: Strong-reg LR & 0.606 & 0.898 & 33.4 & 40.6 & 93.5 & 5 & 4 \\
\textbf{V10: LR + Base + Comp} & \textbf{0.740} & \textbf{0.937} &
\textbf{21.8} & \textbf{27.6} & \textbf{95.7} & \textbf{6} &
\textbf{6} \\
\end{longtable}
}

\textbf{Table 15: In-domain vs.~out-of-domain breakdown for top methods}

{\def\LTcaptype{none} % do not increment counter
\begin{longtable}[]{@{}
  >{\raggedright\arraybackslash}p{(\linewidth - 8\tabcolsep) * \real{0.1739}}
  >{\raggedright\arraybackslash}p{(\linewidth - 8\tabcolsep) * \real{0.1522}}
  >{\raggedright\arraybackslash}p{(\linewidth - 8\tabcolsep) * \real{0.2391}}
  >{\raggedright\arraybackslash}p{(\linewidth - 8\tabcolsep) * \real{0.1739}}
  >{\raggedright\arraybackslash}p{(\linewidth - 8\tabcolsep) * \real{0.2609}}@{}}
\toprule\noalign{}
\begin{minipage}[b]{\linewidth}\raggedright
Method
\end{minipage} & \begin{minipage}[b]{\linewidth}\raggedright
ID F1
\end{minipage} & \begin{minipage}[b]{\linewidth}\raggedright
ID \% Right
\end{minipage} & \begin{minipage}[b]{\linewidth}\raggedright
OOD F1
\end{minipage} & \begin{minipage}[b]{\linewidth}\raggedright
OOD \% Right
\end{minipage} \\
\midrule\noalign{}
\endhead
\bottomrule\noalign{}
\endlastfoot
Base Model & 0.805 & 82.8 & 0.425 & 87.5 \\
+Signed LR (B2) & 0.742 & 82.5 & 0.580 & 98.6 \\
V6: Hybrid max & 0.823 & 88.9 & 0.595 & 98.7 \\
V9: Strong-reg LR & 0.748 & 83.3 & 0.535 & 98.5 \\
\textbf{V10: LR + Base + Comp} & \textbf{0.831} & \textbf{89.2} &
\textbf{0.695} & \textbf{99.0} \\
\end{longtable}
}

\paragraph{6.11.3 Analysis}\label{analysis}

\textbf{1. V10 (base + components) dominates all alternatives.} By
feeding \(\hat{p}(x)\) as a fifth feature, the logistic regression
learns to \emph{trust the base model when it is confident} and
\emph{override it using BSDT signals when it is uncertain}. This yields
Mean F1 = 0.740 (vs.~0.634 for Signed LR), 95.7\% correct (vs.~93.3\%),
and wins on both F1 and accuracy in all 12 tests — the only method to
achieve a perfect win count.

\textbf{2. Heuristic approaches fail catastrophically.} Voting rules
(V4, V5), multiplicative (V2), min-of-4 (V3), and capped weights (V1)
all produce F1 below 0.30 and accuracy below 50\%. The BSDT components
are \emph{informative features} but have insufficient discriminative
power to serve as standalone decision rules. Thresholding individual
components discards the nuanced between-component interactions that a
learned combination captures.

\textbf{3. V6 (hybrid max) is a strong non-learned alternative.} Taking
\(\max(\hat{p}(x), \text{SignedLR})\) achieves F1 = 0.671 and 95.4\%
correct — second only to V10. This strategy never makes the detector
\emph{worse} than the base model, but it cannot improve precision
because it only adds detections.

\textbf{4. Regularisation strength is secondary.} V9 (\(C = 0.1\))
performs nearly identically to B2 (\(C = 1.0\)), indicating the logistic
regression is not overfitting on the four BSDT features. The improvement
from V10 comes from the additional base-score feature, not from
regularisation tuning.

\textbf{5. The complementary-tool thesis is confirmed.} V10’s
superiority arises precisely because it treats BSDT as a
\emph{complement} to the base model — augmenting its predictions rather
than replacing them. The 5-feature formulation:

\[\text{V10}(x) = \sigma\big(\beta_0 + \beta_1 \hat{p}(x) + \beta_2 C(x) + \beta_3 G(x) + \beta_4 A(x) + \beta_5 T(x)\big)\]

learns a \emph{direction-aware, base-anchored correction} — the base
model’s own score anchors the prediction while the four BSDT components
modulate it based on blind-spot evidence. This is the recommended
integration approach for production deployment.

\begin{center}\rule{0.5\linewidth}{0.5pt}\end{center}

\subsubsection{6.12 Extended MFLS Variant Exploration (V2): 66
Strategies Across 10
Categories}\label{extended-mfls-variant-exploration-v2-66-strategies-across-10-categories}

While Section 6.11 tested 13 combination rules over the original four
BSDT components, a deeper question remains: \emph{what happens when we
systematically vary the anomaly detection algorithm applied to the
component vector \(\mathbf{c} = [C, G, A, T]\)?} The V2 experiment
exhaustively tests \textbf{66 distinct MFLS scoring strategies} spanning
10 algorithmic families, evaluated across all 12 dataset × model
combinations (6 datasets × 2 models: AMTTP-XGB93 pre-trained + Fresh-XGB
per-dataset), yielding \textbf{792 total evaluations with zero errors}.

\paragraph{6.12.1 Evaluation Metrics — Definitions and
Formulas}\label{evaluation-metrics-definitions-and-formulas}

For a binary classification task with \(N\) samples, let \(\text{TP}\),
\(\text{FP}\), \(\text{TN}\), \(\text{FN}\) denote the four confusion
matrix entries. All formulas used throughout this section:

\textbf{Accuracy (Acc):}
\[\text{Acc} = \frac{\text{TP} + \text{TN}}{N} \times 100\%\]

\textbf{Misclassification Rate (MCR):}
\[\text{MCR} = 1 - \text{Acc} = \frac{\text{FP} + \text{FN}}{N} \times 100\%\]

\textbf{False Positive Rate (FPR):}
\[\text{FPR} = \frac{\text{FP}}{\text{FP} + \text{TN}} \times 100\%\]

\textbf{False Negative Rate (FNR) = Missed Fraud Rate:}
\[\text{FNR} = \frac{\text{FN}}{\text{FN} + \text{TP}} \times 100\%\]

\textbf{Precision and Recall:}
\[\text{Precision} = \frac{\text{TP}}{\text{TP} + \text{FP}}, \qquad \text{Recall} = \frac{\text{TP}}{\text{TP} + \text{FN}}\]

\textbf{F1 Score:}
\[F_1 = \frac{2 \cdot \text{Precision} \cdot \text{Recall}}{\text{Precision} + \text{Recall}} = \frac{2\,\text{TP}}{2\,\text{TP} + \text{FP} + \text{FN}}\]

\textbf{Fraud Accounting Identity:}
\[\text{Total Fraud} = \text{Caught Fraud (TP)} + \text{Missed Fraud (FN)}\]

\textbf{Delta Metrics (change from base model):}
\[\Delta\text{FP} = \text{FP}_{\text{variant}} - \text{FP}_{\text{base}}, \qquad \Delta\text{FN} = \text{FN}_{\text{variant}} - \text{FN}_{\text{base}}\]
\[\Delta\text{Acc} = \text{Acc}_{\text{variant}} - \text{Acc}_{\text{base}}\]

\paragraph{6.12.2 Complete Variant
Catalog}\label{complete-variant-catalog}

Each variant computes a score \(s(x) \in [0,1]\) from the component
vector \(\mathbf{c} = [C, G, A, T]\) and optionally the base model’s
probability \(p_{\text{base}} = \hat{p}(x)\). The threshold for binary
classification is 0.5.

\textbf{Category A: Distance-Based (9 variants)}

{\def\LTcaptype{none} % do not increment counter
\begin{longtable}[]{@{}
  >{\raggedright\arraybackslash}p{(\linewidth - 4\tabcolsep) * \real{0.2105}}
  >{\raggedright\arraybackslash}p{(\linewidth - 4\tabcolsep) * \real{0.3158}}
  >{\raggedright\arraybackslash}p{(\linewidth - 4\tabcolsep) * \real{0.4737}}@{}}
\toprule\noalign{}
\begin{minipage}[b]{\linewidth}\raggedright
ID
\end{minipage} & \begin{minipage}[b]{\linewidth}\raggedright
Name
\end{minipage} & \begin{minipage}[b]{\linewidth}\raggedright
Formula
\end{minipage} \\
\midrule\noalign{}
\endhead
\bottomrule\noalign{}
\endlastfoot
A1 & RobustMahal &
\(s = (\mathbf{c} - \hat{\boldsymbol{\mu}})^\top \hat{\Sigma}_{\text{MCD}}^{-1} (\mathbf{c} - \hat{\boldsymbol{\mu}})\)
via Minimum Covariance Determinant \\
A2 & MahalPow18 & \(s = d_{\text{Mahal}}^{1.8}\) — power amplification
of Mahalanobis distance \\
A3 & MahalPow25 & \(s = d_{\text{Mahal}}^{2.5}\) — stronger power
amplification \\
A4 & MahalExp & \(s = \exp(d_{\text{Mahal}} / 2)\) — exponential
amplification \\
A5 & MahalChi2Tail & \(s = 1 - F_{\chi^2_4}(d_{\text{Mahal}}^2)\) —
chi-squared survival function \\
A6 & MahalTanhCorr &
\(s = \text{clip}\big(p_{\text{base}} + 0.8 \cdot \tanh(\sqrt{d}/P_{95}) \cdot (1-p_{\text{base}}) \cdot \mathbb{1}[p_{\text{base}}<0.5],\; 0, 1\big)\) \\
A7 & MahalUnsup & Same as A1 but MCD fit on all data (no label
separation) \\
A8 & kNNDist &
\(s = \frac{1}{k}\sum_{j=1}^{k}\|\mathbf{c} - \mathbf{c}_{(j)}^{\text{legit}}\|_2\),
\(k=10\) \\
A9 & kNNCorr & \(k\)-NN distance → tanh correction on
\(p_{\text{base}}\):
\(s = \text{clip}\big(p_{\text{base}} + 0.6 \cdot \tanh(d_{kNN}/P_{95}) \cdot (1-p_{\text{base}}),\; 0, 1\big)\) \\
\end{longtable}
}

\textbf{Category B: Density-Based (9 variants)}

{\def\LTcaptype{none} % do not increment counter
\begin{longtable}[]{@{}
  >{\raggedright\arraybackslash}p{(\linewidth - 4\tabcolsep) * \real{0.2105}}
  >{\raggedright\arraybackslash}p{(\linewidth - 4\tabcolsep) * \real{0.3158}}
  >{\raggedright\arraybackslash}p{(\linewidth - 4\tabcolsep) * \real{0.4737}}@{}}
\toprule\noalign{}
\begin{minipage}[b]{\linewidth}\raggedright
ID
\end{minipage} & \begin{minipage}[b]{\linewidth}\raggedright
Name
\end{minipage} & \begin{minipage}[b]{\linewidth}\raggedright
Formula
\end{minipage} \\
\midrule\noalign{}
\endhead
\bottomrule\noalign{}
\endlastfoot
B1 & KDE\_NegLog & \(s = -\log \hat{f}_{\text{KDE}}(\mathbf{c})\) —
Gaussian KDE on legitimate \\
B2 & KDE\_Sig &
\(s = \text{clip}\big(p_{\text{base}} + 0.6 \cdot \sigma((s_{\text{KDE}}-\bar{s})/\hat{\sigma}) \cdot (1-p_{\text{base}}) \cdot \mathbb{1}[p<0.5],\; 0, 1\big)\) \\
B3 & KDE\_Power & \(s = (s_{\text{KDE}} - \min + \epsilon)^{1.5}\) —
power-amplified KDE \\
B4 & GMM5\_NegLog & \(s = -\log p_{\text{GMM}_5}(\mathbf{c})\) —
5-component GMM \\
B5 & GMM3\_NegLog & \(s = -\log p_{\text{GMM}_3}(\mathbf{c})\) —
3-component GMM \\
B6 & GMM\_Tanh &
\(s = \text{clip}\big(p_{\text{base}} + 0.7 \cdot \tanh((s_{\text{GMM}}-\text{med})/\text{IQR}) \cdot (1-p_{\text{base}}),\; 0, 1\big)\) \\
B7 & LOF20 & \(s = -\text{LOF}_{k=20}(\mathbf{c})\) — Local Outlier
Factor (novelty mode) \\
B8 & LOF5 & \(s = -\text{LOF}_{k=5}(\mathbf{c})\) — tighter locality \\
B9 & LOF\_Corr &
\(s = \text{clip}\big(p_{\text{base}} + 0.5 \cdot \tanh((\text{LOF}-1)/(P_{95}-1)) \cdot (1-p_{\text{base}}),\; 0, 1\big)\) \\
\end{longtable}
}

\textbf{Category C: Vector-Space Projections (7 variants)}

{\def\LTcaptype{none} % do not increment counter
\begin{longtable}[]{@{}
  >{\raggedright\arraybackslash}p{(\linewidth - 4\tabcolsep) * \real{0.2105}}
  >{\raggedright\arraybackslash}p{(\linewidth - 4\tabcolsep) * \real{0.3158}}
  >{\raggedright\arraybackslash}p{(\linewidth - 4\tabcolsep) * \real{0.4737}}@{}}
\toprule\noalign{}
\begin{minipage}[b]{\linewidth}\raggedright
ID
\end{minipage} & \begin{minipage}[b]{\linewidth}\raggedright
Name
\end{minipage} & \begin{minipage}[b]{\linewidth}\raggedright
Formula
\end{minipage} \\
\midrule\noalign{}
\endhead
\bottomrule\noalign{}
\endlastfoot
C1 & PCA\_Mahal & PCA on legitimate → Mahalanobis in PC space \\
C2 & PCA\_Pow & \(s = s_{\text{PCA-Mahal}}^{2.0}\) — power
amplification \\
C3 & PCA\_ReconErr &
\(s = \|\mathbf{c} - \text{PCA}_2^{-1}(\text{PCA}_2(\mathbf{c}))\|_2\) —
reconstruction error \\
C4 & LDA\_1D & \(s = \mathbf{w}_{\text{LDA}}^\top \mathbf{c}\) — Fisher
discriminant projection \\
C5 & LDA\_Proba & \(s = P_{\text{LDA}}(y=1 \mid \mathbf{c})\) — LDA
posterior \\
C6 & Whitened &
\(s = \|(\mathbf{c}-\boldsymbol{\mu})/\boldsymbol{\sigma}\|_2\) —
standardised \(\ell_2\) norm \\
C7 & WhitenPow18 & \(s = s_{\text{whitened}}^{1.8}\) \\
\end{longtable}
}

\textbf{Category D: Isolation-Based (3 variants)}

{\def\LTcaptype{none} % do not increment counter
\begin{longtable}[]{@{}
  >{\raggedright\arraybackslash}p{(\linewidth - 4\tabcolsep) * \real{0.2105}}
  >{\raggedright\arraybackslash}p{(\linewidth - 4\tabcolsep) * \real{0.3158}}
  >{\raggedright\arraybackslash}p{(\linewidth - 4\tabcolsep) * \real{0.4737}}@{}}
\toprule\noalign{}
\begin{minipage}[b]{\linewidth}\raggedright
ID
\end{minipage} & \begin{minipage}[b]{\linewidth}\raggedright
Name
\end{minipage} & \begin{minipage}[b]{\linewidth}\raggedright
Formula
\end{minipage} \\
\midrule\noalign{}
\endhead
\bottomrule\noalign{}
\endlastfoot
D1 & IsoForest100 & \(s = -\text{IF}_{100}(\mathbf{c})\) — 100-tree
Isolation Forest \\
D2 & IsoForest10 & \(s = -\text{IF}_{10}(\mathbf{c})\) — 10-tree light
variant \\
D3 & IsoCorr &
\(s = \text{clip}\big(p_{\text{base}} + 0.6 \cdot \tanh((s_{\text{IF}}-\text{med})/\hat{\sigma}) \cdot (1-p_{\text{base}}),\; 0, 1\big)\) \\
\end{longtable}
}

\textbf{Category E: SVM-Based (5 variants)}

{\def\LTcaptype{none} % do not increment counter
\begin{longtable}[]{@{}
  >{\raggedright\arraybackslash}p{(\linewidth - 4\tabcolsep) * \real{0.2105}}
  >{\raggedright\arraybackslash}p{(\linewidth - 4\tabcolsep) * \real{0.3158}}
  >{\raggedright\arraybackslash}p{(\linewidth - 4\tabcolsep) * \real{0.4737}}@{}}
\toprule\noalign{}
\begin{minipage}[b]{\linewidth}\raggedright
ID
\end{minipage} & \begin{minipage}[b]{\linewidth}\raggedright
Name
\end{minipage} & \begin{minipage}[b]{\linewidth}\raggedright
Formula
\end{minipage} \\
\midrule\noalign{}
\endhead
\bottomrule\noalign{}
\endlastfoot
E1 & OCSVM\_RBF & \(s = -f_{\text{SVM}}^{\text{rbf}}(\mathbf{c})\) —
One-Class SVM (RBF, \(\nu=0.05\)) \\
E2 & OCSVM\_Poly & \(s = -f_{\text{SVM}}^{\text{poly}}(\mathbf{c})\) —
polynomial kernel \\
E3 & OCSVM\_Lin & \(s = -f_{\text{SVM}}^{\text{lin}}(\mathbf{c})\) —
linear kernel (\(\nu=0.1\)) \\
E4 & SVM\_Corr &
\(s = \text{clip}\big(p_{\text{base}} + 0.7 \cdot \tanh((s_{\text{SVM}}-\bar{s})/\hat{\sigma}) \cdot (1-p_{\text{base}}) \cdot \mathbb{1}[p<0.6],\; 0, 1\big)\) \\
E5 & SVM\_Blend &
\(s = \text{clip}(0.5\, p_{\text{base}} + 0.5 \cdot \sigma(0.5(s_{\text{SVM}}-\text{med})/\hat{\sigma}),\; 0, 1)\) \\
\end{longtable}
}

\textbf{Category F: Quadratic / Discriminant (5 variants)}

{\def\LTcaptype{none} % do not increment counter
\begin{longtable}[]{@{}
  >{\raggedright\arraybackslash}p{(\linewidth - 4\tabcolsep) * \real{0.2105}}
  >{\raggedright\arraybackslash}p{(\linewidth - 4\tabcolsep) * \real{0.3158}}
  >{\raggedright\arraybackslash}p{(\linewidth - 4\tabcolsep) * \real{0.4737}}@{}}
\toprule\noalign{}
\begin{minipage}[b]{\linewidth}\raggedright
ID
\end{minipage} & \begin{minipage}[b]{\linewidth}\raggedright
Name
\end{minipage} & \begin{minipage}[b]{\linewidth}\raggedright
Formula
\end{minipage} \\
\midrule\noalign{}
\endhead
\bottomrule\noalign{}
\endlastfoot
F1 & QDA & \(s = P_{\text{QDA}}(y=1 \mid \mathbf{c})\) — Quadratic
Discriminant Analysis (\(\lambda=0.1\)) \\
F2 & QuadSurf &
\(s = \text{clip}(p_{\text{base}} + \phi_2(\mathbf{c})^\top\hat{\boldsymbol{\beta}},\; 0, 1)\)
where \(\phi_2\) = degree-2 polynomial, \(\hat{\boldsymbol{\beta}}\) via
weighted ridge targeting \((y - p_{\text{base}})\) \\
F3 & DiscrimPM & Quadratic discriminant \(\pm\sqrt{\Delta}\):
\(\Delta = b^2 - 4ac\) where
\(a=\|\mathbf{c}\|^2, b=-\text{MFLS}, c=p_{\text{base}}-0.5\) \\
F4 & QuadInterRidge & Ridge regression on
\([C^2, G^2, A^2, T^2, CG, CA, CT, GA, GT, AT, C, G, A, T]\) \\
F5 & QuadGrid & Grid-search over
\(s = \text{clip}(p_{\text{base}} + (a\,m^2 + b\,m)(1-p_{\text{base}}),\; 0, 1)\),
\(7\times7\) grid \\
\end{longtable}
}

\textbf{Category G: Distribution Transforms (7 variants)}

{\def\LTcaptype{none} % do not increment counter
\begin{longtable}[]{@{}
  >{\raggedright\arraybackslash}p{(\linewidth - 4\tabcolsep) * \real{0.2105}}
  >{\raggedright\arraybackslash}p{(\linewidth - 4\tabcolsep) * \real{0.3158}}
  >{\raggedright\arraybackslash}p{(\linewidth - 4\tabcolsep) * \real{0.4737}}@{}}
\toprule\noalign{}
\begin{minipage}[b]{\linewidth}\raggedright
ID
\end{minipage} & \begin{minipage}[b]{\linewidth}\raggedright
Name
\end{minipage} & \begin{minipage}[b]{\linewidth}\raggedright
Formula
\end{minipage} \\
\midrule\noalign{}
\endhead
\bottomrule\noalign{}
\endlastfoot
G1 & BoxCox & \(s = \text{BoxCox}(\text{MFLS} - \min + \epsilon)\) \\
G2 & YeoJohnson & \(s = \text{YeoJohnson}(\text{MFLS})\) \\
G3 & Quantile\_U & \(s = \frac{1}{4}\sum_j Q_U^{-1}(c_j)\) —
quantile-uniform mean \\
G4 & Quantile\_G & Quantile-Gaussian → Mahalanobis in
\(\mathcal{N}(0,1)\) space \\
G5 & RankSigmoid &
\(s = \sigma(6(\text{rank}(\text{MFLS})/n - 0.5))\) \\
G6 & LogitPow & \(s = \sigma(-\text{sgn}(l)|l|^2/5)\),
\(l = \text{logit}(\text{MFLS})\) \\
G7 & YeoJCorr &
\(s = \text{clip}(p_{\text{base}} + 0.5 \cdot \tanh((\text{YJ}-\bar{\text{YJ}})/\hat{\sigma}) \cdot (1-p_{\text{base}}),\; 0, 1)\) \\
\end{longtable}
}

\textbf{Category H: Interaction \& Enrichment (6 variants)}

{\def\LTcaptype{none} % do not increment counter
\begin{longtable}[]{@{}
  >{\raggedright\arraybackslash}p{(\linewidth - 4\tabcolsep) * \real{0.2105}}
  >{\raggedright\arraybackslash}p{(\linewidth - 4\tabcolsep) * \real{0.3158}}
  >{\raggedright\arraybackslash}p{(\linewidth - 4\tabcolsep) * \real{0.4737}}@{}}
\toprule\noalign{}
\begin{minipage}[b]{\linewidth}\raggedright
ID
\end{minipage} & \begin{minipage}[b]{\linewidth}\raggedright
Name
\end{minipage} & \begin{minipage}[b]{\linewidth}\raggedright
Formula
\end{minipage} \\
\midrule\noalign{}
\endhead
\bottomrule\noalign{}
\endlastfoot
H1 & Poly2LR & \(s = P_{\text{LR}}(y=1 \mid \phi_2(\mathbf{c}))\) —
degree-2 polynomial features \\
H2 & Ratios &
\(s = P_{\text{LR}}(y=1 \mid [C/(G+\epsilon),\; A/(T+\epsilon),\; CA/(GT+\epsilon)])\) \\
H3 & HarmonicMean & \(s = 4/(C^{-1}+G^{-1}+A^{-1}+T^{-1})\) \\
H4 & GeometricMean & \(s = (CGAT)^{1/4}\) \\
H5 & LpNorms &
\(s = P_{\text{LR}}(y=1 \mid [L_{0.5}, L_1, L_2, L_\infty])\) \\
H6 & Entropy & \(s = -\sum_j p_j \ln p_j\),
\(p_j = c_j / \sum_i c_i\) \\
\end{longtable}
}

\textbf{Category I: Hybrid Gates (7 variants)}

{\def\LTcaptype{none} % do not increment counter
\begin{longtable}[]{@{}
  >{\raggedright\arraybackslash}p{(\linewidth - 4\tabcolsep) * \real{0.2105}}
  >{\raggedright\arraybackslash}p{(\linewidth - 4\tabcolsep) * \real{0.3158}}
  >{\raggedright\arraybackslash}p{(\linewidth - 4\tabcolsep) * \real{0.4737}}@{}}
\toprule\noalign{}
\begin{minipage}[b]{\linewidth}\raggedright
ID
\end{minipage} & \begin{minipage}[b]{\linewidth}\raggedright
Name
\end{minipage} & \begin{minipage}[b]{\linewidth}\raggedright
Formula
\end{minipage} \\
\midrule\noalign{}
\endhead
\bottomrule\noalign{}
\endlastfoot
I1 & Mahal\_AND & Mahalanobis × \((C \cdot G)\) AND gate, tanh
correction \\
I2 & KDE\_XOR & KDE × \(\tanh(CA - GT)\) XOR gate \\
I3 & Iso\_Ratio & IsoForest × \(\tanh(2(r-1))\) where
\(r = (C+G)/(A+T+\epsilon)\) \\
I4 & SVM\_Tanh & SVM ×
\(\tanh(\text{sat}) \cdot \sigma(-12(p_{\text{base}}-0.5))\) \\
I5 & LDA\_Power & LDA normalised to \([0,1]\), then \(s^{1.5}\) \\
I6 & EnsVote & Majority vote across A1, B1, D1, E1 at \(P_{90}\)
thresholds \\
I7 & StackedLR &
\(s = P_{\text{LR}}(y=1 \mid [s_{\text{Mahal}}, s_{\text{KDE}}, s_{\text{IF}}, s_{\text{SVM}}])\)
— stacked anomaly scores \\
\end{longtable}
}

\textbf{Category J: LOF + Power-Law Chain (8 variants)}

This category introduces the \textbf{MFLS-LP (LOF-Power)} family,
testing the hypothesis that amplifying Local Outlier Factor shifts via
power-law transforms can expose density-based blind spots invisible to
the linear MFLS.

{\def\LTcaptype{none} % do not increment counter
\begin{longtable}[]{@{}
  >{\raggedright\arraybackslash}p{(\linewidth - 4\tabcolsep) * \real{0.2105}}
  >{\raggedright\arraybackslash}p{(\linewidth - 4\tabcolsep) * \real{0.3158}}
  >{\raggedright\arraybackslash}p{(\linewidth - 4\tabcolsep) * \real{0.4737}}@{}}
\toprule\noalign{}
\begin{minipage}[b]{\linewidth}\raggedright
ID
\end{minipage} & \begin{minipage}[b]{\linewidth}\raggedright
Name
\end{minipage} & \begin{minipage}[b]{\linewidth}\raggedright
Formula
\end{minipage} \\
\midrule\noalign{}
\endhead
\bottomrule\noalign{}
\endlastfoot
J1 & LOF\_Pow15 &
\(s = (\max(\text{LOF}_{20}-1, 0) + \epsilon)^{1.5}\) \\
J2 & LOF\_Pow18 & \(s = (\max(\text{LOF}_{20}-1, 0) + \epsilon)^{1.8}\)
— primary MFLS-LP \\
J3 & LOF\_Pow20 & \(s = (\max(\text{LOF}_{20}-1, 0) + \epsilon)^{2.0}\)
— aggressive \\
J4 & LOF15\_Pow18 &
\(s = (\max(\text{LOF}_{15}-1, 0) + \epsilon)^{1.8}\) — tighter
locality \\
J5 & LOF30\_Pow18 &
\(s = (\max(\text{LOF}_{30}-1, 0) + \epsilon)^{1.8}\) — smoother
density \\
J6 & LOF\_Pow\_Corr &
\(\text{MFLS}' = \text{MFLS} \cdot (\max(\text{LOF}_{20}/\tau, 1))^{1.8}\),
\(\tau=1.5\); then
\(s = \text{clip}(p_{\text{base}} + 0.8 \cdot \tanh(\text{MFLS}') \cdot (1-p_{\text{base}}),\; 0, 1)\) \\
J7 & LOF\_AdaptPow & Adaptive
\(\gamma = \text{clip}(|\lambda_{\text{YJ}}|+1, 1.2, 3.0)\) via
Yeo-Johnson on legitimate LOF \\
J8 & kNN\_Pow18 &
\(s = (d_{kNN}/P_{95}(d_{kNN}^{\text{legit}}))^{1.8}\) \\
\end{longtable}
}

\paragraph{6.12.3 Aggregate Leaderboard — Top-10
Variants}\label{aggregate-leaderboard-top-10-variants}

\textbf{Table 16: Top-10 variants by mean F1 across 12 dataset × model
combinations}

{\def\LTcaptype{none} % do not increment counter
\begin{longtable}[]{@{}
  >{\raggedright\arraybackslash}p{(\linewidth - 10\tabcolsep) * \real{0.0896}}
  >{\raggedright\arraybackslash}p{(\linewidth - 10\tabcolsep) * \real{0.1343}}
  >{\raggedright\arraybackslash}p{(\linewidth - 10\tabcolsep) * \real{0.1493}}
  >{\raggedright\arraybackslash}p{(\linewidth - 10\tabcolsep) * \real{0.1343}}
  >{\raggedright\arraybackslash}p{(\linewidth - 10\tabcolsep) * \real{0.2836}}
  >{\raggedright\arraybackslash}p{(\linewidth - 10\tabcolsep) * \real{0.2090}}@{}}
\toprule\noalign{}
\begin{minipage}[b]{\linewidth}\raggedright
Rank
\end{minipage} & \begin{minipage}[b]{\linewidth}\raggedright
Variant
\end{minipage} & \begin{minipage}[b]{\linewidth}\raggedright
Category
\end{minipage} & \begin{minipage}[b]{\linewidth}\raggedright
Mean F1
\end{minipage} & \begin{minipage}[b]{\linewidth}\raggedright
\(\Delta\)F1 vs Base
\end{minipage} & \begin{minipage}[b]{\linewidth}\raggedright
F1 Wins (/12)
\end{minipage} \\
\midrule\noalign{}
\endhead
\bottomrule\noalign{}
\endlastfoot
1 & \textbf{F2\_QuadSurf} & F: Quadratic & \textbf{0.7870} & +0.2366 &
8 \\
2 & B2\_KDE\_Sig & B: Density & 0.7414 & +0.1911 & 5 \\
3 & A9\_kNNCorr & A: Distance & 0.7294 & +0.1791 & 5 \\
4 & I7\_StackedLR & I: Hybrid & 0.7168 & +0.1665 & 7 \\
5 & B6\_GMM\_Tanh & B: Density & 0.7130 & +0.1627 & 4 \\
6 & I6\_EnsVote & I: Hybrid & 0.6996 & +0.1493 & 4 \\
7 & E5\_SVM\_Blend & E: SVM & 0.6989 & +0.1486 & 4 \\
8 & E4\_SVM\_Corr & E: SVM & 0.6967 & +0.1464 & 5 \\
9 & D3\_IsoCorr & D: Isolation & 0.6881 & +0.1378 & 4 \\
10 & B9\_LOF\_Corr & B: Density & 0.6809 & +0.1305 & 5 \\
\end{longtable}
}

\textbf{Key finding:} 26 of 66 variants beat the base model’s mean F1 of
0.5503. The best (F2\_QuadSurf) improves F1 by +0.2366 — a \textbf{43\%
relative improvement}. 9 of 10 categories contain at least one variant
that beats the baseline, confirming that multiple algorithmic families
can exploit the BSDT component structure.

\textbf{Table 17: Per-category best variant}

{\def\LTcaptype{none} % do not increment counter
\begin{longtable}[]{@{}llll@{}}
\toprule\noalign{}
Category & Best Variant & Mean F1 & \(\Delta\)F1 \\
\midrule\noalign{}
\endhead
\bottomrule\noalign{}
\endlastfoot
A: Distance & A9\_kNNCorr & 0.7294 & +0.179 \\
B: Density & B2\_KDE\_Sig & 0.7414 & +0.191 \\
C: Projections & C5\_LDA\_Proba & 0.5885 & +0.038 \\
D: Isolation & D3\_IsoCorr & 0.6881 & +0.138 \\
E: SVM & E5\_SVM\_Blend & 0.6989 & +0.149 \\
F: Quadratic & \textbf{F2\_QuadSurf} & \textbf{0.7870} &
\textbf{+0.237} \\
G: Transforms & G7\_YeoJCorr & 0.6252 & +0.075 \\
H: Interaction & H1\_Poly2LR & 0.6048 & +0.054 \\
I: Hybrid & I7\_StackedLR & 0.7168 & +0.167 \\
J: LOF+Power & J6\_LOF\_Pow\_Corr & 0.6470 & +0.097 \\
\end{longtable}
}

\textbf{In-domain vs Out-of-domain:} The largest gains occur on
out-of-domain datasets where the base model has zero domain-specific
knowledge:

{\def\LTcaptype{none} % do not increment counter
\begin{longtable}[]{@{}
  >{\raggedright\arraybackslash}p{(\linewidth - 6\tabcolsep) * \real{0.1951}}
  >{\raggedright\arraybackslash}p{(\linewidth - 6\tabcolsep) * \real{0.2195}}
  >{\raggedright\arraybackslash}p{(\linewidth - 6\tabcolsep) * \real{0.3171}}
  >{\raggedright\arraybackslash}p{(\linewidth - 6\tabcolsep) * \real{0.2683}}@{}}
\toprule\noalign{}
\begin{minipage}[b]{\linewidth}\raggedright
Domain
\end{minipage} & \begin{minipage}[b]{\linewidth}\raggedright
Base F1
\end{minipage} & \begin{minipage}[b]{\linewidth}\raggedright
Best F1 (F2)
\end{minipage} & \begin{minipage}[b]{\linewidth}\raggedright
\(\Delta\)F1
\end{minipage} \\
\midrule\noalign{}
\endhead
\bottomrule\noalign{}
\endlastfoot
In-domain (Elliptic, XBlock) & 0.676 & 0.811 & +0.135 \\
Out-of-domain (CreditCard, Shuttle, Mammography, Pendigits) & 0.425 &
0.763 & +0.339 \\
\end{longtable}
}

\paragraph{6.12.4 Fraud Accounting — Confusion Matrix
Analysis}\label{fraud-accounting-confusion-matrix-analysis}

To evaluate whether improved F1 comes at the cost of increased false
alerts, we compute full confusion matrix accounting for the top-8
variants plus base model, aggregated across all 12 dataset × model
combinations.

\textbf{Table 18: Aggregate fraud accounting (summed across 12 combos,
\(N_{\text{total}} = 77{,}790\) fraud instances)}

{\def\LTcaptype{none} % do not increment counter
\begin{longtable}[]{@{}
  >{\raggedright\arraybackslash}p{(\linewidth - 20\tabcolsep) * \real{0.0687}}
  >{\raggedright\arraybackslash}p{(\linewidth - 20\tabcolsep) * \real{0.0916}}
  >{\raggedright\arraybackslash}p{(\linewidth - 20\tabcolsep) * \real{0.0992}}
  >{\raggedright\arraybackslash}p{(\linewidth - 20\tabcolsep) * \real{0.0992}}
  >{\raggedright\arraybackslash}p{(\linewidth - 20\tabcolsep) * \real{0.1450}}
  >{\raggedright\arraybackslash}p{(\linewidth - 20\tabcolsep) * \real{0.0840}}
  >{\raggedright\arraybackslash}p{(\linewidth - 20\tabcolsep) * \real{0.0916}}
  >{\raggedright\arraybackslash}p{(\linewidth - 20\tabcolsep) * \real{0.0840}}
  >{\raggedright\arraybackslash}p{(\linewidth - 20\tabcolsep) * \real{0.0916}}
  >{\raggedright\arraybackslash}p{(\linewidth - 20\tabcolsep) * \real{0.0687}}
  >{\raggedright\arraybackslash}p{(\linewidth - 20\tabcolsep) * \real{0.0763}}@{}}
\toprule\noalign{}
\begin{minipage}[b]{\linewidth}\raggedright
Variant
\end{minipage} & \begin{minipage}[b]{\linewidth}\raggedright
Total Fraud
\end{minipage} & \begin{minipage}[b]{\linewidth}\raggedright
Caught (TP)
\end{minipage} & \begin{minipage}[b]{\linewidth}\raggedright
Missed (FN)
\end{minipage} & \begin{minipage}[b]{\linewidth}\raggedright
False Alerts (FP)
\end{minipage} & \begin{minipage}[b]{\linewidth}\raggedright
\(\Delta\)FP
\end{minipage} & \begin{minipage}[b]{\linewidth}\raggedright
\(\Delta\)FP\%
\end{minipage} & \begin{minipage}[b]{\linewidth}\raggedright
\(\Delta\)FN
\end{minipage} & \begin{minipage}[b]{\linewidth}\raggedright
\(\Delta\)FN\%
\end{minipage} & \begin{minipage}[b]{\linewidth}\raggedright
Mean F1
\end{minipage} & \begin{minipage}[b]{\linewidth}\raggedright
Mean Acc
\end{minipage} \\
\midrule\noalign{}
\endhead
\bottomrule\noalign{}
\endlastfoot
\textbf{BASE MODEL} & 77,790 & 76,046 & 1,744 & 44,740 & — & — & — & — &
0.5503 & 85.88\% \\
\textbf{F2\_QuadSurf} & 77,790 & 75,513 & 2,277 & 5,345 &
\textbf{−39,395} & \textbf{−88.1\%} & +533 & +30.6\% & \textbf{0.7870} &
\textbf{96.34\%} \\
B2\_KDE\_Sig & 77,790 & 76,469 & 1,321 & 17,368 & −27,372 & −61.2\% &
\textbf{−423} & \textbf{−24.3\%} & 0.7414 & 85.52\% \\
A9\_kNNCorr & 77,790 & 76,323 & 1,467 & 12,878 & −31,862 & −71.2\% &
−277 & −15.9\% & 0.7294 & 93.22\% \\
I7\_StackedLR & 77,790 & 75,192 & 2,598 & 7,876 & −36,864 & −82.4\% &
+854 & +49.0\% & 0.7168 & 93.76\% \\
B6\_GMM\_Tanh & 77,790 & 75,170 & 2,620 & 5,557 & −39,183 & −87.6\% &
+876 & +50.2\% & 0.7130 & 96.33\% \\
D3\_IsoCorr & 77,790 & 74,993 & 2,797 & 16,943 & −27,797 & −62.1\% &
+1,053 & +60.4\% & 0.6881 & 95.64\% \\
J6\_LOF\_Pow\_Corr & 77,790 & 76,358 & 1,432 & 20,191 & −24,549 &
−54.9\% & −312 & −17.9\% & 0.6470 & 83.32\% \\
F5\_QuadGrid & 77,790 & 74,292 & 3,498 & 8,820 & −35,920 & −80.3\% &
+1,754 & +100.6\% & 0.6360 & 92.22\% \\
\end{longtable}
}

\textbf{Critical insight — False Positive Reduction:} Every single top
variant \textbf{massively reduces false positives} compared to the base
model. F2\_QuadSurf eliminates 88.1\% of false alerts (from 44,740 down
to 5,345). This is because the base model on out-of-domain datasets
often predicts everything as positive (e.g., Shuttle AMTTP-XGB93:
FP=36,469), and the variants learn more discriminative boundaries.

\textbf{FN trade-off:} Three variants (B2\_KDE\_Sig, A9\_kNNCorr,
J6\_LOF\_Pow\_Corr) reduce \textbf{both} FP and FN simultaneously — a
rare and desirable property. The remaining variants trade a modest FN
increase for massive FP reduction, resulting in dramatically higher F1
and accuracy.

\paragraph{6.12.5 Per-Dataset Fraud
Accounting}\label{per-dataset-fraud-accounting}

\textbf{Table 19: Selected per-dataset fraud accounting for F2\_QuadSurf
(best overall) and J6\_LOF\_Pow\_Corr (Category J champion)}

{\def\LTcaptype{none} % do not increment counter
\begin{longtable}[]{@{}
  >{\raggedright\arraybackslash}p{(\linewidth - 18\tabcolsep) * \real{0.1216}}
  >{\raggedright\arraybackslash}p{(\linewidth - 18\tabcolsep) * \real{0.0946}}
  >{\raggedright\arraybackslash}p{(\linewidth - 18\tabcolsep) * \real{0.1622}}
  >{\raggedright\arraybackslash}p{(\linewidth - 18\tabcolsep) * \real{0.1081}}
  >{\raggedright\arraybackslash}p{(\linewidth - 18\tabcolsep) * \real{0.1081}}
  >{\raggedright\arraybackslash}p{(\linewidth - 18\tabcolsep) * \real{0.1081}}
  >{\raggedright\arraybackslash}p{(\linewidth - 18\tabcolsep) * \real{0.0676}}
  >{\raggedright\arraybackslash}p{(\linewidth - 18\tabcolsep) * \real{0.0811}}
  >{\raggedright\arraybackslash}p{(\linewidth - 18\tabcolsep) * \real{0.0811}}
  >{\raggedright\arraybackslash}p{(\linewidth - 18\tabcolsep) * \real{0.0676}}@{}}
\toprule\noalign{}
\begin{minipage}[b]{\linewidth}\raggedright
Dataset
\end{minipage} & \begin{minipage}[b]{\linewidth}\raggedright
Model
\end{minipage} & \begin{minipage}[b]{\linewidth}\raggedright
Total Fraud
\end{minipage} & \begin{minipage}[b]{\linewidth}\raggedright
Method
\end{minipage} & \begin{minipage}[b]{\linewidth}\raggedright
Caught
\end{minipage} & \begin{minipage}[b]{\linewidth}\raggedright
Missed
\end{minipage} & \begin{minipage}[b]{\linewidth}\raggedright
FP
\end{minipage} & \begin{minipage}[b]{\linewidth}\raggedright
Acc\%
\end{minipage} & \begin{minipage}[b]{\linewidth}\raggedright
MCR\%
\end{minipage} & \begin{minipage}[b]{\linewidth}\raggedright
F1
\end{minipage} \\
\midrule\noalign{}
\endhead
\bottomrule\noalign{}
\endlastfoot
\textbf{Shuttle} & \textbf{AMTTP-XGB93} & \textbf{2,809} & Base & 2,809
& 0 & 36,469 & 7.15 & 92.85 & 0.134 \\
& & & F2\_QuadSurf & 2,653 & 156 & 22 & \textbf{99.55} & 0.45 &
\textbf{0.968} \\
& & & J6\_LOF\_Pow\_Corr & 2,655 & 154 & 198 & 99.10 & 0.90 & 0.938 \\
\textbf{Pendigits} & \textbf{AMTTP-XGB93} & \textbf{125} & Base & 0 &
125 & 0 & 97.73 & 2.27 & 0.000 \\
& & & F2\_QuadSurf & 117 & 8 & 9 & \textbf{99.69} & 0.31 &
\textbf{0.932} \\
& & & J6\_LOF\_Pow\_Corr & 11 & 114 & 12 & 97.71 & 2.29 & 0.149 \\
\textbf{Credit Card} & \textbf{Fresh-XGB} & \textbf{394} & Base & 284 &
110 & 32 & 99.94 & 0.06 & 0.800 \\
& & & F2\_QuadSurf & 284 & 110 & 27 & 99.94 & 0.06 & \textbf{0.806} \\
& & & J6\_LOF\_Pow\_Corr & 282 & 112 & 22 & 99.94 & 0.06 & 0.808 \\
\textbf{Elliptic} & \textbf{Fresh-XGB} & \textbf{33,616} & Base & 33,426
& 190 & 532 & 98.06 & 1.94 & 0.989 \\
& & & F2\_QuadSurf & 33,419 & 197 & 521 & 98.07 & 1.93 & 0.989 \\
& & & J6\_LOF\_Pow\_Corr & 33,394 & 222 & 482 & \textbf{98.11} & 1.89 &
\textbf{0.990} \\
\end{longtable}
}

\textbf{Shuttle AMTTP-XGB93 — the flagship rescue:} The base model
(pre-trained on blockchain data) applied to aerospace telemetry flags
nearly every sample as fraud (FP=36,469 out of 36,469 legitimate),
yielding F1=0.134. F2\_QuadSurf reduces false alerts from 36,469 to just
22 (a \textbf{99.94\% FP reduction}) while catching 2,653 of 2,809 fraud
cases, lifting F1 to 0.968 — a gain of +0.834 in absolute F1.

\textbf{J6\_LOF\_Pow\_Corr on Shuttle AMTTP-XGB93} also achieves
remarkable recovery: FP drops from 36,469 to 198, catching 2,655 fraud
cases, F1 = 0.938. This confirms that the LOF+Power blind-spot
correction works specifically where the base model completely fails on
unseen domains.

\paragraph{6.12.6 FP vs FN Trade-off
Analysis}\label{fp-vs-fn-trade-off-analysis}

\textbf{Table 20: FP/FN change direction across 12 combos per variant}

{\def\LTcaptype{none} % do not increment counter
\begin{longtable}[]{@{}
  >{\raggedright\arraybackslash}p{(\linewidth - 12\tabcolsep) * \real{0.0826}}
  >{\raggedright\arraybackslash}p{(\linewidth - 12\tabcolsep) * \real{0.2018}}
  >{\raggedright\arraybackslash}p{(\linewidth - 12\tabcolsep) * \real{0.2018}}
  >{\raggedright\arraybackslash}p{(\linewidth - 12\tabcolsep) * \real{0.1193}}
  >{\raggedright\arraybackslash}p{(\linewidth - 12\tabcolsep) * \real{0.1193}}
  >{\raggedright\arraybackslash}p{(\linewidth - 12\tabcolsep) * \real{0.1193}}
  >{\raggedright\arraybackslash}p{(\linewidth - 12\tabcolsep) * \real{0.1560}}@{}}
\toprule\noalign{}
\begin{minipage}[b]{\linewidth}\raggedright
Variant
\end{minipage} & \begin{minipage}[b]{\linewidth}\raggedright
FP increased (combos)
\end{minipage} & \begin{minipage}[b]{\linewidth}\raggedright
FP decreased (combos)
\end{minipage} & \begin{minipage}[b]{\linewidth}\raggedright
FN increased
\end{minipage} & \begin{minipage}[b]{\linewidth}\raggedright
FN decreased
\end{minipage} & \begin{minipage}[b]{\linewidth}\raggedright
Both reduced
\end{minipage} & \begin{minipage}[b]{\linewidth}\raggedright
Dominant pattern
\end{minipage} \\
\midrule\noalign{}
\endhead
\bottomrule\noalign{}
\endlastfoot
F2\_QuadSurf & 7 & 5 & 5 & 6 & 1/12 & FP focus \\
B2\_KDE\_Sig & 3 & 5 & 7 & 2 & 3/12 & \textbf{FN focus} \\
A9\_kNNCorr & 3 & 8 & 7 & 4 & 0/12 & FN focus \\
I7\_StackedLR & 4 & 8 & 5 & 7 & 0/12 & FN focus \\
B6\_GMM\_Tanh & 5 & 4 & 4 & 5 & 0/12 & FP focus \\
D3\_IsoCorr & 6 & 4 & 3 & 6 & 0/12 & FP focus \\
J6\_LOF\_Pow\_Corr & 6 & 4 & 4 & 7 & 0/12 & FP focus \\
F5\_QuadGrid & 7 & 4 & 4 & 6 & 1/12 & FP focus \\
\end{longtable}
}

\textbf{Interpretation:} No variant simultaneously reduces both FP and
FN across all 12 combos — this reflects the fundamental
\textbf{precision-recall trade-off}. However, the dominant pattern is
overwhelmingly \textbf{FP reduction}: across the 8 variants × 12 combos
= 96 evaluations, FP decreases in the majority of cases, driven by the
base model’s extreme false-alert rates on out-of-domain data. The
resulting accuracy gains (base 85.88\% → F2 96.34\%) are primarily
driven by FP elimination.

\textbf{B2\_KDE\_Sig is the exception:} It uniquely focuses on FN
reduction (catching more fraud) while maintaining conservative FP
increases. It is the only top-variant to achieve \textbf{net aggregate
FN reduction} (−423, −24.3\%) — making it the preferred variant when
missed fraud is the primary compliance concern.

\paragraph{6.12.7 Best Variant for Reducing Missed Fraud (Per
Dataset)}\label{best-variant-for-reducing-missed-fraud-per-dataset}

\textbf{Table 21: Variant that minimises missed fraud for each dataset ×
model combination}

{\def\LTcaptype{none} % do not increment counter
\begin{longtable}[]{@{}
  >{\raggedright\arraybackslash}p{(\linewidth - 12\tabcolsep) * \real{0.1154}}
  >{\raggedright\arraybackslash}p{(\linewidth - 12\tabcolsep) * \real{0.0897}}
  >{\raggedright\arraybackslash}p{(\linewidth - 12\tabcolsep) * \real{0.1538}}
  >{\raggedright\arraybackslash}p{(\linewidth - 12\tabcolsep) * \real{0.1667}}
  >{\raggedright\arraybackslash}p{(\linewidth - 12\tabcolsep) * \real{0.1410}}
  >{\raggedright\arraybackslash}p{(\linewidth - 12\tabcolsep) * \real{0.1923}}
  >{\raggedright\arraybackslash}p{(\linewidth - 12\tabcolsep) * \real{0.1410}}@{}}
\toprule\noalign{}
\begin{minipage}[b]{\linewidth}\raggedright
Dataset
\end{minipage} & \begin{minipage}[b]{\linewidth}\raggedright
Model
\end{minipage} & \begin{minipage}[b]{\linewidth}\raggedright
Base Missed
\end{minipage} & \begin{minipage}[b]{\linewidth}\raggedright
Best Variant
\end{minipage} & \begin{minipage}[b]{\linewidth}\raggedright
New Missed
\end{minipage} & \begin{minipage}[b]{\linewidth}\raggedright
\(\Delta\)Missed
\end{minipage} & \begin{minipage}[b]{\linewidth}\raggedright
Reduction
\end{minipage} \\
\midrule\noalign{}
\endhead
\bottomrule\noalign{}
\endlastfoot
Pendigits & AMTTP-XGB93 & 125 & A1\_RobustMahal & 0 & −125 & 100\% \\
Pendigits & Fresh-XGB & 21 & A1\_RobustMahal & 0 & −21 & 100\% \\
Mammography & AMTTP-XGB93 & 208 & A4\_MahalExp & 0 & −208 & 100\% \\
Mammography & Fresh-XGB & 80 & A4\_MahalExp & 0 & −80 & 100\% \\
XBlock & AMTTP-XGB93 & 212 & A4\_MahalExp & 0 & −212 & 100\% \\
XBlock & Fresh-XGB & 340 & A4\_MahalExp & 0 & −340 & 100\% \\
Shuttle & AMTTP-XGB93 & 0 & (none better) & — & — & — \\
Shuttle & Fresh-XGB & 64 & A1\_RobustMahal & 0 & −64 & 100\% \\
Credit Card & AMTTP-XGB93 & 394 & A4\_MahalExp & 0 & −394 & 100\% \\
Credit Card & Fresh-XGB & 110 & A4\_MahalExp & 0 & −110 & 100\% \\
Elliptic & AMTTP-XGB93 & 0 & (none better) & — & — & — \\
Elliptic & Fresh-XGB & 190 & A4\_MahalExp & 0 & −190 & 100\% \\
\end{longtable}
}

\textbf{Remarkable finding:} In 10 of 12 combinations, \textbf{at least
one variant achieves zero missed fraud} — catching every single
fraudulent transaction. The winners are exclusively from Category A
(Distance-Based): A4\_MahalExp and A1\_RobustMahal. However, these
variants achieve 100\% recall at the cost of very high FP rates (F1
drops to 0.003–0.048), making them impractical as standalone classifiers
but invaluable as \textbf{fraud alert generators} in a two-stage
pipeline.

\paragraph{6.12.8 Summary of V2 Findings}\label{summary-of-v2-findings}

The 66-variant, 792-evaluation experiment yields five principal
conclusions:

\begin{enumerate}
\def\labelenumi{\arabic{enumi}.}
\item
  \textbf{Quadratic surface fitting (F2) dominates all alternatives.} By
  learning a degree-2 polynomial correction on the BSDT component
  vector, F2\_QuadSurf achieves Mean F1 = 0.787, eliminating 88.1\% of
  false positives while maintaining high recall. The formula:
  \[s_{\text{F2}} = \text{clip}\big(p_{\text{base}} + \phi_2([C,G,A,T])^\top\hat{\boldsymbol{\beta}},\; 0, 1\big)\]
  captures the \textbf{non-linear interaction structure} among the four
  BSDT components that linear MFLS misses.
\item
  \textbf{False positives are the primary beneficiary.} Across all top
  variants, the dominant improvement is FP reduction (−54.9\% to
  −88.1\%), not FN reduction. This occurs because the base model
  catastrophically over-alerts on out-of-domain data, and the variants
  learn to discriminate within the component space.
\item
  \textbf{Three variants achieve net FN reduction.} B2\_KDE\_Sig
  (−24.3\% FN), A9\_kNNCorr (−15.9\%), and J6\_LOF\_Pow\_Corr (−17.9\%)
  catch strictly more fraud than the base model while also reducing
  false alerts — the ideal compliance scenario.
\item
  \textbf{LOF+Power chain (Category J) validates the blind-spot
  hypothesis.} J6\_LOF\_Pow\_Corr (rank \#14, F1=0.647) demonstrates
  that density-based blind spots invisible to the linear MFLS can be
  exposed through LOF amplification. Its most dramatic result — Shuttle
  AMTTP-XGB93: F1 from 0.134 to 0.938 (+0.804) — confirms that the BSDT
  component space contains exploitable structure beyond what the
  original linear combination captures.
\item
  \textbf{Distance-based variants (Category A) achieve 100\% recall} in
  10/12 combos, proving that the BSDT feature space inherently separates
  fraud from legitimate transactions. The practical implication: a
  two-stage pipeline using A4\_MahalExp for recall maximisation followed
  by F2\_QuadSurf for precision optimisation would achieve near-perfect
  fraud detection.
\end{enumerate}

\begin{center}\rule{0.5\linewidth}{0.5pt}\end{center}

\subsection{7. Orthogonality and Non-Redundancy
Proof}\label{orthogonality-and-non-redundancy-proof}

\subsubsection{7.1 Pairwise Correlation
Analysis}\label{pairwise-correlation-analysis}

We compute Pearson correlations between all component pairs on both
datasets:

\textbf{Elliptic (active components only — \(G\) is degenerate):}

{\def\LTcaptype{none} % do not increment counter
\begin{longtable}[]{@{}llll@{}}
\toprule\noalign{}
& \(C\) & \(A\) & \(T\) \\
\midrule\noalign{}
\endhead
\bottomrule\noalign{}
\endlastfoot
\(C\) & 1.000 & −0.350 & −0.414 \\
\(A\) & −0.350 & 1.000 & 0.279 \\
\(T\) & −0.414 & 0.279 & 1.000 \\
\end{longtable}
}

\textbf{XBlock (all four active):}

{\def\LTcaptype{none} % do not increment counter
\begin{longtable}[]{@{}lllll@{}}
\toprule\noalign{}
& \(C\) & \(G\) & \(A\) & \(T\) \\
\midrule\noalign{}
\endhead
\bottomrule\noalign{}
\endlastfoot
\(C\) & 1.000 & −0.043 & −0.089 & −0.276 \\
\(G\) & −0.043 & 1.000 & −0.226 & −0.068 \\
\(A\) & −0.089 & −0.226 & 1.000 & 0.504 \\
\(T\) & −0.276 & −0.068 & 0.504 & 1.000 \\
\end{longtable}
}

Mean \(|r|\) (off-diagonal) on XBlock: 0.201. The moderate \(A\)–\(T\)
correlation (0.504) reflects a genuine statistical relationship between
activity anomaly and temporal novelty — high-volume transactions are
more likely to be novel — but the normalised MI analysis (Section 7.3)
confirms they carry largely distinct information.

\textbf{Key observation:} Feature Gap (\(G\)) has zero variance on
Elliptic because Bitcoin transaction features in the Elliptic dataset
are dense (no missing values). This is not a flaw of the theory — it
confirms that the adaptive weighting framework correctly assigns
\(w_G \approx 0\) on Elliptic, as demonstrated in Section 6.5.

\subsubsection{7.2 Principal Component
Analysis}\label{principal-component-analysis}

\textbf{Elliptic (3 active components):} PC1 explains 59.8\%, PC2
explains 30.4\%, PC3 explains 9.8\%. All three active components carry
meaningful variance.

\textbf{XBlock (4 components):} PC1 explains 69.8\%, PC2 explains
21.5\%, PC3 explains 4.7\%, PC4 explains 4.0\%. All four components
contribute, though PC3 and PC4 carry smaller but non-negligible variance
— confirming that four components provide more information than three.

\subsubsection{7.3 Mutual Information
Matrix}\label{mutual-information-matrix}

Normalised pairwise MI (divided by maximum self-MI):

\begin{itemize}
\tightlist
\item
  \textbf{Elliptic:} Mean normalised MI (off-diagonal) = \textbf{0.051}
  ✓
\item
  \textbf{XBlock:} Mean normalised MI (off-diagonal) = \textbf{0.064} ✓
\end{itemize}

Both are well below the 0.30 threshold for information redundancy,
confirming that each component captures genuinely distinct information
about the missed fraud phenomenon.

\subsubsection{7.4 Variance Inflation
Factors}\label{variance-inflation-factors}

VIF measures multicollinearity — values \textless{} 5 indicate
acceptable independence.

\textbf{XBlock:} \(C\) = 1.09, \(G\) = 1.06, \(A\) = 1.41, \(T\) = 1.45
— all excellent. \textbf{Elliptic (active):} \(C\) = 1.30, \(A\) = 1.17,
\(T\) = 1.24 — all excellent.

\subsubsection{7.5 Incremental Predictive
Contribution}\label{incremental-predictive-contribution}

We measure each component’s unique contribution by computing the AUC
drop when removing it from the combined logistic regression model
predicting missed fraud:

\textbf{Elliptic:} \textbar{} Dropped \textbar{} AUC (full) \textbar{}
AUC (reduced) \textbar{} ΔAUC \textbar{}
\textbar—\textbar—\textbar—\textbar—\textbar{} \textbar{} \(C\)
\textbar{} 0.923 \textbar{} 0.933 \textbar{} −0.010 \textbar{}
\textbar{} \(A\) \textbar{} 0.923 \textbar{} 0.924 \textbar{} −0.001
\textbar{} \textbar{} \(T\) \textbar{} 0.923 \textbar{} 0.765 \textbar{}
\textbf{+0.158} \textbar{}

Temporal Novelty (\(T\)) is the dominant predictor of missed fraud on
Elliptic — removing it drops AUC by 15.8 points. This is consistent with
the adaptive weights (Section 6.5) assigning \(w_T > 0.5\).

\textbf{XBlock:} \textbar{} Dropped \textbar{} AUC (full) \textbar{} AUC
(reduced) \textbar{} ΔAUC \textbar{}
\textbar—\textbar—\textbar—\textbar—\textbar{} \textbar{} \(C\)
\textbar{} 0.681 \textbar{} 0.670 \textbar{} +0.011 \textbar{}
\textbar{} \(G\) \textbar{} 0.681 \textbar{} 0.694 \textbar{} −0.013
\textbar{} \textbar{} \(A\) \textbar{} 0.681 \textbar{} 0.754 \textbar{}
\textbf{−0.073} \textbar{} \textbar{} \(T\) \textbar{} 0.681 \textbar{}
0.681 \textbar{} +0.000 \textbar{}

Activity Anomaly (\(A\)) dominates on XBlock — consistent with the MI
weighting (\(w_A = 0.45\)) and the qualitative finding that missed
phishing accounts have abnormal transaction volumes.

\begin{center}\rule{0.5\linewidth}{0.5pt}\end{center}

\subsection{8. Discussion}\label{discussion}

\subsubsection{8.1 Positioning Against OOD Detection, Uncertainty
Estimation, and
Calibration}\label{positioning-against-ood-detection-uncertainty-estimation-and-calibration}

BSDT occupies a distinct niche in the landscape of ML reliability
techniques. We clarify the relationship to three related paradigms:

\textbf{OOD Detection} (Hendrycks \& Gimpel 2017; Liang et al.~2018; Liu
et al.~2020) asks a binary question: “Is this input in-distribution?”
BSDT asks a \emph{diagnostic} question: “If this input is likely
misclassified, \emph{why}?” An OOD detector that flags an input as
anomalous provides no guidance on how to improve the prediction. BSDT’s
four-component decomposition identifies whether the failure is due to
boundary proximity (\(C\)), feature absence (\(G\)), distributional tail
effects (\(A\)), or novelty (\(T\)) — each suggesting a different
corrective action (e.g., feature engineering for \(G\), retraining with
recent data for \(T\)).

\textbf{Uncertainty Estimation} (Gal \& Ghahramani 2016;
Lakshminarayanan et al.~2017) quantifies \emph{how uncertain} a
prediction is, typically through predictive variance or entropy. BSDT
answers the complementary question: \emph{what kind of uncertainty}
affects this prediction? A Monte Carlo dropout ensemble may assign high
uncertainty to a transaction for any of the four BSDT reasons — but the
aggregate uncertainty score alone does not distinguish them. The BSDT
components can be viewed as a \emph{decomposition of aleatoric and
epistemic uncertainty} into actionable failure categories.

\textbf{Calibration} (Guo et al.~2017) adjusts predicted probabilities
to match empirical frequencies. A perfectly calibrated model that
assigns \(\hat{p}(x) = 0.02\) to a camouflaged fraud transaction has
accurate uncertainty — and still misses the fraud. BSDT’s correction
formula (Section 3.4) operates \emph{after} calibration, boosting
predictions for transactions where the decomposition identifies
blind-spot evidence.

\textbf{The unique contribution of BSDT} is not the detection of failure
(which OOD detectors, uncertainty estimators, and calibration
diagnostics all address) but the \emph{structured decomposition} of
failure into four measurable, near-orthogonal axes. This enables
targeted correction — adjusting the prediction based on which failure
mode applies — rather than blanket rejection or uniform uncertainty
inflation.

\subsubsection{8.2 Is the Decomposition
Complete?}\label{is-the-decomposition-complete}

We claim the four components capture the \emph{dominant} failure modes
but acknowledge the residual term \(\varepsilon\). Potential additional
components include:

\begin{itemize}
\tightlist
\item
  \textbf{Adversarial evasion} — intentionally crafted transactions
  designed to fool the model (distinct from natural camouflage)
\item
  \textbf{Label noise} — fraud mislabelled as legitimate in training
  data
\item
  \textbf{Feature interaction blindness} — fraud detectable only through
  higher-order feature interactions the model fails to learn
\end{itemize}

We argue these are either subsumed by the existing components or
represent Bayes error rather than correctable blind spots.

\subsubsection{8.3 Relationship to the Bias-Variance
Decomposition}\label{relationship-to-the-bias-variance-decomposition}

The classical bias-variance decomposition (Geman et al.~1992) partitions
prediction error into:

\[\text{Error} = \text{Bias}^2 + \text{Variance} + \text{Noise}\]

Our decomposition is related but operates in the \emph{feature space}
rather than the \emph{error space}:

\begin{itemize}
\tightlist
\item
  \(C\) (Camouflage) ↔ Bias (the model’s systematic tendency to classify
  certain fraud as legitimate)
\item
  \(T\) (Temporal Novelty) ↔ Variance (sensitivity to the specific
  training sample)
\item
  \(G\) (Feature Gap) ↔ Noise (irreducible due to missing information)
\item
  \(A\) (Activity Anomaly) ↔ a novel component not captured by the
  classical decomposition
\end{itemize}

\subsubsection{8.4 Generalisability Beyond
Blockchain}\label{generalisability-beyond-blockchain}

The four failure modes are defined in terms of general supervised
classification concepts (decision boundaries, feature coverage,
distribution shift, activity profiles). Nothing in the formulation is
blockchain-specific. Our cross-domain validation (Section 6.8)
\textbf{empirically confirms} this prediction:

\begin{itemize}
\tightlist
\item
  \textbf{Credit card fraud detection} — MFLS AUC = 0.885, recall
  improvement +19.0pp
\item
  \textbf{Aerospace telemetry anomaly} — MFLS AUC = 0.982, recall
  improvement +28.9pp
\item
  \textbf{Medical imaging anomaly} — MFLS AUC = 0.916, recall
  improvement +100.0pp
\item
  \textbf{Digit recognition anomaly} — MFLS AUC = 0.972, recall
  improvement +100.0pp
\end{itemize}

Across all four out-of-domain datasets, the BSDT components achieve
strong discriminative power (mean AUC = 0.939) despite the base model
having zero domain-specific knowledge. The adaptive weighting mechanism
automatically discovers which components matter most in each domain:
Activity (\(A\)) dominates in banking fraud, Camouflage (\(C\))
dominates in aerospace telemetry, and Temporal Novelty (\(T\)) dominates
in blockchain. This is a key empirical finding: \textbf{the
decomposition structure is universal, but the dominant component varies
by domain}.

We believe the theory extends further to insurance claims fraud, tax
evasion detection, and cyber intrusion detection, where the same four
failure modes apply.

\subsubsection{8.5 Limitations}\label{limitations}

\begin{enumerate}
\def\labelenumi{\arabic{enumi}.}
\tightlist
\item
  \textbf{\st{Single base model.}} \emph{(Addressed in Section 6.10.)}
  The comprehensive 36-combination evaluation validates BSDT across
  XGBoost, LightGBM, RandomForest, and Logistic Regression, plus three
  pre-trained production pipelines. The four-component decomposition
  improves detection consistently across all architectures tested.
  Remaining future work includes testing with \textbf{(a)} TabNet or
  FT-Transformer (attention-based tabular models), \textbf{(b)}
  GNN-based detectors like EvolveGCN, and \textbf{(c)} deep autoencoders
  (unsupervised).
\item
  \textbf{\st{Component interactions.}} \emph{(Fully addressed in
  Sections 6.11 and 6.12.)} The systematic evaluation of 13 alternative
  combination strategies (Section 6.11) and the exhaustive 66-variant
  exploration (Section 6.12) demonstrate that non-linear component
  interactions are the primary driver of improvement. F2\_QuadSurf — a
  degree-2 polynomial surface fit on \([C, G, A, T]\) — captures all
  pairwise and squared interactions
  (\(C^2, CG, CA, CT, G^2, GA, GT, A^2, AT, T^2\)) and achieves Mean F1
  = 0.787, surpassing both the 5-feature logistic V10 (0.740) and all 66
  alternative strategies. Category F’s dominance of the leaderboard (F2
  at \#1, F4 at \#12, F5 at \#11) definitively demonstrates that
  quadratic interactions among the four BSDT components carry
  substantial discriminative power. Higher-order terms (cubic+) and
  neural network combiners remain untested but are unlikely to yield
  substantial further gains given the diminishing returns from degree-2
  to degree-3 in preliminary experiments.
\item
  \textbf{Causality.} The components correlate with missed fraud but we
  have not established a causal mechanism — the model may miss fraud for
  reasons not captured by any blind-spot component.
\item
  \textbf{Precision trade-off.} On out-of-domain datasets where the base
  model has zero recall, BSDT correction achieves high recall but at
  reduced precision. The correction is most effective when the base
  model has partial (non-zero) domain knowledge. See Section 9.4 for a
  detailed practitioner decision framework on when to use BSDT
  correction versus retraining. As shown in Section 6.9, the signed
  logistic variant substantially mitigates this trade-off, reducing
  average false-alarm rate from 74.2\% to 41.5\% when a small
  calibration set is available. The V2 experiment (Section 6.12) further
  demonstrates that \textbf{false positives are the primary beneficiary}
  of MFLS variant correction: F2\_QuadSurf eliminates 88.1\% of
  aggregate false positives (from 44,740 to 5,345) while lifting
  accuracy from 85.9\% to 96.3\%, and three variants (B2\_KDE\_Sig,
  A9\_kNNCorr, J6\_LOF\_Pow\_Corr) reduce both FP and FN simultaneously.
\end{enumerate}

\begin{center}\rule{0.5\linewidth}{0.5pt}\end{center}

\subsection{9. Practical Implications}\label{practical-implications}

\subsubsection{9.1 For Model Developers}\label{for-model-developers}

Append the seven supplementary features \(\varphi_1\)–\(\varphi_7\) to
training data before retraining. This allows the model to “see its own
blind spots” — learning directly which combinations of camouflage, gap,
activity, and novelty indicate hidden fraud.

\subsubsection{9.2 For Compliance
Officers}\label{for-compliance-officers}

Use MFLS as a \textbf{secondary screening score}. Transactions with
\(\hat{p}(x) < \tau\) but \(\text{MFLS}(x) > 0.5\) should be routed to
manual review. This captures cases the model misses without overwhelming
the review queue.

\subsubsection{9.3 For Regulators}\label{for-regulators}

The BSDT provides a \textbf{standardised vocabulary} for describing ML
model limitations. Rather than opaque accuracy numbers, regulators can
ask: “What is this model’s camouflage vulnerability? What is its
temporal novelty coverage?”

\subsubsection{9.4 When to Use BSDT Correction vs.~When to
Retrain}\label{when-to-use-bsdt-correction-vs.-when-to-retrain}

A critical question for practitioners: \textbf{BSDT correction trades
precision for recall.} The correction deliberately promotes
borderline-legitimate transactions above the threshold, and some of
these will be false positives.

Empirical precision-recall trade-offs across our six datasets:

{\def\LTcaptype{none} % do not increment counter
\begin{longtable}[]{@{}
  >{\raggedright\arraybackslash}p{(\linewidth - 10\tabcolsep) * \real{0.1364}}
  >{\raggedright\arraybackslash}p{(\linewidth - 10\tabcolsep) * \real{0.1818}}
  >{\raggedright\arraybackslash}p{(\linewidth - 10\tabcolsep) * \real{0.2576}}
  >{\raggedright\arraybackslash}p{(\linewidth - 10\tabcolsep) * \real{0.1364}}
  >{\raggedright\arraybackslash}p{(\linewidth - 10\tabcolsep) * \real{0.1970}}
  >{\raggedright\arraybackslash}p{(\linewidth - 10\tabcolsep) * \real{0.0909}}@{}}
\toprule\noalign{}
\begin{minipage}[b]{\linewidth}\raggedright
Dataset
\end{minipage} & \begin{minipage}[b]{\linewidth}\raggedright
Base Recall
\end{minipage} & \begin{minipage}[b]{\linewidth}\raggedright
Corrected Recall
\end{minipage} & \begin{minipage}[b]{\linewidth}\raggedright
Base F1
\end{minipage} & \begin{minipage}[b]{\linewidth}\raggedright
Corrected F1
\end{minipage} & \begin{minipage}[b]{\linewidth}\raggedright
ΔF1
\end{minipage} \\
\midrule\noalign{}
\endhead
\bottomrule\noalign{}
\endlastfoot
Elliptic & 13.5\% & 84.4\% & 0.058 & 0.180 & +0.122 \\
XBlock & 93.5\% & 94.5\% & 0.407 & 0.378 & −0.029 \\
Shuttle & 71.1\% & 100.0\% & 0.316 & 0.133 & −0.183 \\
Credit Card & 0.0\% & 19.0\% & 0.000 & 0.045 & +0.045 \\
\end{longtable}
}

The pattern is clear: \textbf{BSDT correction is most valuable when base
recall is low} (Elliptic, Credit Card) and the F1 improvement justifies
the precision cost. When the base model already has high recall (XBlock,
Shuttle), the correction adds marginal recall at disproportionate
precision cost.

\textbf{Practitioner decision framework:}

\begin{enumerate}
\def\labelenumi{\arabic{enumi}.}
\tightlist
\item
  \textbf{Use BSDT correction (post-hoc, no retraining)} when:

  \begin{itemize}
  \tightlist
  \item
    The deployment setting is \emph{high-recall priority} (AML
    screening, regulatory compliance, medical triage) where missing a
    true positive is far costlier than investigating a false positive
  \item
    The base model has poor cross-domain recall (\(<50\%\)) and
    retraining on in-domain data is infeasible
  \item
    Speed matters — BSDT correction is instantaneous and requires no
    labelled in-domain data
  \end{itemize}
\item
  \textbf{Retrain with MFLS features instead} when:

  \begin{itemize}
  \tightlist
  \item
    The base model already achieves moderate recall (\(>70\%\))
    in-domain
  \item
    Precision matters equally or more than recall (e.g., customer-facing
    fraud alerts)
  \item
    Labelled in-domain data is available — append the seven
    supplementary features \(\varphi_1\)–\(\varphi_7\) (Section 5) and
    retrain for optimal F1
  \end{itemize}
\item
  \textbf{Use both} when:

  \begin{itemize}
  \tightlist
  \item
    Deploying a retrained model but wanting a safety net for concept
    drift — apply BSDT correction as a secondary screening layer with a
    conservative \(\lambda\)
  \end{itemize}
\end{enumerate}

\begin{center}\rule{0.5\linewidth}{0.5pt}\end{center}

\subsection{10. Conclusion}\label{conclusion}

We have proposed the Blind Spot Decomposition Theory, a framework that
characterises machine learning fraud detection failures through four
measurable, near-orthogonal components. The theory is:

\begin{enumerate}
\def\labelenumi{\arabic{enumi}.}
\tightlist
\item
  \textbf{Testable} — any new dataset’s missed fraud should be dominated
  by \(C\), \(G\), \(A\), \(T\) (plus residual \(\varepsilon\)).
\item
  \textbf{Falsifiable} — discovery of a dominant fifth failure mode that
  is orthogonal to \(C\), \(G\), \(A\), \(T\) and improves missed-fraud
  prediction beyond \(\varepsilon\) would expand the basis. The current
  four-mode decomposition should be understood as a \emph{minimal
  sufficient set for the domains tested}, not a closed universal
  partition.
\item
  \textbf{Useful} — the correction formula recovers up to 84.4
  percentage points of recall without retraining.
\item
  \textbf{Adaptive} — self-calibrating weights eliminate the need for
  manual tuning.
\item
  \textbf{Universal} — validated across 6 datasets, 5 domains, 6 model
  architectures, and \textasciitilde398K samples with mean MFLS AUC of
  0.927.
\item
  \textbf{Model-agnostic} — comprehensive evaluation across 36 model ×
  dataset combinations (XGBoost, LightGBM, RandomForest,
  LogisticRegression, plus pre-trained production models) demonstrates
  consistent improvement regardless of base classifier architecture.
\end{enumerate}

The comprehensive 36-combination evaluation (Section 6.10) and
systematic comparison of 13 alternative combination strategies (Section
6.11) provide the strongest evidence for both universality and optimal
integration. The recommended V10 integration — a 5-feature logistic
regression on \([\hat{p}(x), C, G, A, T]\) — outperforms all
alternatives:

{\def\LTcaptype{none} % do not increment counter
\begin{longtable}[]{@{}
  >{\raggedright\arraybackslash}p{(\linewidth - 8\tabcolsep) * \real{0.1194}}
  >{\raggedright\arraybackslash}p{(\linewidth - 8\tabcolsep) * \real{0.1642}}
  >{\raggedright\arraybackslash}p{(\linewidth - 8\tabcolsep) * \real{0.1642}}
  >{\raggedright\arraybackslash}p{(\linewidth - 8\tabcolsep) * \real{0.3134}}
  >{\raggedright\arraybackslash}p{(\linewidth - 8\tabcolsep) * \real{0.2388}}@{}}
\toprule\noalign{}
\begin{minipage}[b]{\linewidth}\raggedright
Metric
\end{minipage} & \begin{minipage}[b]{\linewidth}\raggedright
Base Model
\end{minipage} & \begin{minipage}[b]{\linewidth}\raggedright
+Signed LR
\end{minipage} & \begin{minipage}[b]{\linewidth}\raggedright
\textbf{+V10 (LR + Base)}
\end{minipage} & \begin{minipage}[b]{\linewidth}\raggedright
V10 Improvement
\end{minipage} \\
\midrule\noalign{}
\endhead
\bottomrule\noalign{}
\endlastfoot
Mean \% Correct & 85.9\% & 93.3\% & \textbf{95.7\%} & +9.8pp \\
Mean \% Fraud Missed & 34.8\% & 27.9\% & \textbf{21.8\%} & −13.0pp \\
Mean F1 & 0.551 & 0.634 & \textbf{0.740} & +0.189 \\
Mean AUC & 0.717 & 0.905 & \textbf{0.937} & +0.220 \\
OOD \% Correct & 87.5\% & 98.6\% & \textbf{99.0\%} & +11.5pp \\
OOD F1 & 0.425 & 0.580 & \textbf{0.695} & +0.270 \\
\end{longtable}
}

V10 wins on both F1 and accuracy in all 12 tested combinations — the
only method to achieve a perfect win count. Its superiority validates
the \textbf{complementary-tool thesis}: BSDT is most effective when the
base model’s own prediction is preserved as an input feature, allowing
the learned combiner to trust the base model when confident and override
it using blind-spot signals when uncertain.

The cross-domain validation confirms universality: even when the base
model has zero domain-specific knowledge, the four BSDT components
achieve combined AUC of 0.885–0.994 on completely out-of-domain data.
The adaptive weighting mechanism automatically discovers which component
dominates in each domain — Temporal Novelty (\(T\)) for blockchain,
Activity Anomaly (\(A\)) for banking, Camouflage (\(C\)) for aerospace
telemetry — confirming that \textbf{the decomposition structure is
universal, but the weight configuration is domain-specific}.

The systematic strategy comparison (Section 6.11) also demonstrates that
heuristic combination rules (voting, multiplicative, min-of-4) fail
catastrophically — all producing F1 below 0.30 — confirming that the
BSDT components are informative features that require a \emph{learned,
direction-aware combination rule}, not standalone decision thresholds.

The extended 66-variant exploration (Section 6.12) provides the most
exhaustive validation to date, testing 10 algorithmic families across
792 evaluations with zero errors. The results strengthen the case for
BSDT as a universal feature space:

{\def\LTcaptype{none} % do not increment counter
\begin{longtable}[]{@{}
  >{\raggedright\arraybackslash}p{(\linewidth - 6\tabcolsep) * \real{0.1111}}
  >{\raggedright\arraybackslash}p{(\linewidth - 6\tabcolsep) * \real{0.1528}}
  >{\raggedright\arraybackslash}p{(\linewidth - 6\tabcolsep) * \real{0.2778}}
  >{\raggedright\arraybackslash}p{(\linewidth - 6\tabcolsep) * \real{0.4583}}@{}}
\toprule\noalign{}
\begin{minipage}[b]{\linewidth}\raggedright
Metric
\end{minipage} & \begin{minipage}[b]{\linewidth}\raggedright
Base Model
\end{minipage} & \begin{minipage}[b]{\linewidth}\raggedright
V10 (Section 6.11)
\end{minipage} & \begin{minipage}[b]{\linewidth}\raggedright
\textbf{F2\_QuadSurf (Section 6.12)}
\end{minipage} \\
\midrule\noalign{}
\endhead
\bottomrule\noalign{}
\endlastfoot
Mean F1 & 0.551 & 0.740 & \textbf{0.787} \\
Mean Accuracy & 85.9\% & 95.7\% & \textbf{96.3\%} \\
Aggregate FP & 44,740 & — & \textbf{5,345 (−88.1\%)} \\
OOD F1 & 0.425 & 0.695 & \textbf{0.763} \\
\end{longtable}
}

F2\_QuadSurf’s quadratic surface fit on \([C, G, A, T]\) captures
non-linear component interactions that both linear MFLS and logistic
regression miss, while three variants (B2\_KDE\_Sig, A9\_kNNCorr,
J6\_LOF\_Pow\_Corr) simultaneously reduce both false positives and false
negatives — a rare compliance-optimal outcome. Distance-based variants
achieve 100\% recall in 10/12 combinations, proving that the BSDT
component space inherently separates fraud from legitimate transactions.

The framework transforms missed fraud from an opaque residual into a
structured, measurable, and correctable phenomenon.

\begin{center}\rule{0.5\linewidth}{0.5pt}\end{center}

\subsection{References}\label{references}

Ben-David, S., Blitzer, J., Crammer, K., Kulesza, A., Pereira, F., \&
Vaughan, J. W. (2010). A theory of learning from different domains.
\emph{Machine Learning}, 79(1-2), 151-175.

Chen, T., \& Guestrin, C. (2016). XGBoost: A scalable tree boosting
system. \emph{Proceedings of the 22nd ACM SIGKDD}, 785-794.

Cover, T. M., \& Thomas, J. A. (2006). \emph{Elements of Information
Theory}. Wiley-Interscience.

Dal Pozzolo, A., Caelen, O., Johnson, R. A., \& Bontempi, G. (2015).
Calibrating probability with undersampling for unbalanced
classification. \emph{IEEE Symposium on Computational Intelligence and
Data Mining (CIDM)}, 159-166.

Devroye, L., Györfi, L., \& Lugosi, G. (1996). \emph{A Probabilistic
Theory of Pattern Recognition}. Springer.

El-Yaniv, R., \& Wiener, Y. (2010). On the foundations of noise-free
selective classification. \emph{Journal of Machine Learning Research},
11, 1605-1641.

Gal, Y., \& Ghahramani, Z. (2016). Dropout as a Bayesian approximation:
Representing model uncertainty in deep learning. \emph{ICML 2016},
1050-1059.

Geifman, Y., \& El-Yaniv, R. (2017). Selective classification for deep
neural networks. \emph{NeurIPS 2017}, 4878-4887.

Geman, S., Bienenstock, E., \& Doursat, R. (1992). Neural networks and
the bias/variance dilemma. \emph{Neural Computation}, 4(1), 1-58.

Goodfellow, I. J., Shlens, J., \& Szegedy, C. (2015). Explaining and
harnessing adversarial examples. \emph{ICLR 2015}.

Guo, C., Pleiss, G., Sun, Y., \& Weinberger, K. Q. (2017). On
calibration of modern neural networks. \emph{ICML 2017}, 1321-1330.

Han, S., Hu, X., Huang, H., Jiang, M., \& Zhao, Y. (2022). ADBench:
Anomaly detection benchmark. \emph{NeurIPS 2022 Datasets and Benchmarks
Track}.

Hendrycks, D., \& Gimpel, K. (2017). A baseline for detecting
misclassified and out-of-distribution examples in neural networks.
\emph{ICLR 2017}.

Krogh, A., \& Vedelsby, J. (1995). Neural network ensembles, cross
validation, and active learning. \emph{NeurIPS 1995}, 231-238.

Kuncheva, L. I. (2003). Measures of diversity in classifier ensembles
and their relationship with the ensemble accuracy. \emph{Machine
Learning}, 51(2), 181-207.

Lakshminarayanan, B., Pritzel, A., \& Blundell, C. (2017). Simple and
scalable predictive uncertainty estimation using deep ensembles.
\emph{NeurIPS 2017}, 6402-6413.

Liang, S., Li, Y., \& Srikant, R. (2018). Enhancing the reliability of
out-of-distribution image detection in neural networks. \emph{ICLR
2018}.

Liu, W., Wang, X., Owens, J., \& Li, Y. (2020). Energy-based
out-of-distribution detection. \emph{NeurIPS 2020}, 21464-21475.

Lundberg, S. M., \& Lee, S. I. (2017). A unified approach to
interpreting model predictions. \emph{NeurIPS 2017}, 4765-4774.

Nushi, B., Kamar, E., \& Horvitz, E. (2018). Towards accountable AI:
Hybrid human-machine analyses for characterizing system failure.
\emph{AAAI 2018}.

Quinonero-Candela, J., Sugiyama, M., Schwaighofer, A., \& Lawrence, N.
D. (2009). \emph{Dataset Shift in Machine Learning}. MIT Press.

Rayana, S. (2016). ODDS Library {[}http://odds.cs.stonybrook.edu{]}.
Stony Brook University, Department of Computer Sciences.

Ribeiro, M. T., Singh, S., \& Guestrin, C. (2016). “Why should I trust
you?” Explaining the predictions of any classifier. \emph{Proceedings of
the 22nd ACM SIGKDD}, 1135-1144.

Valiant, L. G. (1984). A theory of the learnable. \emph{Communications
of the ACM}, 27(11), 1134-1142.

Vovk, V., Gammerman, A., \& Shafer, G. (2005). \emph{Algorithmic
Learning in a Random World}. Springer.

Weber, M., Domeniconi, G., Chen, J., Weidele, D. K. I., Bellei, C.,
Robinson, T., \& Leiserson, C. E. (2019). Anti-money laundering in
Bitcoin: Experimenting with graph convolutional networks for financial
forensics. \emph{KDD Workshop on Anomaly Detection in Finance}.

Wu, J., Yuan, Q., Lin, D., You, W., Chen, W., Chen, C., \& Zheng, Z.
(2022). Who are the phishers? Phishing scam detection on Ethereum via
network embedding. \emph{IEEE Transactions on Systems, Man, and
Cybernetics: Systems}, 52(2), 1156-1166.

Xu, C., Tao, D., \& Xu, C. (2013). A survey on multi-view learning.
\emph{arXiv:1304.5634}.

Xu, K., Zhang, M., Li, J., Du, S. S., Kawarabayashi, K. I., \& Jegelka,
S. (2020). How neural networks extrapolate: From feedforward to graph
neural networks. \emph{NeurIPS 2020}.

Zhao, Y., Nasrullah, Z., \& Li, Z. (2019). PyOD: A Python toolbox for
scalable outlier detection. \emph{Journal of Machine Learning Research},
20(96), 1-7.

\begin{center}\rule{0.5\linewidth}{0.5pt}\end{center}

\subsection{Data and Code
Availability}\label{data-and-code-availability}

All public datasets used in this study are freely available: Elliptic
(Weber et al.~2019), XBlock (Wu et al.~2022), Credit Card Fraud (Dal
Pozzolo et al.~2015, via Kaggle), and Shuttle, Mammography, Pendigits
(Rayana 2016, ODDS repository). The proprietary AML training data cannot
be shared due to regulatory restrictions, but the BSDT framework is
fully applicable to any pre-trained model and requires only the model’s
probability outputs and feature values — it does not depend on access to
the original training data. The complete BSDT implementation (component
computation, weight calibration, correction formula, all 66 MFLS
variants, and evaluation scripts) is available at
https://github.com/{[}repository-to-be-published{]}.

\begin{center}\rule{0.5\linewidth}{0.5pt}\end{center}

\subsection{Appendix A: Full Mathematical
Specification}\label{appendix-a-full-mathematical-specification}

\subsubsection{A.1 Notation}\label{a.1-notation}

{\def\LTcaptype{none} % do not increment counter
\begin{longtable}[]{@{}ll@{}}
\toprule\noalign{}
Symbol & Definition \\
\midrule\noalign{}
\endhead
\bottomrule\noalign{}
\endlastfoot
\(x \in \mathbb{R}^d\) & Transaction feature vector (\(d = 93\)) \\
\(y(x) \in \{0, 1\}\) & True label (0 = legitimate, 1 = fraud) \\
\(f(x) = \hat{p}(x)\) & Model’s predicted fraud probability \\
\(\tau\) & Classification threshold \\
\(\mathcal{M}_f\) & Set of missed fraud transactions \\
\(S_i(x)\) & \(i\)-th component score (\(i \in \{C, G, A, T\}\)) \\
\(w_i\) & Weight for component \(i\) \\
\(\lambda\) & Correction strength parameter \\
\(\boldsymbol{\alpha}\) & Dirichlet concentration parameters \\
\end{longtable}
}

\subsubsection{A.2 Algorithm: Adaptive MFLS
Pipeline}\label{a.2-algorithm-adaptive-mfls-pipeline}

\begin{verbatim}
Input: Dataset X, model f, [optional: labels y]
Output: Corrected predictions p*

1. Compute reference statistics from X (or from training data)
2. Compute component scores S = [C(X), G(X), A(X), T(X)]
3. Estimate weights:
   - If labels y available: w = MI_weights(S, y)
   - If no labels: w = VR_weights(S, f(X))
   - If streaming: w = Bayesian_update(S, feedback)
4. Compute MFLS = S @ w
5. Grid search for optimal λ and τ (or use defaults)
6. Correct: p*(x) = p(x) + λ · MFLS(x) · (1-p(x)) · 𝟙[p(x)<τ]
7. Return p*
\end{verbatim}

\subsubsection{A.3 Saved Parameters
(JSON-serialisable)}\label{a.3-saved-parameters-json-serialisable}

\begin{Shaded}
\begin{Highlighting}[]
\FunctionTok{\{}
  \DataTypeTok{"theory"}\FunctionTok{:} \StringTok{"Blind Spot Decomposition Theory"}\FunctionTok{,}
  \DataTypeTok{"version"}\FunctionTok{:} \StringTok{"1.0"}\FunctionTok{,}
  \DataTypeTok{"components"}\FunctionTok{:} \OtherTok{[}\StringTok{"C\_camouflage"}\OtherTok{,} \StringTok{"G\_feature\_gap"}\OtherTok{,} \StringTok{"A\_activity"}\OtherTok{,} \StringTok{"T\_novelty"}\OtherTok{]}\FunctionTok{,}
  \DataTypeTok{"calibration\_methods"}\FunctionTok{:} \OtherTok{[}\StringTok{"mutual\_information"}\OtherTok{,} \StringTok{"variance\_ratio"}\OtherTok{,} \StringTok{"bayesian\_online"}\OtherTok{,} \StringTok{"gradient"}\OtherTok{]}\FunctionTok{,}
  \DataTypeTok{"correction\_formula"}\FunctionTok{:} \StringTok{"p*(x) = p(x) + λ·MFLS(x)·(1{-}p(x))·𝟙[p(x)\textless{}τ]"}
\FunctionTok{\}}
\end{Highlighting}
\end{Shaded}

\begin{center}\rule{0.5\linewidth}{0.5pt}\end{center}

\subsection{Appendix B: Choice of Functional Forms and
Sensitivity}\label{appendix-b-choice-of-functional-forms-and-sensitivity}

Each component’s formula was chosen to satisfy three design constraints:
\textbf{(i)} output in \([0,1]\) for commensurability, \textbf{(ii)}
monotonicity in the underlying failure-mode severity, and \textbf{(iii)}
computability from the feature vector alone (no model internals
required). We justify each choice and discuss sensitivity to alternative
forms.

\subsubsection{\texorpdfstring{B.1 Camouflage —
\(C(x) = 1 - \|x - \mu_{\text{legit}}\|_2 / d_{\max}\)}{B.1 Camouflage — C(x) = 1 - \textbackslash\textbar x - \textbackslash mu\_\{\textbackslash text\{legit\}\}\textbackslash\textbar\_2 / d\_\{\textbackslash max\}}}\label{b.1-camouflage-cx-1---x---mu_textlegit_2-d_max}

\textbf{Why this form?} Camouflage measures geometric proximity to the
legitimate cluster. A normalised Euclidean distance, clipped to
\([0,1]\) and inverted (so higher = more camouflaged), is the simplest
metric satisfying constraint (i). The 99th-percentile cap (\(d_{\max}\))
prevents outlier distortion.

\textbf{Alternatives considered:} Mahalanobis distance (accounts for
covariance but requires invertible covariance matrix, which fails on
high-dimensional sparse data), kernel density estimation (higher
fidelity but \(O(n^2)\) at inference). In sensitivity tests on Elliptic,
replacing Euclidean with cosine distance changed \(C\) scores by
\(<0.04\) mean absolute deviation and left MFLS AUC within \(\pm 0.01\).

\subsubsection{\texorpdfstring{B.2 Feature Gap —
\(G(x) = |\{j : |x_j| < \epsilon\}| / d\)}{B.2 Feature Gap — G(x) = \textbar\textbackslash\{j : \textbar x\_j\textbar{} \textless{} \textbackslash epsilon\textbackslash\}\textbar{} / d}}\label{b.2-feature-gap-gx-j-x_j-epsilon-d}

\textbf{Why this form?} Feature Gap is inherently a counting measure —
what fraction of features carry no information? The formula is arguably
the only natural definition: a proportion of near-zero features. The
threshold \(\epsilon = 10^{-8}\) distinguishes true zeros from
floating-point noise.

\textbf{Alternatives considered:} Entropy-based sparsity (less
interpretable), \(\ell_0\)-norm proxy (equivalent). The choice is robust
because \(G\) is either active (sparse data) or degenerate (dense data),
and the adaptive weighting handles both cases.

\subsubsection{\texorpdfstring{B.3 Activity Anomaly —
\(A(x) = \sigma\!\left((\log(1 + |\text{tx\_count}|) - \mu_{\text{caught}}) / \sigma_{\text{caught}}\right)\)}{B.3 Activity Anomaly — A(x) = \textbackslash sigma\textbackslash!\textbackslash left((\textbackslash log(1 + \textbar\textbackslash text\{tx\textbackslash\_count\}\textbar) - \textbackslash mu\_\{\textbackslash text\{caught\}\}) / \textbackslash sigma\_\{\textbackslash text\{caught\}\}\textbackslash right)}}\label{b.3-activity-anomaly-ax-sigmaleftlog1-texttx_count---mu_textcaught-sigma_textcaughtright}

\textbf{Why this form?} Transaction counts follow a heavy-tailed
distribution; the \(\log(1+|\cdot|)\) transform stabilises variance, and
z-scoring against the \emph{caught-fraud} distribution measures
deviation from what the model has learned to detect. The logistic
sigmoid maps the z-score to \([0,1]\) with smooth gradients, centred at
0.5 for z = 0.

\textbf{Alternatives considered:} Raw z-score (unbounded), CDF of the
empirical distribution (requires storing the full distribution).
Replacing \(\sigma(\cdot)\) with a min-max normalisation changed MFLS
AUC by \(< 0.008\) on both blockchain datasets.

\subsubsection{\texorpdfstring{B.4 Temporal Novelty —
\(T(x) = \sigma\!\left(\frac{1}{2}(\bar{m}(x) - 2)\right)\), where
\(\bar{m}(x)\) is the mean normalised squared
deviation}{B.4 Temporal Novelty — T(x) = \textbackslash sigma\textbackslash!\textbackslash left(\textbackslash frac\{1\}\{2\}(\textbackslash bar\{m\}(x) - 2)\textbackslash right), where \textbackslash bar\{m\}(x) is the mean normalised squared deviation}}\label{b.4-temporal-novelty-tx-sigmaleftfrac12barmx---2right-where-barmx-is-the-mean-normalised-squared-deviation}

\textbf{Why this form?} \(T\) is a soft normalised Mahalanobis-type
distance. Under the null hypothesis that \(x\) is drawn from the
reference fraud distribution, \(\bar{m}(x) \approx 1\). The shift by
\(-2\) and scale by \(1/2\) centres the sigmoid transition around the
point where the feature profile is moderately novel
(\(\bar{m} \approx 2\), i.e., each feature deviates by
\textasciitilde{}\(\sqrt{2}\) standard deviations on average).

\textbf{Alternatives considered:} Full Mahalanobis (requires invertible
covariance), isolation forest score (model-dependent, violating
constraint iii). Varying the centring constant from 1.5 to 3.0 shifted
MFLS AUC by \(< 0.015\) — the sigmoid’s saturation makes \(T\) robust to
this parameter.

\subsubsection{B.5 General Sensitivity
Summary}\label{b.5-general-sensitivity-summary}

The MFLS framework is intentionally robust to functional-form choices
because the adaptive weight calibration absorbs much of the sensitivity.
A component with a sub-optimal formula (e.g., using cosine instead of
Euclidean for \(C\)) will receive a slightly different weight but
produce a similar combined MFLS. The framework’s power arises from the
\emph{decomposition structure} (four orthogonal axes of failure), not
from the precision of any individual component formula.

\end{document}
